\documentclass[11pt,a4paper,openany]{book}

% ── Packages ──────────────────────────────────────────────────
\usepackage[utf8]{inputenc}
\usepackage[T1]{fontenc}
\usepackage{amsmath,amssymb,amsthm}
\usepackage{mathtools}
\usepackage{mathrsfs}
\usepackage{bbm}
\usepackage{bm}
\usepackage[margin=1in]{geometry}
\usepackage[colorlinks=true,linkcolor=blue,citecolor=blue]{hyperref}
\usepackage{graphicx}
\usepackage{booktabs}
\usepackage{array}
\usepackage{enumitem}
\usepackage{fancyhdr}
\usepackage{tocloft}
\usepackage{algorithm}
\usepackage{algpseudocode}

% ── Theorem Environments ─────────────────────────────────────
\theoremstyle{definition}
\newtheorem{definition}{Definition}[chapter]
\newtheorem{example}[definition]{Example}

\theoremstyle{plain}
\newtheorem{theorem}{Theorem}[chapter]
\newtheorem{lemma}[theorem]{Lemma}
\newtheorem{proposition}[theorem]{Proposition}
\newtheorem{corollary}[theorem]{Corollary}

\theoremstyle{remark}
\newtheorem{remark}[theorem]{Remark}

% ── Math Operators ───────────────────────────────────────────
\DeclareMathOperator{\E}{\mathbb{E}}
\DeclareMathOperator{\Var}{Var}
\DeclareMathOperator{\Cov}{Cov}
\DeclareMathOperator{\tr}{tr}
\DeclareMathOperator*{\esssup}{ess\,sup}
\DeclareMathOperator{\id}{id}
\DeclareMathOperator*{\argmax}{arg\,max}
\DeclareMathOperator*{\argmin}{arg\,min}
\newcommand{\1}{\mathbbm{1}}
\newcommand{\R}{\mathbb{R}}
\newcommand{\F}{\mathcal{F}}
\newcommand{\B}{\mathcal{B}}
\newcommand{\N}{\mathcal{N}}
\newcommand{\M}{\mathcal{M}}
\newcommand{\Q}{\mathbb{Q}}
\newcommand{\Pbb}{\mathbb{P}}
\newcommand{\calL}{\mathcal{L}}
\newcommand{\calG}{\mathcal{G}}
\newcommand{\calH}{\mathcal{H}}
\newcommand{\calD}{\mathcal{D}}
\newcommand{\calE}{\mathcal{E}}
\newcommand{\calC}{\mathcal{C}}
\newcommand{\calT}{\mathcal{T}}
\newcommand{\calS}{\mathcal{S}}
\newcommand{\calP}{\mathcal{P}}

% ── Header / Footer ──────────────────────────────────────────
\pagestyle{fancy}
\fancyhf{}
\fancyhead[LE,RO]{\thepage}
\fancyhead[LO]{\nouppercase{\leftmark}}
\fancyhead[RE]{\nouppercase{\rightmark}}
\renewcommand{\headrulewidth}{0.4pt}
\setlength{\headheight}{26pt}

% ── Title ────────────────────────────────────────────────────
\title{\Huge Continuous-Time Markov Chain Approximation\\[4pt] for American Option Pricing\\[12pt]
\large From Generator Convergence to Optimal Stopping}
\author{}
\date{}

\begin{document}

\maketitle

\setcounter{tocdepth}{1}
\tableofcontents

\chapter*{Reader Warning: What This Project Actually Requires}
\addcontentsline{toc}{chapter}{Reader Warning: What This Project Actually Requires}

\begin{quote}
\itshape This document is intentionally written at full mathematical resolution. The problem is not ``use a Markov chain and do backward induction.'' The problem is to replace a continuous-state stochastic process by a finite-state continuous-time Markov chain in a way that preserves the generator, the semigroup, the weak law of the path, and finally the optimal-stopping value. Every one of those arrows hides a theorem.
\end{quote}

\section*{Prerequisite Audit}

A reader who has only seen elementary probability should not expect the project to be conceptually linear. The following list is the minimum background needed to read the argument without reducing it to a collection of formulas:

\begin{center}
\begin{tabular}{>{\raggedright\arraybackslash}p{4cm}>{\raggedright\arraybackslash}p{8.5cm}}
\toprule
\textbf{Layer} & \textbf{Objects that must be controlled} \\
\midrule
Measure-theoretic probability & $\sigma$-algebras, conditional expectation, filtrations, stopping times, martingales. \\
Stochastic processes & Markov property, transition kernels, Feller processes, Skorokhod-space weak convergence. \\
Operator theory & $C_0$-semigroups, infinitesimal generators, domains, cores, Hille--Yosida, Trotter--Kato. \\
CTMC approximation & Conservative $Q$-matrices, local consistency, moment matching, boundary treatment, matrix exponentials. \\
Option pricing & Snell envelopes, Bermudan approximation, variational inequalities, error propagation through dynamic programming. \\
Estimation & Likelihood construction, state-dependent observation density, filtering for latent regimes, model diagnostics. \\
\bottomrule
\end{tabular}
\end{center}

\noindent The purpose of the document is therefore not merely to compute prices. Its purpose is to explain why the computation is legitimate. Without the generator-level argument, a finite-state Markov chain is only a numerical guess; with the generator, semigroup, and stopping-stability arguments in place, it becomes a convergent approximation scheme.

\section*{The Non-Negotiable Chain}

The project lives or dies by the following chain:
\[
\boxed{
\begin{array}{c}
\text{local moment matching} \\
\Downarrow \\
\text{generator consistency on a core} \\
\Downarrow \\
\text{semigroup convergence} \\
\Downarrow \\
\text{weak convergence of Markov processes} \\
\Downarrow \\
\text{stability of the Snell envelope} \\
\Downarrow \\
\text{convergence of American option prices.}
\end{array}}
\]

\noindent Omitting any link changes the nature of the project. It is no longer a rigorous CTMC approximation framework; it becomes an implementation heuristic.

% ═══════════════════════════════════════════════════════════════
% CHAPTER 0 — THEORY ROADMAP
% ═══════════════════════════════════════════════════════════════
\chapter*{The Generator Approximation Paradigm}
\addcontentsline{toc}{chapter}{The Generator Approximation Paradigm}

\begin{quote}
\itshape This chapter outlines the conceptual architecture of the monograph. It answers three questions: What are we trying to prove? How do the pieces fit together? In what order should the chapters be read? The reader may return to this chapter at any point to reorient themselves within the larger argument.
\end{quote}

\section{The Central Idea}

The entire monograph is organized around a single unifying principle---the correspondence between a Markov process and its infinitesimal generator, and the approximation of the latter by a finite-state continuous-time Markov chain (CTMC):

\[
\boxed{\begin{array}{c}
\text{Markov Process } \{X_{t}\} \;\longrightarrow\; \text{Generator } \mathcal{L} \;\longrightarrow\; \text{CTMC Generator } Q^{(N)} \\
\downarrow \\
\text{Semigroup Convergence (Trotter--Kato)} \\
\downarrow \\
\text{Transition Operators } \Pi_N P_{t}^{(N)}\pi_N f \to P_t f \\
\downarrow \\
\text{American Option Price } V_{0}^{(N,L)} \to V(0,X_{0})
\end{array}}
\]

\noindent Each arrow represents a well-defined mathematical operation:
\begin{enumerate}[label=(\arabic*)]
    \item \textbf{Process $\to$ Generator:} Given the SDE or probabilistic description of $\{X_{t}\}$, extract its infinitesimal generator $\mathcal{L}$ (Chapters~5--9).
    \item \textbf{Generator $\to$ CTMC Generator:} Construct a finite matrix $Q^{(N)}$ whose action on test functions approximates that of $\mathcal{L}$ via local moment matching (Chapters~10--14).
    \item \textbf{CTMC Generator $\to$ Semigroup Convergence:} Apply the Trotter--Kato theorem to conclude that $\Pi_N e^{tQ^{(N)}}\pi_N f \to P_{t}f$ as $N \to \infty$ for admissible test functions (Chapter~9).
    \item \textbf{Semigroup Convergence $\to$ Pricing:} Combine the semigroup convergence with Bermudan dynamic programming to obtain convergence of American option prices (Chapters~15--18).
\end{enumerate}

\section{Part Dependency Diagram}

The logical dependence among the seven parts is:

\[
\begin{array}{ccccccc}
\text{I. Probability} & \longrightarrow & \text{II. Markov Processes} & \longrightarrow & \text{III. Generator Theory} & \longrightarrow & \text{IV. CTMC Approximation} \\
&&&&&& \downarrow \\
&&&&&& \text{V. Optimal Stopping} \\
&&&&&& \downarrow \\
&&&&&& \text{VI. Unified Pricing} \\
&&&&&& \downarrow \\
&&&&&& \text{VII. Parameter Estimation}
\end{array}
\]

\noindent Part~I (Probability) is the measure-theoretic foundation. Part~II (Markov Processes) introduces the Markov property and transition semigroups. Part~III (Generator Theory) defines the central object $\mathcal{L}$ and the convergence theorems (Hille--Yosida, Trotter--Kato). Part~IV (CTMC Approximation) contains the constructive core---how to build $Q^{(N)}$ from $\mathcal{L}$. Part~V (Optimal Stopping) reformulates the American option as a free-boundary problem. Part~VI (Unified Pricing) assembles the pipeline. Part~VII (Parameter Estimation) applies the framework to data.

\section{How to Read This Monograph}

\begin{itemize}
    \item \textbf{First reading (practitioner):} Chapters~1--4 (skim), 5--6, 8, 10--13, 15--16, 18.
    \item \textbf{Second reading (theorist):} Add Chapters~7, 9, 14, 17, plus selected advanced topics.
    \item \textbf{Reference use:} Each chapter is self-contained; consult the ``Purpose'' block at the start and ``Used in'' block at the end.
\end{itemize}

\section{Main Results (Three Theorems)}

The monograph is organized around three hierarchically nested theorems, each depending on the previous:

\[
\boxed{\begin{array}{c}
\textbf{Main Theorem A} \quad \text{(Generator Approximation)} \\
\|\Pi_N Q^{(N)} \pi_N f - \mathcal{L} f\|_{\infty} \leq C_{f} h^{2} \quad \text{on a core} \\
\downarrow \\
\textbf{Main Theorem B} \quad \text{(Weak Convergence)} \\
\Pi_N e^{t Q^{(N)}} \pi_N f \to P_{t}f \;\longrightarrow\; \widetilde{X}^{(N)} \Rightarrow X \text{ in } D([0,T],E) \\
\downarrow \\
\textbf{Main Theorem C} \quad \text{(Pricing Convergence)} \\
\displaystyle \lim_{L \to \infty} \lim_{N \to \infty} V_{0}^{(N,L)} = V(0, X_{0})
\end{array}}
\]

\noindent\textbf{Theorem A} (Chapter~13) establishes that the moment-matching construction produces a CTMC generator $Q^{(N)}$ whose action on smooth test functions approximates the continuous generator $\mathcal{L}$ pointwise, with a local error of order $h^{2}$.

\noindent Here $\pi_N$ denotes restriction/interpolation of a continuous test function onto the finite grid and $\Pi_N$ denotes the corresponding reconstruction back to the ambient function space. This notation is deliberately explicit: a finite matrix and an unbounded differential operator do not live on the same space until the comparison maps are specified.

\noindent\textbf{Theorem B} (Chapter~9, applied in Chapter~12) invokes the Trotter--Kato theorem to lift local generator convergence to global semigroup convergence, and from there to weak convergence of the CTMC process $\widetilde{X}^{(N)}$ to the target process $X$.

\noindent\textbf{Theorem C} (Chapter~18) combines semigroup convergence with the Bermudan dynamic programming recursion to prove convergence of American option prices, with an explicit error decomposition into generator, Bermudan, and boundary components.


\noindent The reader may view the monograph as providing the proof of Theorems~A, B, and C in sequence. Part~I supplies the probabilistic vocabulary. Part~II defines the generator. Part~III proves Theorem~A. Part~IV extends it to all model classes. Part~V builds the optimal stopping framework. Part~VI assembles Theorem~C.

\medskip
\noindent\textbf{Theorem dependency map.} The formal dependency chain is:
\[
\begin{array}{c}
\text{Operator Framework and Hille--Yosida assumptions} \\
\Downarrow \\
\text{Core consistency and conservativity of } Q^{(n)} \\
\Downarrow \\
\text{Theorem A: generator consistency on a core of } C_0(E) \\
\Downarrow \\
\text{Tightness + martingale-problem uniqueness} \\
\Downarrow \\
\text{Theorem B: weak convergence in } D([0,T],E) \\
\Downarrow \\
\text{Snell-envelope stability and uniform integrability} \\
\Downarrow \\
\text{Theorem C: American pricing convergence.}
\end{array}
\]

\section{Conventions on Theorem Hierarchy}

The following hierarchy is observed throughout the monograph:

\begin{center}
\begin{tabular}{>{\raggedright\arraybackslash}p{2.5cm}>{\raggedright\arraybackslash}p{9cm}}
\toprule
\textbf{Environment} & \textbf{Usage} \\
\midrule
\textbf{Definition} & Introduces a new mathematical object or concept. \\
\textbf{Example} & Illustrates a definition or theorem with a concrete instance. \\
\textbf{Lemma} & A technical intermediate result used to prove a theorem or proposition. \\
\textbf{Proposition} & A self-contained result of moderate importance; often a direct consequence of definitions or straightforward computation. \\
\textbf{Theorem} & A major result requiring a non-trivial proof; the statement has structural significance within the theory. \\
\textbf{Corollary} & An immediate consequence of a theorem or proposition. \\
\textbf{Remark} & A comment on mathematical subtleties, limitations, or interpretative guidance. Not a formal result. \\
\textbf{Proof} & A complete verification. Inline derivations (e.g., solving a $2 \times 2$ linear system) are presented as proofs; shorter calculations appear within the text. \\
\bottomrule
\end{tabular}
\end{center}

\noindent In particular: algebraic derivations of transition rates from moment-matching conditions are designated as \textbf{Propositions} (not Theorems), as they follow from elementary linear algebra. Convergence guarantees (Trotter--Kato, consistency order, main pricing theorem) are designated as \textbf{Theorems}. The reader may note that in earlier drafts some Propositions were labeled as Theorems; the current version adopts the strict hierarchy above.

\newpage

% ═══════════════════════════════════════════════════════════════
% PART I — PROBABILITY FOUNDATIONS
% ═══════════════════════════════════════════════════════════════
\part{Foundations}

\begin{quote}
\itshape The material in this part constitutes the measure-theoretic substrate upon which the entire theory of Markov processes, generators, and CTMC approximation is erected. No concession is made to readability at the expense of precision. The reader is assumed to possess, or to be prepared to acquire, facility with measure-theoretic probability at the level of Billingsley~(1995) or Durrett~(2019).
\end{quote}

% ──────────────────────────────────────────────────────────────
\newpage

% ═══════════════════════════════════════════════════════════════
% NOTATION INDEX
% ═══════════════════════════════════════════════════════════════
\chapter*{Notation Index}
\addcontentsline{toc}{chapter}{Notation Index}

\noindent The following notation is used throughout the monograph. Entries are grouped by category: probability spaces, processes, operators, and model parameters. For asymptotic CTMC sequences, the canonical notation is $Q^{(n)}$ for generators, $P_t^{(n)}$ for semigroups, and $E_n$ for grids; superscripts such as $Q^{(N)}$ are retained only in model-specific numerical chapters where $N$ denotes the number of states.

\subsection*{Probability and Measure Theory}

\begin{center}
\begin{tabular}{>{$}l<{$} >{\raggedright\arraybackslash}p{10cm}}
\toprule
\multicolumn{1}{l}{\bfseries Symbol} & \bfseries Meaning \\
\midrule
(\Omega, \mathcal{F}, \mathbb{P}) & Probability space (Chapter~1) \\
(\Omega, \mathcal{F}, \{\mathcal{F}_{t}\}, \mathbb{P}) & Filtered probability space (stochastic basis) \\
\mathbb{E}[\,\cdot\,] & Expectation under $\mathbb{P}$ \\
\mathbb{E}[\,\cdot\mid\mathcal{G}] & Conditional expectation given $\sigma$-algebra $\mathcal{G}$ \\
\mathbb{Q} & Risk-neutral (martingale) measure \\
\mathcal{N} & Collection of $\mathbb{P}$-null sets \\
\mathcal{B}(\mathbb{R}) & Borel $\sigma$-algebra on $\mathbb{R}$ \\
C_{0}(E) & Continuous functions vanishing at infinity on $E$ \\
C_{c}^{\infty}(\mathbb{R}) & Smooth, compactly supported functions (core for generators) \\
L^{p} & Space of $p$-integrable random variables \\
\bottomrule
\end{tabular}
\end{center}

\subsection*{Stochastic Processes}

\begin{center}
\begin{tabular}{>{$}l<{$} >{\raggedright\arraybackslash}p{10cm}}
\toprule
\multicolumn{1}{l}{\bfseries Symbol} & \bfseries Meaning \\
\midrule
\{X_{t}\}_{t \geq 0} & Target continuous-time Markov process \\
\{\widetilde{X}_{t}^{(N)}\}_{t \geq 0} & Approximating CTMC with $N$ states \\
W_{t} & Standard Brownian motion \\
N_{t} & Poisson process \\
N(dt, dy) & Poisson random measure \\
\widetilde{N}(dt, dy) & Compensated Poisson random measure \\
\{R_{t}\} & Regime process (finite-state CTMC) \\
\tau & Stopping time \\
\mathcal{T}_{[t,T]} & Set of stopping times taking values in $[t,T]$ \\
V(t,x) & American option value function \\
\phi(x) & Option payoff function \\
\bottomrule
\end{tabular}
\end{center}

\subsection*{Operators and Semigroups}

\begin{center}
\begin{tabular}{>{$}l<{$} >{\raggedright\arraybackslash}p{10cm}}
\toprule
\multicolumn{1}{l}{\bfseries Symbol} & \bfseries Meaning \\
\midrule
\mathcal{L} & Infinitesimal generator of $\{X_{t}\}$ \\
\mathcal{D}(\mathcal{L}) & Domain of the generator \\
P_{t} & Transition semigroup of $\{X_{t}\}$: $P_{t}f(x) = \mathbb{E}[f(X_{t})\mid X_{0}=x]$ \\
T_{t} & Generic $C_{0}$-semigroup \\
Q^{(N)} & $N \times N$ CTMC generator matrix \\
P_{t}^{(N)} = e^{tQ^{(N)}} & CTMC transition matrix over time $t$ \\
Q_{D}^{(N)} & Diffusion (tridiagonal) part of $Q^{(N)}$ \\
Q_{J}^{(N)} & Jump part of $Q^{(N)}$ \\
Q^{R} & Regime-switching generator matrix \\
\mathcal{L}_{D}, \mathcal{L}_{J} & Diffusion and jump parts of $\mathcal{L}$ \\
\bottomrule
\end{tabular}
\end{center}

\subsection*{CTMC Construction}

\begin{center}
\begin{tabular}{>{$}l<{$} >{\raggedright\arraybackslash}p{10cm}}
\toprule
\multicolumn{1}{l}{\bfseries Symbol} & \bfseries Meaning \\
\midrule
\mathcal{S}^{(N)} = \{x_{1},\ldots,x_{N}\} & Discrete state space (grid) \\
h_{i} = x_{i+1} - x_{i} & Local grid spacing \\
h = \max_{i} h_{i} & Grid fineness parameter \\
q_{ij} & Transition rate from state $i$ to $j$ \\
q_{i}^{+} = q_{i,i+1} & Upward transition rate \\
q_{i}^{-} = q_{i,i-1} & Downward transition rate \\
\mu(x) & Drift coefficient \\
\sigma^{2}(x) & Diffusion (squared volatility) coefficient \\
D_{i} = \sigma^{2}(x_{i}) & Local diffusion variance \\
M_{1}^{J}, M_{2}^{J} & First and second jump moment contributions \\
\nu(dy) & L\'{e}vy measure \\
\bottomrule
\end{tabular}
\end{center}

\subsection*{Dynamic Programming and Pricing}

\begin{center}
\begin{tabular}{>{$}l<{$} >{\raggedright\arraybackslash}p{10cm}}
\toprule
\multicolumn{1}{l}{\bfseries Symbol} & \bfseries Meaning \\
\midrule
L & Number of Bermudan exercise dates \\
\Delta t = T / L & Time step for DP recursion \\
\mathbf{V}^{(\ell)} & Value vector at time step $\ell$ \\
\mathbf{P}_{\Delta t}^{(N)} & CTMC transition matrix over $\Delta t$ \\
\boldsymbol{\phi} & Payoff vector on the grid \\
r & Risk-free discount rate \\
\varepsilon_{\text{gen}}, \varepsilon_{\text{semi}}, \ldots & Error decomposition terms (Appendix~F) \\
\bottomrule
\end{tabular}
\end{center}

\subsection*{Model Parameters}

\begin{center}
\begin{tabular}{>{$}l<{$} >{\raggedright\arraybackslash}p{10cm}}
\toprule
\multicolumn{1}{l}{\bfseries Symbol} & \bfseries Meaning \\
\midrule
\sigma & Volatility parameter (CEV/GBM) \\
\beta & CEV elasticity parameter \\
\lambda & Jump intensity \\
p & Probability of upward jump (DEJD) \\
\eta_{1}, \eta_{2} & Exponential decay rates (positive and negative jumps) \\
\lambda_{mn} & Regime transition rate from $m$ to $n$ \\
M & Number of regimes \\
K = N \cdot M & Total number of CTMC states (regime-switching) \\
\bottomrule
\end{tabular}

\end{center}

\newpage

\chapter*{Operator Framework}
\addcontentsline{toc}{chapter}{Operator Framework}

\section{Purpose and Scope}

This chapter fixes the functional-analytic framework used throughout the monograph. All convergence statements are interpreted on a common state space and Banach space; this prevents the CTMC construction from relying on informal pointwise consistency alone.

\section{State Space and Banach Space}

Let $E$ be a locally compact separable metric space. In the one-dimensional diffusion chapters $E \subset \R$; in the regime-switching chapters $E = \R \times \{1,\ldots,M\}$; after truncation, $E_n \subset E$ denotes a finite grid.

The ambient Banach space is
\[
C_0(E) := \{f:E\to\R \;:\; f \text{ is continuous and vanishes at infinity}\},
\]
with norm
\[
\|f\|_\infty := \sup_{x\in E} |f(x)|.
\]
For bounded payoff functions we also use $C_b(E)$, the space of bounded continuous functions on $E$ with the same supremum norm. For finite grids $E_n$, $C_0(E_n)=\R^{|E_n|}$.

\section{Contraction Semigroups and Generators}

\begin{definition}[Strongly Continuous Contraction Semigroup]
A family of linear operators $(P_t)_{t\ge 0}$ on $C_0(E)$ is a strongly continuous contraction semigroup if
\begin{enumerate}[label=(\roman*)]
    \item $P_0=I$ and $P_{t+s}=P_tP_s$ for all $s,t\ge0$;
    \item $\|P_t f\|_\infty \le \|f\|_\infty$ for all $f\in C_0(E)$;
    \item $\lim_{t\downarrow0}\|P_tf-f\|_\infty=0$ for all $f\in C_0(E)$.
\end{enumerate}
\end{definition}

\begin{definition}[Infinitesimal Generator]
The infinitesimal generator $(\mathcal{L},D(\mathcal{L}))$ of $(P_t)$ is defined by
\[
D(\mathcal{L}) := \left\{ f\in C_0(E): \lim_{t\downarrow0}\frac{P_tf-f}{t}\text{ exists in }\|\cdot\|_\infty \right\},
\]
and
\[
\mathcal{L}f := \lim_{t\downarrow0}\frac{P_tf-f}{t}, \qquad f\in D(\mathcal{L}).
\]
The domain $D(\mathcal{L})$ is part of the operator, not an optional technicality.
\end{definition}

\begin{definition}[Core]
A subspace $\mathcal{D}\subset D(\mathcal{L})$ is a core for $\mathcal{L}$ if the closure of $\mathcal{L}|_{\mathcal{D}}$ in the graph norm
\[
\|f\|_{\mathcal{L}} := \|f\|_\infty + \|\mathcal{L}f\|_\infty
\]
is $(\mathcal{L},D(\mathcal{L}))$.
\end{definition}

\section{Hille--Yosida Applicability}

Whenever this monograph writes that a process has semigroup $(P_t)$ generated by $\mathcal{L}$, the following assumptions are in force:
\begin{enumerate}[label=(HY\arabic*)]
    \item $D(\mathcal{L})$ is dense in $C_0(E)$;
    \item $\mathcal{L}$ is closed, or is the closure of its restriction to an explicitly stated core;
    \item $\mathcal{L}$ is dissipative and satisfies the range/resolvent condition of the Hille--Yosida theorem;
    \item the resulting semigroup is a strongly continuous contraction semigroup on $C_0(E)$.
\end{enumerate}
Thus $P_t$ is not defined as a literal matrix exponential of an unbounded operator. The notation $e^{t\mathcal{L}}$ is used only as formal shorthand for ``the semigroup generated by $\mathcal{L}$.'' In contrast, for a finite CTMC generator matrix $Q^{(n)}$, the expression $e^{tQ^{(n)}}$ is the ordinary matrix exponential.

\section{Generator Convergence on Varying Grids}

For a grid $E_n\subset E$, functions on $E$ are restricted to $E_n$ before applying $Q^{(n)}$. Operator convergence means convergence on a core:
\[
\sup_{x\in E_n}|Q^{(n)}f(x)-\mathcal{L}f(x)|\to0, \qquad f\in\mathcal{D}.
\]
The core property then extends the convergence from $\mathcal{D}$ to the closed generator via the graph-norm closure. This is the mathematical step that separates true operator convergence from mere pointwise moment matching.

\section{Role of Moment Matching}

Moment matching is used only to construct candidate matrices $Q^{(n)}$ and verify local consistency. It is not itself a convergence theorem. The logical chain in this monograph is
\[
\text{moment matching} \Rightarrow \text{core consistency} \Rightarrow \text{semigroup convergence} \Rightarrow \text{weak convergence} \Rightarrow \text{pricing convergence}.
\]

\newpage

\chapter*{Standing Assumptions}
\addcontentsline{toc}{chapter}{Standing Assumptions}

\section{Purpose of This Chapter}
This chapter collects all structural assumptions used throughout the monograph. Every theorem stated after this chapter operates under one of the assumption families defined below. When a result requires a specific assumption, it is referenced explicitly (e.g., ``under Assumption~D''). The purpose is to eliminate ambiguity: without a stated assumption framework, theorems in the sequel would rest on implicit hypotheses scattered across individual statements.

\section{Assumption D (Diffusion Process)}

\begin{description}
    \item[D1 -- Regularity.] $\mu \in C(\R)$ and $\sigma \in C(\R)$.
    \item[D2 -- Ellipticity.] $\sigma(x) > 0$ for all $x \in \R$.
    \item[D3 -- Local Lipschitz.] For every compact $K \subset \R$, there exists $L_{K} < \infty$ such that
    \[
    |\mu(x) - \mu(y)| + |\sigma(x) - \sigma(y)| \leq L_{K} |x-y|, \qquad \forall x,y \in K.
    \]
    \item[D4 -- Linear Growth.] There exists $C < \infty$ such that
    \[
    |\mu(x)| + |\sigma(x)| \leq C(1 + |x|), \qquad \forall x \in \R.
    \]
\end{description}

\noindent Under Assumption~D, the SDE
\[
dX_{t} = \mu(X_{t}) dt + \sigma(X_{t}) dW_{t}
\]
admits a unique strong solution, which is a Feller process. Its infinitesimal generator on $C_{c}^{\infty}(\R)$ is
\[
\mathcal{L} f(x) = \mu(x) f'(x) + \frac{1}{2} \sigma^{2}(x) f''(x).
\]

\section{Assumption J (Jump-Diffusion Process)}

\begin{description}
    \item[J1 -- Diffusion part.] $\mu, \sigma$ satisfy Assumption~D. (Thus D1--D4 hold.)
    \item[J2 -- Jump intensity.] $\lambda \in C_{b}(\R)$ (bounded and continuous) with $\lambda(x) \geq 0$.
    \item[J3 -- L\'{e}vy measure.] $\nu$ is a $\sigma$-finite measure on $\R \setminus \{0\}$ satisfying $\int_{\R} (1 \wedge y^{2}) \, \nu(dy) < \infty$.
    \item[J4 -- Finite second moment of jumps.] $\int_{\R} y^{2} \, \nu(dy) < \infty$.
\end{description}

\noindent Under Assumption~J, the process
\[
dX_{t} = \mu(X_{t-}) dt + \sigma(X_{t-}) dW_{t} + \int_{\R} h(X_{t-}, y) \widetilde{N}(dt, dy)
\]
is well-defined and is a Feller process. Its generator on $C_{c}^{\infty}(\R)$ is
\[
\mathcal{L} f(x) = \mu(x) f'(x) + \frac{1}{2} \sigma^{2}(x) f''(x) + \int_{\R} \bigl[f(x + h(x,y)) - f(x)\bigr] \, \nu(dy).
\]
The decomposition $\mathcal{L}=\mathcal{L}_{D}+\mathcal{L}_{J}$ is understood on the common core $C_c^\infty(\R)$, where both terms are well-defined. Extensions beyond this core are always taken through graph-norm closure; truncation of the L\'{e}vy measure is treated as an approximation of $\mathcal{L}_{J}$ on that core.

\section{Assumption R (Regime-Switching Process)}

\begin{description}
    \item[R1 -- Per-regime diffusion.] For each $m \in \{1,\ldots,M\}$, the coefficients $\mu_{m}, \sigma_{m}$ satisfy Assumption~D.
    \item[R2 -- Regime process.] $\{R_{t}\}$ is an irreducible finite-state CTMC on $\{1,\ldots,M\}$ with generator matrix $Q^{R}$.
    \item[R3 -- Independence.] The Brownian motion $W_{t}$ is independent of $\{R_{t}\}$.
\end{description}

\noindent Under Assumption~R, the joint process $(X_{t}, R_{t})$ is a Feller process on $\R \times \{1,\ldots,M\}$. Its generator on $C_{c}^{\infty}(\R \times \{1,\ldots,M\})$ is
\[
\mathcal{L} f(x, m) = \mu_{m}(x) \frac{\partial f}{\partial x}(x,m) + \frac{1}{2} \sigma_{m}^{2}(x) \frac{\partial^{2} f}{\partial x^{2}}(x,m) + \sum_{n \neq m} \lambda_{mn} \bigl[f(x,n) - f(x,m)\bigr].
\]

\section{Remarks on Generality}

\begin{remark}
Assumption~D (specifically D3--D4) is standard in the SDE literature (Karatzas \& Shreve, 1991; \O ksendal, 2003). It guarantees existence, uniqueness, and the Feller property. For the CEV model with $\beta < 0$, the local Lipschitz condition fails at $x=0$; the process is treated on $(0,\infty)$ with an absorbing or reflecting boundary at zero---see Chapter~\ref{ch:boundary} for the boundary theory.
\end{remark}

\begin{remark}
Assumption~J4 (finite second moment of jumps) is necessary for the CTMC moment-matching construction. If $\int y^{2} \nu(dy) = \infty$, the jump contribution to the second moment diverges, and the moment-matching system has no finite solution. This is why the double-exponential distribution requires $\eta_{1} > 2$ (see Appendix~\ref{app:dejd-mgf}).
\end{remark}

\begin{remark}
Assumption~R is stated for finitely many regimes. The extension to countably many regimes is possible but requires additional tightness conditions; it is not pursued in this monograph.
\end{remark}

\newpage

% ──────────────────────────────────────────────────────────────
\chapter{Probability Space and Random Variables}

\noindent\textbf{Purpose.} Establishes the measure-theoretic foundation: probability spaces, random variables, distributions, expectation, $L^{p}$ spaces, stochastic processes, and modes of convergence. Every subsequent chapter operates on the objects defined here.
\textbf{Prerequisites.} None (starting point of the monograph).
\textbf{Used in.} All chapters. Chapters~2--4 extend these concepts to conditional expectation, martingales, and stopping times.

\section{The Probability Triple}

\begin{definition}[Measurable Space]
A measurable space is a pair $(\Omega, \F)$ where $\Omega$ is a nonempty set and $\F \subseteq 2^{\Omega}$ is a $\sigma$-algebra, i.e., a collection of subsets satisfying
\begin{enumerate}[label=(\arabic*)]
    \item $\Omega \in \F$,
    \item $A \in \F \implies A^{c} \in \F$,
    \item $\{A_{n}\}_{n=1}^{\infty} \subseteq \F \implies \bigcup_{n=1}^{\infty} A_{n} \in \F$.
\end{enumerate}
\end{definition}

\medskip
\noindent\textbf{Why this matters.} A measurable space is the canvas of a random experiment before any color touches it: $\Omega$ is the fabric containing every possible outcome, and $\mathcal{F}$ is the zoning rule that declares which regions of the canvas are legitimate to paint on and talk about. The three axioms of a $\sigma$-algebra guarantee that negating an event or taking a countable union of events never produces a region outside the measurable territory. Without this scaffolding, probability has no footing---any more than a painting can exist without a canvas and a partition of the pigment area.

\begin{definition}[Probability Measure]
Let $(\Omega, \F)$ be a measurable space. A function $\Pbb: \F \to [0,1]$ is a probability measure if
\begin{enumerate}[label=(\arabic*)]
    \item $\Pbb(\Omega) = 1$,
    \item For every countable collection $\{A_{n}\}_{n=1}^{\infty} \subseteq \F$ of pairwise disjoint sets,
    \[
    \Pbb\left(\bigcup_{n=1}^{\infty} A_{n}\right) = \sum_{n=1}^{\infty} \Pbb(A_{n}).
    \]
\end{enumerate}
\end{definition}

\medskip
\noindent\textbf{Why this matters.} The heart of a probability measure is additivity: the probability of two non-overlapping events equals the sum of their individual probabilities. The intuition is the same as splitting a pizza---if the guests at two separate tables have nothing to do with each other, the total head count is just the sum of the two table counts. Countable additivity extends this intuition to infinitely many possibilities, providing the backbone of rigorous probabilistic reasoning.

\begin{definition}[Probability Space]
The triple $(\Omega, \F, \Pbb)$ is called a probability space.
\end{definition}

\medskip
\noindent\textbf{Why this matters.} One can think of a probability space as the complete blueprint of a random experiment: $\Omega$ is the catalog of all possible outcomes, $\mathcal{F}$ specifies which events are observable and measurable, and $\mathbb{P}$ assigns a weight between 0 and 1 to each observable event. The triple is indispensable---remove any one component, and the mathematical description of the random experiment collapses.

\begin{definition}[Complete Probability Space]
A probability space $(\Omega, \F, \Pbb)$ is \textbf{complete} if $\F$ contains all subsets of $\Pbb$-null sets: for every $A \in \F$ with $\Pbb(A) = 0$ and every $B \subseteq A$, one has $B \in \F$. Every probability space admits a completion; in what follows, we always work with a complete probability space, augmenting filtrations accordingly.
\end{definition}

\medskip
\noindent\textbf{Why this matters.} Completeness is a technical safety net that prevents pathological edge cases from undermining the theory. If an event has probability zero, every sub-event of it ought to be negligible as well. Without completeness, zero-measure sets can harbor bizarre counterexamples---much like an idealized physics experiment that assumes zero air resistance: once a real-world draft slips in, the whole derivation crumbles. Completing the probability space is therefore like installing a surge protector on the mathematical wiring.

\section{Random Variables and Their Distributions}

\begin{definition}[Random Variable]
Let $(\Omega, \F)$ and $(E, \calE)$ be measurable spaces. A function $X: \Omega \to E$ is $\F/\calE$-measurable (or simply a random variable when $E = \R$ and $\calE = \B(\R)$) if $X^{-1}(B) \in \F$ for every $B \in \calE$.
\end{definition}

\medskip
\noindent\textbf{Why this matters.} A random variable is not ``a variable that is random'' but a measuring instrument---it maps the abstract outcome space into a numerical space that we can read. Consider a thermometer that maps the state of the air onto a degree reading: the air may be in infinitely many configurations, but the thermometer reports a single number. Measurability guarantees that every meaningful numerical interval corresponds to a measurable event, so we can discuss ``the probability that the temperature falls between 20 and 25 degrees.''

\begin{definition}[Distribution of a Random Variable]
The distribution, or law, of a random variable $X: \Omega \to E$ is the pushforward measure $\mu_{X} := \Pbb \circ X^{-1}$ defined on $(E, \calE)$ by
\[
\mu_{X}(B) = \Pbb(X \in B), \qquad B \in \calE.
\]
\end{definition}

\medskip
\noindent\textbf{Why this matters.} The distribution is the fingerprint of a random variable---it describes the probabilities of falling into various numerical intervals while remaining completely indifferent to what $\Omega$ looks like internally. Two people flipping coins with different materials and tossing styles share the same distribution as long as both coins land heads with probability 50\%. This abstraction frees us from the underlying mechanism: the distribution alone encodes everything needed about the statistical behavior of a random variable.

\begin{definition}[Expectation]
Let $X: \Omega \to \R$ be a random variable. If $X \geq 0$, its expectation is defined via the Lebesgue integral:
\[
\E[X] := \int_{\Omega} X \, d\Pbb.
\]
For a general $X$, decompose $X = X^{+} - X^{-}$ and set $\E[X] := \E[X^{+}] - \E[X^{-}]$ whenever at least one of the terms is finite. If $\E[|X|] < \infty$, we write $X \in L^{1}(\Omega, \F, \Pbb)$.
\end{definition}

\medskip
\noindent\textbf{Why this matters.} The expectation is the mathematically precise version of a long-run average. Picture a roulette wheel at a casino: each spin is unpredictable, but after thousands of spins the frequency of each number settles near a fixed value---the expectation is that stable center of mass. The Lebesgue integral is more powerful than the Riemann integral because it partitions the range of function values rather than the domain, which lets it handle a much wider and stranger class of random variables.

\begin{definition}[$L^{p}$ Spaces]
For $p \in [1, \infty)$, define
\[
L^{p}(\Omega, \F, \Pbb) := \left\{X : \Omega \to \R \text{ measurable} : \E[|X|^{p}] < \infty \right\},
\]
equipped with the norm $\|X\|_{L^{p}} := (\E[|X|^{p}])^{1/p}$. For $p = \infty$, $L^{\infty}(\Omega, \F, \Pbb)$ is the space of essentially bounded random variables with norm $\|X\|_{L^{\infty}} := \inf\{M \geq 0 : \Pbb(|X| > M) = 0\}$. Each $(L^{p}, \|\cdot\|_{L^{p}})$ is a Banach space; $L^{2}$ is a Hilbert space under the inner product $\langle X, Y \rangle = \E[XY]$.
\end{definition}

\medskip
\noindent\textbf{Why this matters.} $L^{p}$ spaces classify random variables by their degree of fluctuation. $L^{1}$ is like a steady neighborhood grocery---it has a predictable average revenue but no guarantee that the variance stays finite. $L^{2}$ is like a firm with a well-documented budget volatility---not only can you compute the mean, but you can also calculate variance and correlations. The fact that $L^{2}$ is a Hilbert space is especially powerful because it lets us understand relationships among random variables geometrically: covariance is essentially an inner product, and orthogonality corresponds to uncorrelatedness.

\begin{proposition}[Basic Inequalities]
Let $X, Y \in L^{1}(\Omega, \F, \Pbb)$. Then:
\begin{enumerate}[label=(\arabic*)]
    \item \textbf{Markov Inequality:} $\Pbb(|X| \geq a) \leq a^{-1} \E[|X|]$ for $a > 0$.
    \item \textbf{Chebyshev Inequality:} $\Pbb(|X - \E[X]| \geq a) \leq a^{-2} \Var(X)$ for $a > 0$.
    \item \textbf{H\"{o}lder Inequality:} $|\E[XY]| \leq \|X\|_{L^{p}} \|Y\|_{L^{q}}$ for $p, q \in (1, \infty)$ with $1/p + 1/q = 1$.
    \item \textbf{Jensen Inequality:} If $\varphi$ is convex and $X, \varphi(X) \in L^{1}$, then $\varphi(\E[X]) \leq \E[\varphi(X)]$.
\end{enumerate}
\end{proposition}

\medskip
\noindent\textbf{Why this matters.} These four inequalities are the Swiss army knife of probability theory. Markov and Chebyshev supply upper bounds on tail probabilities---in plain terms, events that deviate far from the mean are rare, the same way that in a class where the average score is 70, the probability of scoring 150 is crushed by a formula involving an 80-point deviation. The intuition behind Jensen is that convex functions amplify uncertainty: evaluating the convex function after taking the expectation always yields a value smaller than or equal to taking the convex function first and then the expectation---much like the volatility drag in investing, where uncertainty itself carries a cost.

\section{Stochastic Processes}

\begin{definition}[Stochastic Process]
Let $T \subseteq [0, \infty)$ be an index set. A stochastic process is a family $\{X_{t}\}_{t \in T}$ of $E$-valued random variables defined on a common probability space. The process is said to be
\begin{enumerate}[label=(\arabic*)]
    \item \textbf{Measurable} if the mapping $(t, \omega) \mapsto X_{t}(\omega)$ is $\B(T) \otimes \F / \calE$-measurable.
    \item \textbf{Adapted} to a filtration $\{\F_{t}\}_{t \in T}$ if $X_{t}$ is $\F_{t}$-measurable for each $t \in T$.
    \item \textbf{Progressively measurable} (when $T = [0, \infty)$) if for each $t \geq 0$, the restriction $X|_{[0,t] \times \Omega}$ is $\B([0,t]) \otimes \F_{t} / \calE$-measurable.
\end{enumerate}
\end{definition}

\medskip
\noindent\textbf{Why this matters.} A stochastic process is a random experiment that unfolds over time. If a random variable is a snapshot, a process is a movie---each frame is a random variable, and consecutive frames are linked by dependence. Adaptedness is the most basic requirement: the value at time $t$ can use only information revealed up to time $t$; it cannot sneak a look at the future---just as a weather forecast cannot be issued today based on tomorrow's actual temperature. Progressive measurability is stronger: it demands that the whole process, viewed as a two-dimensional function of time and randomness, be jointly measurable within any finite time window.

\begin{proposition}
Every right-continuous (or left-continuous) adapted process is progressively measurable.
\end{proposition}

\medskip
\noindent\textbf{Why this matters.} This proposition gives us a practical shortcut: whenever a process is right-continuous---as virtually every financial process we care about is---adaptedness automatically upgrades to progressive measurability. It is like saying ``if the house has a front door, every room inside naturally has a door as well''---right-continuity is a stronger condition than progressive measurability, but precisely because it is stronger, it is easier to verify.

\section{Modes of Convergence}

\begin{definition}
Let $\{X_{n}\}_{n=1}^{\infty}$ be a sequence of random variables. We distinguish:
\begin{enumerate}[label=(\arabic*)]
    \item \textbf{Almost sure convergence:} $X_{n} \to X$ a.s.\ if $\Pbb(\lim_{n} X_{n} = X) = 1$.
    \item \textbf{Convergence in probability:} $X_{n} \xrightarrow{\Pbb} X$ if $\lim_{n} \Pbb(|X_{n} - X| > \varepsilon) = 0$ for every $\varepsilon > 0$.
    \item \textbf{Convergence in $L^{p}$:} $X_{n} \xrightarrow{L^{p}} X$ if $\lim_{n} \|X_{n} - X\|_{L^{p}} = 0$.
    \item \textbf{Convergence in distribution (weak convergence):} $X_{n} \xrightarrow{d} X$ if $\lim_{n} \E[f(X_{n})] = \E[f(X)]$ for every bounded continuous $f$.
\end{enumerate}
\end{definition}

\medskip
\noindent\textbf{Why this matters.} The four convergence modes capture different dimensions of approximation. Almost-sure convergence is the strongest---barring a few disastrously unlucky scenarios, every single sample path eventually converges. Convergence in probability relaxes this: it only demands that the chance of making a large error shrinks to zero. $L^{p}$ convergence focuses on the average size of the error rather than the probability itself. Convergence in distribution is the weakest---it only requires that the statistical portrait of the distribution matches, not that the outcomes of each experiment are close. These four modes are like four criteria for judging an impersonation contest: resemblance in spirit (distribution), resemblance in form (probability), matching gait ($L^{p}$), and perfect recreation of every performance (almost surely).

\begin{proposition}[Relations among Modes]
Almost sure convergence and $L^{p}$ convergence each imply convergence in probability. Convergence in probability implies convergence in distribution. No other implications hold unconditionally.
\end{proposition}

\medskip
\noindent\textbf{Why this matters.} This implication diagram is the road-sign system of probability theory. Almost-sure and $L^{p}$ convergence both point to convergence in probability, which in turn points to convergence in distribution. The reverse implications generally fail---knowing that daily stock prices share the same statistical distribution (convergence in distribution) does not mean yesterday's and today's specific prices are close (convergence in probability). Extra conditions such as uniform integrability or the dominated convergence theorem are needed to travel the reverse direction.

% ──────────────────────────────────────────────────────────────
\chapter{Conditional Expectation and Filtration}

\noindent\textbf{Purpose.} Defines conditional expectation with respect to $\sigma$-algebras and random variables, establishes its fundamental properties (tower, pull-out, Jensen), and introduces filtrations as the mathematical model of information flow.
\textbf{Prerequisites.} Chapter~1 (probability spaces, random variables, $L^{1}$).
\textbf{Used in.} Chapters~3 (martingales),~5 (Markov processes),~8 (generator theory).

\section{Definition of Conditional Expectation}

\begin{definition}[Conditional Expectation with Respect to a $\sigma$-Algebra]
Let $X \in L^{1}(\Omega, \F, \Pbb)$ and let $\calG \subseteq \F$ be a sub-$\sigma$-algebra. The conditional expectation of $X$ given $\calG$, denoted $\E[X \mid \calG]$, is the almost surely unique $\calG$-measurable random variable $Y$ satisfying
\[
\int_{G} Y \, d\Pbb = \int_{G} X \, d\Pbb, \qquad \forall G \in \calG.
\]
Existence follows from the Radon-Nikod\'{y}m theorem: the signed measure $\nu(G) = \int_{G} X \, d\Pbb$ on $(\Omega, \calG)$ is absolutely continuous with respect to $\Pbb|_{\calG}$, so its Radon-Nikod\'{y}m derivative $d\nu/d\Pbb|_{\calG}$ is the conditional expectation.
\end{definition}

\medskip
\noindent\textbf{Why this matters.} Conditional expectation is the best guess based on partial information. Picture yourself estimating the number of people in a half-lit room: you can see the left side clearly but the right side is pitch dark. The conditional expectation uses the count on the left to infer the most likely total for the entire room. Here $\mathcal{G}$ is the region your flashlight illuminates, and $\mathbb{E}[X \mid \mathcal{G}]$ is the unbiased optimal estimate you can make within that illuminated region. The integral equality guarantees that the estimate does not systematically over- or under-shoot.

\begin{definition}[Conditional Expectation with Respect to a Random Variable]
If $Y$ is a random variable and $\sigma(Y)$ denotes the $\sigma$-algebra generated by $Y$, then $\E[X \mid Y] := \E[X \mid \sigma(Y)]$. By the Doob-Dynkin lemma, there exists a Borel-measurable function $g$ such that $\E[X \mid Y] = g(Y)$ almost surely. One writes $\E[X \mid Y = y] := g(y)$.
\end{definition}

\medskip
\noindent\textbf{Why this matters.} The key point of this definition is: once you know the value of $Y$, your best guess for $X$ depends only on that numerical value of $Y$. It is like guessing weight $X$ from height $Y$---you do not need to know how the height was measured; you only need the height number itself. The Doob-Dynkin lemma guarantees that the resulting $g(y)$ is a genuine, usable function.

\begin{definition}[Conditional Probability]
For $A \in \F$, define $\Pbb(A \mid \calG) := \E[\1_{A} \mid \calG]$.
\end{definition}

\medskip
\noindent\textbf{Why this matters.} Conditional probability is nothing more than a special case of conditional expectation---plug the indicator function $\mathbf{1}_A$ into it. Event $A$ either happens (1) or does not (0), and the conditional probability is the best guess within this black-and-white framework. Think of estimating the chance of rain tomorrow ($A$) based on the weather forecast ($\mathcal{G}$): the answer is not 0 or 1, but a number between 0 and 1 that reflects what is known.

\section{Fundamental Properties}

\begin{theorem}[Properties of Conditional Expectation]
Let $X, Y \in L^{1}(\Omega, \F, \Pbb)$ and let $\calG, \calH \subseteq \F$ be sub-$\sigma$-algebras. The following hold almost surely:
\begin{enumerate}[label=(\arabic*)]
    \item \textbf{Linearity:} $\E[aX + bY \mid \calG] = a \E[X \mid \calG] + b \E[Y \mid \calG]$ for $a, b \in \R$.
    \item \textbf{Monotonicity:} $X \leq Y$ a.s.\ $\implies \E[X \mid \calG] \leq \E[Y \mid \calG]$ a.s.
    \item \textbf{Monotone Convergence:} If $0 \leq X_{n} \uparrow X$ a.s.\ and each $X_{n} \in L^{1}$, then $\E[X_{n} \mid \calG] \uparrow \E[X \mid \calG]$ a.s.
    \item \textbf{Dominated Convergence:} If $|X_{n}| \leq Y \in L^{1}$ and $X_{n} \to X$ a.s., then $\E[X_{n} \mid \calG] \to \E[X \mid \calG]$ a.s.\ and in $L^{1}$.
    \item \textbf{Tower Property:} If $\calH \subseteq \calG$, then $\E[\E[X \mid \calG] \mid \calH] = \E[X \mid \calH]$.
    \item \textbf{Pull-Out Property:} If $Y$ is $\calG$-measurable and $Y X \in L^{1}$, then $\E[Y X \mid \calG] = Y \E[X \mid \calG]$.
    \item \textbf{Independence:} If $X$ is independent of $\calG$, then $\E[X \mid \calG] = \E[X]$.
    \item \textbf{Conditional Jensen:} If $\varphi$ is convex and $\varphi(X) \in L^{1}$, then $\varphi(\E[X \mid \calG]) \leq \E[\varphi(X) \mid \calG]$.
\end{enumerate}
\end{theorem}

\medskip
\noindent\textbf{Why this matters.} These eight properties form the toolkit of conditional expectation. The Tower property is especially important: estimating first with finer information and then coarsening to cruder information yields the same result as estimating directly on the crude information---like glancing at the entire exam paper and then closing your eyes to guess the total score, which is equivalent to guessing with your eyes closed from the start. The Pull-out property says that if a quantity is already known (i.e., $\mathcal{G}$-measurable), it can be pulled out of the conditional expectation like a constant. Together these properties endow conditional expectation with the geometric identity of a projection operator, the foundation for martingale theory and Markov processes.

\begin{definition}[Filtration]
A filtration $\{\F_{t}\}_{t \geq 0}$ is an increasing family of sub-$\sigma$-algebras of $\F$:
\[
\F_{s} \subseteq \F_{t} \subseteq \F, \qquad 0 \leq s \leq t < \infty.
\]
The filtration encodes the flow of information: $\F_{t}$ represents the collection of events observable by time $t$.
\end{definition}

\medskip
\noindent\textbf{Why this matters.} A filtration is the mathematical portrait of information revealed over time. Think of a detective solving a case: as time passes, more clues accumulate and the questions that can be answered grow sharper. $\mathcal{F}_t$ is the totality of facts that can be deduced from all clues in the detective's hands at time $t$. The filtration grows larger because information only accumulates---you never ``forget'' what you have already learned.

\begin{definition}[Usual Conditions]
A filtration $\{\F_{t}\}_{t \geq 0}$ satisfies the usual conditions (conditions habituelles) if
\begin{enumerate}[label=(\arabic*)]
    \item \textbf{Right-Continuity:} $\F_{t} = \F_{t+} := \bigcap_{s > t} \F_{s}$ for all $t \geq 0$.
    \item \textbf{Completeness:} $\F_{0}$ contains all $\Pbb$-null sets of $\F$ (and hence every $\F_{t}$ does as well).
\end{enumerate}
\end{definition}

\medskip
\noindent\textbf{Why this matters.} The usual conditions are two technical patches. Right-continuity ensures that information ``right after time $t$'' is the same as information ``at time $t$''---this prevents bizarre information jumps at time boundaries. Completeness labels every probability-zero event as ``known,'' so that null sets can never produce surprises. Though these conditions appear technical, removing them would invalidate a large number of standard theorems, such as the optional stopping theorem.

\begin{definition}[Filtered Probability Space]
The quadruple $(\Omega, \F, \{\F_{t}\}_{t \geq 0}, \Pbb)$ satisfying the usual conditions is called a filtered probability space, or stochastic basis.
\end{definition}

\medskip
\noindent\textbf{Why this matters.} The filtered probability space is the foundation on which everything that follows is built. Just as a building needs a foundation before construction begins, specifying a filtered probability space means we have settled not only which events can occur and with what probability, but also how information unfolds over time. This is the starting posture for studying any dynamic random phenomenon.

\begin{proposition}[Augmentation]
Given any filtration $\{\calG_{t}\}$, the augmented filtration defined by
\[
\F_{t} := \sigma\left(\calG_{t+} \cup \N\right),
\]
where $\N$ is the collection of all $\Pbb$-null sets, satisfies the usual conditions. Throughout this monograph, all filtrations are assumed to satisfy the usual conditions.
\end{proposition}

\medskip
\noindent\textbf{Why this matters.} The augmentation operation is like adding a porch to a small cabin---any filtration that fails the usual conditions can be turned into a valid basis by throwing in the null sets and enforcing right-continuity. This means we never need to worry about whether the usual conditions hold; they can always be repaired automatically through this universal procedure.

% ──────────────────────────────────────────────────────────────
\chapter{Martingales}

\noindent\textbf{Purpose.} Defines martingales, submartingales, and supermartingales in continuous time. States Doob's optional stopping theorem and maximal inequality. These are the probabilistic building blocks for generator theory and optimal stopping.
\textbf{Prerequisites.} Chapters~2 (conditional expectation),~4 (stopping times).
\textbf{Used in.} Chapters~8 (Dynkin's formula),~15 (optimal stopping, Snell envelope).

\section{Definition and First Properties}

\begin{definition}[Martingale]
Let $\{X_{t}\}_{t \geq 0}$ be an adapted, integrable process on a filtered probability space $(\Omega, \F, \{\F_{t}\}, \Pbb)$. The process is a
\begin{enumerate}[label=(\arabic*)]
    \item \textbf{Martingale} if $\E[X_{t} \mid \F_{s}] = X_{s}$ for all $0 \leq s \leq t$,
    \item \textbf{Submartingale} if $\E[X_{t} \mid \F_{s}] \geq X_{s}$ for all $0 \leq s \leq t$,
    \item \textbf{Supermartingale} if $\E[X_{t} \mid \F_{s}] \leq X_{s}$ for all $0 \leq s \leq t$.
\end{enumerate}
\end{definition}

\medskip
\noindent\textbf{Why this matters.} A martingale is the mathematical incarnation of a fair game. If you play a fair roulette wheel at a casino, betting one dollar each round, your wealth process is a martingale---regardless of past wins or losses, your expected wealth after the next round equals your current wealth. A submartingale corresponds to a game tilted in your favor (like cheating at cards), and a supermartingale corresponds to a game tilted against you (like the house edge at a casino). Virtually all of modern asset-pricing theory rests on the single idea that discounted asset prices are martingales.

\begin{proposition}
If $X$ is a martingale, then $\E[X_{t}] = \E[X_{0}]$ for all $t \geq 0$. The analogous statements for sub- and supermartingales involve inequalities.
\end{proposition}

\medskip
\noindent\textbf{Why this matters.} The constant-expectation property of martingales is intuitive: play a fair game for any length of time and your average wealth stays unchanged. But the critical distinction is that the conditional expectation $\mathbb{E}[X_t \mid \mathcal{F}_s] = X_s$ is much stronger than the unconditional statement $\mathbb{E}[X_t] = \mathbb{E}[X_0]$, because it holds along every single information path.

\begin{proposition}[Convex Functions of Martingales]
If $X$ is a martingale and $\varphi$ is convex with $\varphi(X_{t}) \in L^{1}$ for all $t$, then $\{\varphi(X_{t})\}$ is a submartingale.
\end{proposition}

\medskip
\noindent\textbf{Why this matters.} The key intuition here is that convex functions amplify uncertainty. Applying a convex transformation to a fair game turns it into a game biased in your favor. The classic example: the absolute value is convex---a fair one-dimensional random walk (a martingale), when absolute-valued, tends to drift away from zero (a submartingale), because distance from the origin is a one-way street. Similarly, the option payoff $\max(S-K, 0)$ is convex, so even when the stock is a martingale under the risk-neutral measure, the discounted option value becomes a submartingale.

\begin{theorem}[Doob's Optional Stopping Theorem]
Let $X$ be a martingale (resp.\ submartingale) and let $\tau$ be a bounded stopping time. Then $\E[X_{\tau}] = \E[X_{0}]$ (resp.\ $\E[X_{\tau}] \geq \E[X_{0}]$). The result extends to unbounded stopping times under uniform integrability conditions.
\end{theorem}

\medskip
\noindent\textbf{Why this matters.} The optional stopping theorem answers the question: can a player who exits according to a pre-specified strategy beat a fair game? If the player cannot peek into the future (the exit rule is a stopping time) and the game has a definite end (bounded stopping time), the answer is no. Extension to unbounded stopping times requires uniform integrability---much like an exit strategy cannot rely on waiting until the infinite future; a decision must be made eventually.

\begin{theorem}[Doob's Maximal Inequality]
Let $X$ be a submartingale. Then for any $\lambda > 0$,
\[
\Pbb\left(\sup_{0 \leq s \leq t} |X_{s}| \geq \lambda\right) \leq \frac{1}{\lambda} \E[|X_{t}|].
\]
\end{theorem}

\medskip
\noindent\textbf{Why this matters.} Doob's maximal inequality delivers a simple yet powerful bound: the peak value along the entire path of a nonnegative submartingale is controlled entirely by the magnitude at the final time. Think of it this way: a process that can only drift upward cannot first skyrocket and then crash---the highest water level is reined in by the final water level, because a submartingale is biased upward and cannot generate a spike that later collapses.

\section{Continuous-Time Martingale Examples}

\begin{example}[Brownian Motion]
Standard Brownian motion $\{W_{t}\}$ is a martingale: $\E[W_{t} \mid \F_{s}] = W_{s}$.
\end{example}

\begin{example}[Compensated Poisson Process]
If $\{N_{t}\}$ is a Poisson process with intensity $\lambda$, then $\widetilde{N}_{t} = N_{t} - \lambda t$ is a martingale. More generally, for a Poisson random measure $N(dt, dy)$ with compensator $\nu(dy) dt$, the compensated measure $\widetilde{N}(dt, dy) = N(dt, dy) - \nu(dy) dt$ satisfies that integrals with respect to $\widetilde{N}$ are martingales.
\end{example}

\begin{example}[Stochastic Integrals]
If $H \in L^{2}_{\text{loc}}(W)$ (square-integrable and adapted), then $\int_{0}^{t} H_{s} dW_{s}$ is a martingale. This is the fundamental property of the It\^{o} integral.
\end{example}

% ──────────────────────────────────────────────────────────────
\chapter{Stopping Times}

\noindent\textbf{Purpose.} Defines stopping times and the pre-$\tau$ $\sigma$-algebra $\mathcal{F}_{\tau}$. Establishes first hitting times as the canonical example. Stopping times formalize exercise decisions in American options.
\textbf{Prerequisites.} Chapters~2 (filtrations),~3 (martingales).
\textbf{Used in.} Chapters~8 (Dynkin with stopping times),~15 (optimal stopping),~16 (Bermudan DP).

\section{Definition}

\begin{definition}[Stopping Time]
A random variable $\tau: \Omega \to [0, \infty]$ is a stopping time with respect to $\{\F_{t}\}$ if $\{\tau \leq t\} \in \F_{t}$ for all $t \geq 0$. Under right-continuity of the filtration, this is equivalent to $\{\tau < t\} \in \F_{t}$ for all $t$.
\end{definition}

\medskip
\noindent\textbf{Why this matters.} A stopping time is a decision rule that cannot peek into the future. Think of pressing the close-door button in an elevator---you decide ``press now'' or ``wait a bit'' based only on current information, not on whether someone will dash in three seconds later. The condition $\{\tau \leq t\} \in \mathcal{F}_t$ guarantees that the judgment ``has the stopping time occurred yet?'' uses only information available up to time $t$. Every option exercise decision and every stop-loss or take-profit strategy must be a stopping time---otherwise it amounts to insider trading.

\begin{definition}[Pre-$\tau$ $\sigma$-Algebra]
The $\sigma$-algebra of events determined by time $\tau$ is
\[
\F_{\tau} := \left\{A \in \F_{\infty} : A \cap \{\tau \leq t\} \in \F_{t} \text{ for all } t \geq 0\right\},
\]
where $\F_{\infty} := \sigma(\bigcup_{t \geq 0} \F_{t})$.
\end{definition}

\medskip
\noindent\textbf{Why this matters.} $\mathcal{F}_\tau$ contains all historical information up to the moment the stopping time $\tau$ occurs---like everything you know about the elevator the instant before you press the button. This is not a fixed time $t$ but a snapshot of information taken at a random instant that varies with the stopping time. When $\tau$ happens to equal $t$ identically, $\mathcal{F}_\tau$ reduces to the ordinary $\mathcal{F}_t$.

\begin{proposition}[Properties of $\F_{\tau}$]
Under the usual conditions:
\begin{enumerate}[label=(\arabic*)]
    \item $\F_{\tau}$ is a $\sigma$-algebra.
    \item If $\tau \equiv t$ almost surely, then $\F_{\tau} = \F_{t}$.
    \item If $\tau_{1} \leq \tau_{2}$ a.s., then $\F_{\tau_{1}} \subseteq \F_{\tau_{2}}$.
    \item The stopped process $\{X_{\tau \wedge t}\}_{t \geq 0}$ is progressively measurable if $\{X_{t}\}$ is.
\end{enumerate}
\end{proposition}

\medskip
\noindent\textbf{Why this matters.} These properties guarantee the consistency of the information structure at stopping times. The third property is especially intuitive: the later you stop, the more information you accumulate---wait longer, know more. The fourth property (the stopped process remains measurable) is particularly important for computation: it legitimizes integrating and taking expectations of a process stopped at exercise time.

\begin{example}[First Hitting Times]
For a continuous adapted process $\{X_{t}\}$, the first hitting time of a closed set $F$,
\[
\tau_{F} := \inf\{t \geq 0 : X_{t} \in F\},
\]
is a stopping time under the usual conditions. The first hitting time of an open set is a stopping time without additional conditions if the filtration is right-continuous.
\end{example}

% ──────────────────────────────────────────────────────────────
\chapter{Markov Processes and the Markov Property}

\noindent\textbf{Purpose.} Defines Markov processes and the time-homogeneous Markov property. Introduces transition kernels and the Chapman-Kolmogorov equation. This is the structural property that makes generator-based approximation possible.
\textbf{Prerequisites.} Chapters~1 (probability),~2 (conditional expectation).
\textbf{Used in.} Chapters~6 (semigroups),~7 (Feller processes),~8 (generators).

\section{Definition}

\begin{definition}[Markov Process]
A stochastic process $\{X_{t}\}_{t \geq 0}$ with values in a measurable state space $(E, \calE)$ is a Markov process with respect to $\{\F_{t}\}$ if
\[
\E[f(X_{t}) \mid \F_{s}] = \E[f(X_{t}) \mid X_{s}], \qquad 0 \leq s \leq t,
\]
for every bounded measurable function $f: E \to \R$. Equivalently, the conditional distribution of $X_{t}$ given $\F_{s}$ depends only on $X_{s}$:
\[
\Pbb(X_{t} \in A \mid \F_{s}) = \Pbb(X_{t} \in A \mid X_{s}) =: P_{s,t}(X_{s}, A), \quad A \in \calE.
\]
\end{definition}

\medskip
\noindent\textbf{Why this matters.} The core intuition is that a Markov process cares only about the present, not the past. Think of a chess game---your next move depends solely on the current position on the board, not on the sequence of moves that produced it. This property radically simplifies analysis: you do not need to store the entire history; knowing the current state is enough to predict the future.

The family $\{P_{s,t}\}_{0 \leq s \leq t}$ is the \textbf{transition kernel} of the process. For each $s, t$, $P_{s,t}(\cdot, \cdot)$ is a probability kernel: for each fixed $x \in E$, $P_{s,t}(x, \cdot)$ is a probability measure on $(E, \calE)$, and for each fixed $A \in \calE$, $P_{s,t}(\cdot, A)$ is a measurable function on $E$.

\begin{definition}[Time-Homogeneous Markov Process]
A Markov process is \textbf{time-homogeneous} if its transition kernel depends only on the time difference: $P_{s,t}(x, A) = P_{t-s}(x, A)$. In this case, we denote $\{P_{t}\}_{t \geq 0}$ as the transition semigroup. For any bounded measurable $f$ and any $s,t \geq 0$:
\[
\E[f(X_{t+s}) \mid \F_{t}] = \E[f(X_{t+s}) \mid X_{t}] = P_{s}f(X_{t}),
\]
where $P_{t}f(x) := \int_{E} f(y) P_{t}(x, dy)$.
\end{definition}

\medskip
\noindent\textbf{Why this matters.} Time-homogeneity is the mathematical statement that the laws of physics do not change with time. If the relationship between today's temperature and tomorrow's temperature depends only on the number of days between them, not on whether today is Monday or Wednesday, the process is time-homogeneous. For financial applications, time-homogeneity means an option price depends only on the remaining time to expiry, not on calendar date---a dramatic reduction in dimensionality.

% ──────────────────────────────────────────────────────────────
\chapter{Transition Semigroups}

\noindent\textbf{Purpose.} Defines the operator semigroup $\{P_{t}\}$ induced by a Markov process. Establishes the semigroup property $P_{t+s} = P_{t}P_{s}$ and strong continuity---the analytic bridge from the probabilistic description to operator theory.
\textbf{Prerequisites.} Chapter~5 (Markov processes, Chapman-Kolmogorov).
\textbf{Used in.} Chapters~7 (Feller semigroups),~8 (infinitesimal generator),~9 (Hille-Yosida),~11 (CTMC transition matrices).

\begin{definition}[Operator Semigroup]
Let $B(E)$ denote the Banach space of bounded measurable functions on $E$ equipped with the supremum norm $\|f\| = \sup_{x \in E} |f(x)|$. The transition operators $\{P_{t}\}_{t \geq 0}$ defined by
\[
P_{t} f(x) = \E[f(X_{t}) \mid X_{0} = x] = \int_{E} f(y) P_{t}(x, dy)
\]
form a \textbf{one-parameter semigroup} of bounded linear operators on $B(E)$:
\begin{enumerate}[label=(\arabic*)]
    \item \textbf{Identity:} $P_{0} = I$ (the identity operator).
    \item \textbf{Semigroup property:} $P_{t} P_{s} = P_{t+s}$ for all $s,t \geq 0$.
    \item \textbf{Contraction:} $\|P_{t} f\| \leq \|f\|$ for all $t \geq 0$.
    \item \textbf{Positivity:} $f \geq 0 \implies P_{t} f \geq 0$.
\end{enumerate}
\end{definition}

\medskip
\noindent\textbf{Why this matters.} The operator semigroup is the shadow of a Markov process cast onto function space. Instead of tracking the process path at every moment, we track how a function evolves over time. The semigroup property $P_{t+s} = P_t P_s$ is the Markov property expressed at the operator level: evolving $s$ steps and then $t$ more steps is the same as evolving $t+s$ steps in one go. This echoes the additivity of the exponential function $e^{a+b} = e^a e^b$---and indeed the matrix exponential $e^{tQ}$ is the semigroup induced by the generator $Q$.

\begin{theorem}[Chapman-Kolmogorov Equation]
For a time-homogeneous Markov process, the transition probabilities satisfy
\[
P_{t+s}(x, A) = \int_{E} P_{t}(y, A) P_{s}(x, dy), \quad \forall s,t \geq 0, \; \forall A \in \calE.
\]
In operator notation: $P_{t+s} = P_{t} P_{s} = P_{s} P_{t}$. This follows directly from the Markov property and the law of total probability.
\end{theorem}

\medskip
\noindent\textbf{Why this matters.} The Chapman-Kolmogorov equation is the path-concatenation formula for Markov processes. Traveling from $x$ to $A$ in time $t+s$ is equivalent to: first jumping from $x$ to some intermediate state $y$ in time $s$, then jumping from $y$ to $A$ in the remaining $t$, and finally taking a weighted average over all possible $y$. This is like a connecting flight from New York to Tokyo via Anchorage---the statistics of the full journey are composed from the statistics of the two legs.

\begin{definition}[Strongly Continuous Semigroup]
A semigroup $\{P_{t}\}$ on a Banach space $B$ is \textbf{strongly continuous} (or $C_{0}$) if $\lim_{t \to 0} \|P_{t} f - f\| = 0$ for all $f \in B$. For Markov semigroups, strong continuity requires the additional property that functions are mapped into a subspace of $B(E)$ on which the semigroup acts continuously at $t = 0$; this subspace is precisely $C_{0}(E)$ for Feller processes.
\end{definition}

\medskip
\noindent\textbf{Why this matters.} Strong continuity guarantees that the process cannot undergo a catastrophic jump in an arbitrarily short time. If $P_t f$ returns to $f$ in the supremum norm as $t$ approaches zero, the process paths are nearly continuous in the relevant function-space sense. This is the prerequisite for defining the infinitesimal generator---if the limit does not exist, the definition of the generator is meaningless.

% ──────────────────────────────────────────────────────────────
\chapter{Feller Processes}

\noindent\textbf{Purpose.} Defines Feller semigroups and Feller processes---the well-behaved class of Markov processes for which the generator is a closed, densely defined operator on $C_{0}$. States the martingale problem characterization, the central link between generator and process.
\textbf{Prerequisites.} Chapters~5 (Markov),~6 (semigroups).
\textbf{Used in.} Chapters~8 (generator),~9 (Hille-Yosida),~10--14 (CTMC approximation relies on Feller property).

\begin{definition}[Feller Semigroup]
A Markov semigroup $\{P_{t}\}$ on a locally compact separable metric space $E$ is a \textbf{Feller semigroup} if
\begin{enumerate}[label=(\arabic*)]
    \item $P_{t}$ maps $C_{0}(E)$ into $C_{0}(E)$ for all $t \geq 0$,
    \item $\lim_{t \to 0} \|P_{t} f - f\| = 0$ for all $f \in C_{0}(E)$.
\end{enumerate}
A Markov process whose transition semigroup satisfies these properties is called a \textbf{Feller process}.
\end{definition}

\medskip
\noindent\textbf{Why this matters.} A Feller process is a well-behaved Markov process. The first condition ($C_0$ invariance) guarantees that if a function decays to zero in the far reaches of the state space---like the strength of an electric field fading with distance from the charge---then its conditional expectation after time $t$ still respects that decay; the process cannot spontaneously manufacture influence at infinity. The second condition (strong continuity) ensures that the generator can be defined. Most practical financial models (GBM, OU, CIR, CEV) are Feller processes.

\begin{theorem}[Martingale Problem Characterization]
Let $E$ be a complete separable metric space and let $\calL: \calD(\calL) \subseteq C_{0}(E) \to C_{0}(E)$ be a linear operator with $\calD(\calL)$ dense. The following are equivalent:
\begin{enumerate}[label=(\arabic*)]
    \item $\calL$ generates a Feller semigroup.
    \item For each $\nu \in \calP(E)$ (the space of probability measures on $E$), there exists a unique (in law) progressively measurable process $\{X_{t}\}$ with initial distribution $\nu$ such that for every $f \in \calD(\calL)$,
    \[
    M_{t}^{f} := f(X_{t}) - f(X_{0}) - \int_{0}^{t} \calL f(X_{s}) \, ds
    \]
    is an $\{\F_{t}^{X}\}$-martingale.
\end{enumerate}
\end{theorem}

\medskip
\noindent\textbf{Why this matters.} The martingale-problem characterization is the central methodological bridge of this monograph. It tells us that to construct a Feller process, we need not solve a stochastic differential equation; we only need to find an operator $\mathcal{L}$ such that $f(X_t) - \int_0^t \mathcal{L}f(X_s) ds$ becomes a martingale. Put differently: $\mathcal{L}$ describes the average rate of change of $f(X_t)$ after stripping away the martingale noise. The CTMC approximation that follows exploits exactly this equivalence: we do not approximate the process paths directly; we approximate its generator $\mathcal{L}$.

This equivalence is the central structural result linking the analytic object $\calL$ to the probabilistic object $\{X_{t}\}$. It justifies the approximation strategy: if one constructs operators $\calL^{(N)}$ converging to $\calL$, then the associated processes converge in law.

% ═══════════════════════════════════════════════════════════════
% PART III — GENERATOR THEORY
% ═══════════════════════════════════════════════════════════════
\part{Generator Theory}

\begin{quote}
\itshape The infinitesimal generator is the central object of this monograph. Its construction, computation, and approximation form the conceptual thread connecting every subsequent chapter. The generator encodes the entire local behavior of the process. It is the gateway from specification of dynamics (via SDEs or their parameterizations) to the construction of approximating CTMCs.
\end{quote}

% ──────────────────────────────────────────────────────────────
\chapter{The Infinitesimal Generator}

\section{Definition and Domain}

\begin{definition}[Infinitesimal Generator]
Let $\{T_{t}\}_{t \geq 0}$ be a $C_{0}$-semigroup of contractions on a Banach space $B$. The infinitesimal generator $\calL$ of $\{T_{t}\}$ is the linear operator defined by
\[
\calL f := \lim_{t \to 0} \frac{T_{t} f - f}{t},
\]
with domain
\[
\calD(\calL) := \left\{f \in B : \lim_{t \to 0} \frac{T_{t} f - f}{t} \text{ exists in the norm topology of } B\right\}.
\]
For a Feller process $\{X_{t}\}$ with state space $E$, the generator on $C_{0}(E)$ takes the concrete form
\[
\calL f(x) = \lim_{t \to 0} \frac{\E[f(X_{t}) \mid X_{0} = x] - f(x)}{t},
\]
for $f \in \calD(\calL)$, where the limit is in the supremum norm over $E$.
\end{definition}

\medskip
\noindent\textbf{Why this matters.} The generator describes the average behavior of a process over an infinitesimally short time interval. Think of a stochastic process as a car driving through fog: the drift $\mu$ is the tendency of the steering wheel (which direction it biases toward), and the diffusion $\sigma^2$ is the random jitter of the gas pedal (how hard it wobbles). The generator condenses both into a single operator, forming the central bridge between process description and CTMC approximation.

\begin{proposition}[Properties of the Generator]
\begin{enumerate}[label=(\arabic*)]
    \item $\calD(\calL)$ is a dense linear subspace of $B$.
    \item $\calL$ is a closed operator: if $f_{n} \in \calD(\calL)$, $f_{n} \to f$, and $\calL f_{n} \to g$, then $f \in \calD(\calL)$ and $\calL f = g$.
    \item For $f \in \calD(\calL)$, the mapping $t \mapsto T_{t} f$ is continuously differentiable, $T_{t} f \in \calD(\calL)$, and
    \[
    \frac{d}{dt} T_{t} f = \calL T_{t} f = T_{t} \calL f.
    \]
\end{enumerate}
\end{proposition}

\medskip
\noindent\textbf{Why this matters.} The third property---$\frac{d}{dt} T_t f = \mathcal{L} T_t f$---is the reason for the name ``generator.'' The semigroup $T_t$ satisfies a differential equation analogous to $\frac{dy}{dt} = ay$, with the generator $\mathcal{L}$ playing the role of the coefficient $a$. Closedness ensures that limit operations stay within the domain---just as the limit of continuous functions is continuous, the limit of generators is a generator.

\section{Remarks on Domains of Generators}

\subsection{Why the Domain Matters}

The generator $\mathcal{L}$ is an \emph{unbounded} operator on $B$. Unlike a bounded operator (e.g., a finite matrix), $\mathcal{L}$ cannot be applied to every element of $B$; it is defined only on a proper dense subspace $\mathcal{D}(\mathcal{L}) \subsetneq B$. Writing $\mathcal{L} f$ without verifying $f \in \mathcal{D}(\mathcal{L})$ is mathematically meaningless---much as one cannot differentiate a non-differentiable function and claim the result is a derivative in the classical sense.

In the CTMC approximation framework, this issue is partially hidden because the approximating operator $Q^{(N)}$ is a finite matrix and therefore bounded (its domain is all of $\R^{N}$). The convergence $Q^{(N)} f \to \mathcal{L} f$ in the Trotter--Kato theorem is only asserted for $f$ belonging to a \textbf{core} of $\mathcal{L}$. The identification of a suitable core is thus a prerequisite for any rigorous convergence statement.

\subsection{Common Cores in Applications}

For diffusion processes (Assumption~D), the standard core is $C_{c}^{\infty}(\R)$, the space of compactly supported smooth functions. On this core, the generator takes the familiar differential form:
\[
\mathcal{L} f(x) = \mu(x) f'(x) + \frac{1}{2} \sigma^{2}(x) f''(x), \qquad f \in C_{c}^{\infty}(\R).
\]
The full domain $\mathcal{D}(\mathcal{L})$ is the larger space
\[
\mathcal{D}(\mathcal{L}) = \{f \in C_{0}(\R) \cap C^{2}(\R) : \mu f' + \tfrac{1}{2}\sigma^{2} f'' \in C_{0}(\R)\},
\]
but for approximation purposes it suffices to work on the core.

For jump-diffusions (Assumption~J), the core is again $C_{c}^{\infty}(\R)$, and the generator expression includes the integral term:
\[
\mathcal{L} f(x) = \mu(x) f'(x) + \frac{1}{2} \sigma^{2}(x) f''(x) + \int_{\R} [f(x+h(x,y)) - f(x)] \nu(dy).
\]
The non-local integral term maps $C_{c}^{\infty}(\R)$ into $C_{0}(\R)$ under Assumption~J.

For regime-switching processes (Assumption~R), the core is $C_{c}^{\infty}(\R \times \{1,\ldots,M\})$, i.e., functions that are smooth in $x$ for each fixed regime $m$ and vanish as $|x| \to \infty$.

\subsection{Treatment in This Monograph}

Throughout this monograph, we adopt the following convention:
\begin{itemize}
    \item All generator expressions are understood to hold on the core $C_{c}^{\infty}$ (or its discrete analog) unless explicitly stated otherwise.
    \item When we write $\mathcal{L} f(x) = \mu(x) f'(x) + \frac{1}{2}\sigma^{2}(x) f''(x)$, we mean this identity holds for all $f \in C_{c}^{\infty}(\R)$ and all $x \in \R$.
    \item The CTMC generator $Q^{(N)}$ acts on all grid functions $\mathbf{f} \in \R^{N}$ (its domain is the whole space). Convergence $Q^{(N)} f \to \mathcal{L} f$ is proved for test functions $f$ restricted to the grid and sufficiently smooth.
    \item Full domain characterization (e.g., via the graph norm or resolvent estimates) is not required for the approximation results in Parts~IV--VI; the core is sufficient.
\end{itemize}

\begin{remark}
The choice of core is not unique. For diffusions, $C_{c}^{2}(\R)$ or the Schwartz space $\mathcal{S}(\R)$ also serve as cores. The essential requirement is that the core be dense in $C_{0}(\R)$ and invariant under the resolvent. Our use of $C_{c}^{\infty}(\R)$ follows Ethier \& Kurtz (1986).
\end{remark}

\begin{remark}
A common pitfall is to write $\mathcal{L} f$ for $f$ being a payoff function such as $\phi(x) = (K-x)^{+}$, which is not in $C_{c}^{\infty}$ (it is not even $C^{1}$ at the strike). Such expressions must be interpreted in the viscosity sense (see Chapter~\ref{ch:viscosity}) or via smooth approximation. The CTMC pricing recursion never requires evaluating $\mathcal{L} \phi$ directly.
\end{remark}

\begin{theorem}[Dynkin's Formula]
Let $X$ be a Feller process with generator $\calL$, and let $f \in \calD(\calL)$. Then for any $t \geq 0$,
\[
\E[f(X_{t}) \mid X_{0} = x] = f(x) + \E\left[\int_{0}^{t} \calL f(X_{s}) \, ds \mid X_{0} = x\right].
\]
More generally, if $\tau$ is a stopping time with $\E[\tau] < \infty$, then
\[
\E[f(X_{\tau}) \mid X_{0} = x] = f(x) + \E\left[\int_{0}^{\tau} \calL f(X_{s}) \, ds \mid X_{0} = x\right].
\]
\end{theorem}

\medskip
\noindent\textbf{Why this matters.} Dynkin's formula is the fundamental theorem of stochastic calculus---it decomposes the expectation of $f(X_t)$ into the starting value $f(x)$ plus the accumulated effect of the generator along the path. The analogy is the Newton-Leibniz formula $f(b) - f(a) = \int_a^b f'(x) dx$. Here $\mathcal{L} f$ plays the role of the derivative, and the stopping-time version generalizes the variable-limit integral. This formula is the departure point for all subsequent pricing, including Feynman-Kac and dynamic programming.

\begin{corollary}[Martingale Property of Compensated $f(X)$]
For $f \in \calD(\calL)$, the process
\[
M_{t}^{f} := f(X_{t}) - f(X_{0}) - \int_{0}^{t} \calL f(X_{s}) \, ds
\]
is a martingale.
\end{corollary}

\medskip
\noindent\textbf{Why this matters.} This corollary is the martingale-form expression of Dynkin's formula: after subtracting the integrated generator from $f(X_t)$, what remains is the unpredictable martingale noise. This supplies a practical test of whether a candidate operator really is the generator of a process: check whether $f(X_t) - \int_0^t \mathcal{L}f(X_s)ds$ is a martingale.

\begin{corollary}[Generator as Derivative of the Semigroup at Zero]
For $f \in \calD(\calL)$, $\calL f = \left.\frac{d}{dt} P_{t} f\right|_{t=0}$.
\end{corollary}

\medskip
\noindent\textbf{Why this matters.} This corollary gives the instantaneous-rate interpretation of the generator---it is the rate at which the semigroup changes the function $f$ at the very moment $t=0$. Just as velocity is the derivative of position and acceleration is the derivative of velocity, $\mathcal{L} f$ measures the initial drift of $f$ under the motion of the process.

\begin{theorem}[Kolmogorov Backward Equation]
For $f \in \calD(\calL)$ and $t \geq 0$,
\[
\frac{\partial}{\partial t} P_{t} f = \calL P_{t} f, \qquad P_{0}f = f.
\]
\end{theorem}

\medskip
\noindent\textbf{Why this matters.} The Kolmogorov backward equation tells us: to compute how the expected value of $f$ changes over time starting from $x$, apply the generator to the starting point $x$. It is called ``backward'' because the generator acts on the departure position $x$. This connects directly to the Black-Scholes equation in option pricing---replace the generator with its concrete form (e.g., $\mu f' + \frac{1}{2}\sigma^2 f''$) and you obtain a parabolic PDE.

\begin{theorem}[Feynman-Kac Formula]
Let $r: E \to [0, \infty)$ be a bounded measurable discount function. For $f \in B(E)$, define
\[
u(t, x) := \E\left[\exp\left(-\int_{0}^{t} r(X_{s}) ds\right) f(X_{t}) \mid X_{0} = x\right].
\]
Then $u$ satisfies (in a suitable sense)
\[
\frac{\partial u}{\partial t} = \calL u - r u, \qquad u(0, x) = f(x).
\]
\end{theorem}

\medskip
\noindent\textbf{Why this matters.} The Feynman-Kac formula is the cornerstone of financial pricing. It turns an option price expressed as an expectation (with discounting) into the solution of a PDE. Here $r$ is the interest rate (discount factor) and $\mathcal{L}$ is the generator of the underlying asset. In intuitive terms: in a risk-neutral world, the option price satisfies $\frac{\partial V}{\partial t} + \mathcal{L}V - rV = 0$---the time decay of the price equals the diffusion effect minus discounting. Every equivalence between Monte Carlo and PDE methods rests on this foundation.

\begin{proof}[Proof Sketch]
From Proposition~8.1.2, $d/dt \, \E[f(X_{t}) \mid X_{0}=x] = \E[\calL f(X_{t}) \mid X_{0}=x]$. Integrating from $0$ to $t$ and applying Fubini's theorem yields the result. The extension to stopping times requires the optional stopping theorem applied to the martingale $M_{t}^{f} = f(X_{t}) - f(X_{0}) - \int_{0}^{t} \calL f(X_{s}) ds$.
\end{proof}

\begin{corollary}[Martingale Property of Compensated $f(X)$]
For $f \in \calD(\calL)$, the process
\[
M_{t}^{f} := f(X_{t}) - f(X_{0}) - \int_{0}^{t} \calL f(X_{s}) \, ds
\]
is a martingale.
\end{corollary}

\begin{corollary}[Generator as Derivative of the Semigroup at Zero]
For $f \in \calD(\calL)$, $\calL f = \left.\frac{d}{dt} P_{t} f\right|_{t=0}$.
\end{corollary}

\section{The Kolmogorov Equations}

\begin{theorem}[Kolmogorov Backward Equation]
For $f \in \calD(\calL)$ and $t \geq 0$,
\[
\frac{\partial}{\partial t} P_{t} f = \calL P_{t} f, \qquad P_{0}f = f.
\]
\end{theorem}

\begin{theorem}[Feynman-Kac Formula]
Let $r: E \to [0, \infty)$ be a bounded measurable discount function. For $f \in B(E)$, define
\[
u(t, x) := \E\left[\exp\left(-\int_{0}^{t} r(X_{s}) ds\right) f(X_{t}) \mid X_{0} = x\right].
\]
Then $u$ satisfies (in a suitable sense)
\[
\frac{\partial u}{\partial t} = \calL u - r u, \qquad u(0, x) = f(x).
\]
\end{theorem}

% ──────────────────────────────────────────────────────────────
\chapter{Generators of Specific Processes}

\noindent\textbf{Purpose.} Computes the generator for concrete processes: Brownian motion, It\^{o} diffusions, jump-diffusions, and regime-switching processes. This is where the abstract generator $\mathcal{L}$ acquires its parametric form $\mu f' + \frac{1}{2}\sigma^{2} f''$ (plus jump integral).
\textbf{Prerequisites.} Chapter~8 (generator definition and domain).
\textbf{Used in.} Chapters~12--14 (moment matching requires the concrete form of $\mathcal{L}$).

\section{Brownian Motion}

\begin{definition}[Standard Brownian Motion]
A standard one-dimensional Brownian motion $\{W_{t}\}$ satisfies $W_{0} = 0$, has independent increments with $W_{t} - W_{s} \sim \mathcal{N}(0, t-s)$, and has continuous sample paths almost surely.
\end{definition}

\medskip
\noindent\textbf{Why this matters.} Brownian motion is the atom of the stochastic-process universe. Just as physicists derive the behavior of complex molecules from the properties of the hydrogen atom, probabilists build every elaborate diffusion and jump model as a variant or composite of Brownian motion. The erratic dance of pollen grains in water, the minute-by-minute fluctuations of stock prices, and thermal noise in electrical circuits all take Brownian motion as their fundamental model.

\begin{theorem}[Generator of Brownian Motion]
On the core $C_{c}^{\infty}(\R)$,
\[
\calL_{\text{BM}} f(x) = \frac{1}{2} f''(x).
\]
\end{theorem}

\medskip
\noindent\textbf{Why this matters.} The generator $\frac{1}{2}f''$ reveals a core truth: Brownian motion is pure diffusion---no drift, no jumps, only random jitter. The second derivative $f''$ measures the curvature of $f$: where $f''$ is positive (convex), diffusion pushes the expectation upward; where concave, it pushes downward. This is precisely the core operator of the heat equation.

\begin{definition}[It\^{o} Diffusion]
A one-dimensional It\^{o} diffusion is the unique strong solution to the SDE
\[
dX_{t} = \mu(X_{t}) \, dt + \sigma(X_{t}) \, dW_{t},
\]
where $\mu, \sigma: \R \to \R$ satisfy a Lipschitz condition: $|\mu(x) - \mu(y)| + |\sigma(x) - \sigma(y)| \leq K|x-y|$, and a linear growth condition: $|\mu(x)| + |\sigma(x)| \leq K(1 + |x|)$.
\end{definition}

\medskip
\noindent\textbf{Why this matters.} An It\^{o} diffusion is Brownian motion with a steering wheel. $\mu$ is the direction of the wind (drift), and $\sigma$ is the strength of the turbulence (diffusion). The Lipschitz condition guarantees uniqueness of the solution---same starting point, same wind field, the car follows the same path. The linear-growth condition prevents the solution from exploding to infinity in finite time---it imposes a reasonable speed limit on the vehicle.

\begin{theorem}[Generator of an It\^{o} Diffusion]
Under the above conditions, on $C_{c}^{\infty}(\R)$,
\[
\boxed{\calL f(x) = \mu(x) f'(x) + \frac{1}{2} \sigma^{2}(x) f''(x)}.
\]
\end{theorem}

\medskip
\noindent\textbf{Why this matters.} This is one of the most important formulas in the entire monograph. The generator splits into two pieces: $\mu f'$ is the drift-driven tendency---the first derivative captures the slope of $f$ in the drift direction; $\frac{1}{2}\sigma^2 f''$ is the diffusion-driven spreading---the second derivative captures how the curvature of $f$ is amplified by the random shaking. Every CTMC construction that follows amounts to matching these two moments (first and second) exactly on a discrete grid.

\begin{corollary}[Drift and Diffusion Coefficients from the Generator]
The coefficients are recovered by applying the generator to specific test functions:
\[
\mu(x) = \calL(\id)(x), \qquad \sigma^{2}(x) = \calL((\id - x)^{2})(x),
\]
where $\id(x) = x$. This observation provides the moment-matching intuition exploited in Chapter~12.
\end{corollary}

\medskip
\noindent\textbf{Why this matters.} This corollary reveals the essence of moment matching: $\mu$ and $\sigma^2$ correspond to the values obtained by applying the generator to the test functions $x$ and $(x - x_0)^2$, respectively. This explains why CTMC construction requires matching only the first and second moments---these two numbers fully determine the behavior of the generator; no third- or higher-order information is needed.

\section{Jump-Diffusion Processes}

\begin{definition}[Jump-Diffusion]
Let $N(dt, dy)$ be a Poisson random measure on $[0, \infty) \times \R$ with compensator $\nu(dy) dt$. A jump-diffusion process satisfies
\[
dX_{t} = \mu(X_{t-}) \, dt + \sigma(X_{t-}) \, dW_{t} + \int_{\R} h(X_{t-}, y) \, N(dt, dy),
\]
where $h: \R \times \R \to \R$ is measurable. The compensated Poisson random measure is $\widetilde{N}(dt, dy) := N(dt, dy) - \nu(dy) dt$.
\end{definition}

\medskip
\noindent\textbf{Why this matters.} A jump-diffusion process adds sudden leaps on top of continuous diffusion. Think of a car driving smoothly most of the time but occasionally hitting a speed bump---the ride is smooth drift plus random vibration, punctuated by sharp jolts at random intervals. In finance, jump-diffusion models capture price jumps triggered by news events or central-bank announcements, fitting the fat tails of real market data better than pure-diffusion models.

\begin{theorem}[Generator of a Jump-Diffusion]
Under suitable integrability conditions,
\[
\boxed{\calL f(x) = \mu(x) f'(x) + \frac{1}{2} \sigma^{2}(x) f''(x) + \int_{\R} \left[f(x + h(x,y)) - f(x)\right] \nu(dy)}.
\]
If one uses the martingale-compensated formulation, the jump integral takes the alternative form
\[
\int_{\R} \left[f(x + h(x,y)) - f(x) - h(x,y) f'(x)\right] \nu(dy).
\]
\end{theorem}

\medskip
\noindent\textbf{Why this matters.} The jump-diffusion generator consists of three building blocks: the drift part ($\mu f'$), the continuous-diffusion part ($\frac{1}{2}\sigma^2 f''$), and the jump part (the integral). The intuition for the integral is: for each possible jump size $y$, compute the change in $f$ caused by the jump, $f(x+h) - f(x)$, and take a weighted average using the jump intensity $\nu(dy)$. This is analogous to tallying the damage from random machine breakdowns---each failure mode has its own frequency and severity distribution.

\begin{proof}
We proceed step by step.

\textbf{Step 1 --- It\^{o}'s formula.} For $f \in C^{2}$,
\[
\begin{aligned}
f(X_{t}) - f(x) &= \int_{0}^{t} f'(X_{s-}) \mu(X_{s-}) ds + \int_{0}^{t} f'(X_{s-}) \sigma(X_{s-}) dW_{s} \\
&\quad + \frac{1}{2} \int_{0}^{t} f''(X_{s-}) \sigma^{2}(X_{s-}) ds \\
&\quad + \int_{0}^{t} \int_{\R} \left[f(X_{s-} + h(X_{s-}, y)) - f(X_{s-})\right] N(ds, dy).
\end{aligned}
\]

\textbf{Step 2 --- Compensator decomposition.} Decompose the Poisson integral using $\widetilde{N}(ds, dy) = N(ds, dy) - \nu(dy) ds$:
\[
\int_{0}^{t} \int_{\R} [f(X_{s-} + h) - f(X_{s-})] N(ds, dy) = \underbrace{\int_{0}^{t} \int_{\R} [f(X_{s-} + h) - f(X_{s-})] \widetilde{N}(ds, dy)}_{=: M_{t}^{\text{jump}}} + \int_{0}^{t} \int_{\R} [f(X_{s-} + h) - f(X_{s-})] \nu(dy) ds.
\]
The $\widetilde{N}$-integral is a martingale; taking conditional expectation given $X_{0} = x$, it vanishes.

\textbf{Step 3 --- Conditional expectation.} Collecting the non-martingale (finite-variation) terms:
\[
\begin{aligned}
\E[f(X_{t}) \mid X_{0} = x] - f(x) &= \E\left[\int_{0}^{t} \mu(X_{s-}) f'(X_{s-}) ds \mid X_{0} = x\right] \\
&\quad + \frac{1}{2} \E\left[\int_{0}^{t} \sigma^{2}(X_{s-}) f''(X_{s-}) ds \mid X_{0} = x\right] \\
&\quad + \E\left[\int_{0}^{t} \int_{\R} [f(X_{s-} + h(X_{s-}, y)) - f(X_{s-})] \nu(dy) ds \mid X_{0} = x\right].
\end{aligned}
\]

\textbf{Step 4 --- Passage to the limit.} Divide by $t$ and take $t \to 0$. By right-continuity of $X_{t}$, for almost every $\omega$, $X_{s-} \to x$ as $s \to 0$. The dominated convergence theorem yields:
\[
\lim_{t \to 0} \frac{\E[f(X_{t}) \mid X_{0} = x] - f(x)}{t} = \mu(x) f'(x) + \frac{1}{2} \sigma^{2}(x) f''(x) + \int_{\R} [f(x + h(x,y)) - f(x)] \nu(dy).
\]
This is precisely $\calL f(x)$.
\end{proof}

\section{Regime-Switching Diffusion}

\begin{definition}[Regime-Switching Process]
Let $\{R_{t}\}$ be a finite-state CTMC on $\M := \{1, 2, \ldots, M\}$ with generator matrix $Q^{R} = (\lambda_{mn})_{m,n \in \M}$ satisfying $\lambda_{mn} \geq 0$ for $m \neq n$ and $\sum_{n} \lambda_{mn} = 0$. Conditionally on $R_{t} = m$, the process $\{X_{t}\}$ follows
\[
dX_{t} = \mu_{m}(X_{t}) dt + \sigma_{m}(X_{t}) dW_{t}.
\]
The joint process $(X_{t}, R_{t})$ takes values in $E = \R \times \M$.
\end{definition}

\medskip
\noindent\textbf{Why this matters.} Regime-switching models capture the reality that market conditions can flip abruptly. The weather has sunny, cloudy, and rainy modes, each with its own temperature dynamics. Financial markets likewise have bull regimes (low volatility, positive drift) and bear regimes (high volatility, negative drift) as distinct weather patterns. A hidden Markov chain $R_t$ drives the switching between these modes, fitting the sudden shifts in variance seen in market data far better than a single-regime model.

\begin{theorem}[Generator of a Regime-Switching Process]
For $f \in C_{c}^{\infty}(\R \times \M)$,
\[
\boxed{\calL f(x, m) = \mu_{m}(x) \frac{\partial f}{\partial x}(x, m) + \frac{1}{2} \sigma_{m}^{2}(x) \frac{\partial^{2} f}{\partial x^{2}}(x, m) + \sum_{n \neq m} \lambda_{mn} \big[f(x, n) - f(x, m)\big]}.
\]
\end{theorem}

\medskip
\noindent\textbf{Why this matters.} This generator superimposes three ingredients: the diffusion generator within regime $m$ ($\mu_m f_x + \frac{1}{2}\sigma_m^2 f_{xx}$), plus the regime-switching effect $\sum_n \lambda_{mn}[f(x,n) - f(x,m)]$. The intuition for the switching term is straightforward: when currently in regime $m$, the process jumps to regime $n$ at rate $\lambda_{mn}$, and each jump contributes a change of $f(x,n) - f(x,m)$. This two-layer structure---spatial diffusion plus regime switching---is among the most complex generators in the monograph, and it is the reason the CTMC representation requires a block-matrix format.

% ──────────────────────────────────────────────────────────────
\chapter{The Hille-Yosida Theorem}
\label{ch:hille-yosida}

\noindent\textbf{Purpose.} States the Hille-Yosida and Trotter-Kato theorems, which provide the rigorous foundation for generator-to-semigroup convergence. The Trotter-Kato theorem is the mathematical justification for the CTMC approximation strategy.
\textbf{Prerequisites.} Chapters~6 (semigroups),~8 (generator).
\textbf{Used in.} Chapter~12 (convergence chain),~18 (main convergence theorem). This is the theoretical backbone of Part~III.

\section{Statement of the Hille-Yosida Theorem}

\begin{theorem}[Hille-Yosida]
A linear operator $\calL$ on a Banach space $B$ is the infinitesimal generator of a $C_{0}$-semigroup of contractions $\{T_{t}\}$ if and only if
\begin{enumerate}[label=(\arabic*)]
    \item $\calD(\calL)$ is dense in $B$.
    \item $\calL$ is closed.
    \item The resolvent set $\rho(\calL)$ contains $(0, \infty)$, and for all $\lambda > 0$,
    \[
    \|R_{\lambda}\| \leq \frac{1}{\lambda}, \qquad R_{\lambda} := (\lambda I - \calL)^{-1}.
    \]
\end{enumerate}
Moreover, the semigroup is uniquely determined by $\calL$, and is given by the exponential formula
\[
T_{t} f = \lim_{n \to \infty} \left(I - \frac{t}{n}\calL\right)^{-n} f, \qquad f \in B,
\]
which is the Yosida approximation.
\end{theorem}

\medskip
\noindent\textbf{Why this matters.} The Hille-Yosida theorem is the constitution of generator theory---it prescribes exactly which operators can legally generate a semigroup. The three conditions (density, closedness, and the resolvent bound) are not technical footnotes; they guarantee from first principles that the semigroup exists and is unique. The Yosida approximation formula $(I - \frac{t}{n}\mathcal{L})^{-n}$ provides a computable path from generator to semigroup---discretizing continuous time into $n$ infinitesimal steps.

\begin{corollary}[Uniqueness]
The generator $\calL$ uniquely determines the semigroup $\{T_{t}\}$, and hence the finite-dimensional distributions of the associated Markov process. Two Markov processes whose generators agree on a common core have the same law.
\end{corollary}

\medskip
\noindent\textbf{Why this matters.} This corollary is the ultimate justification for the CTMC approximation strategy: if we can construct a CTMC generator $Q^{(N)}$ that agrees with $\mathcal{L}$ on a function core, then by the Hille-Yosida guarantee, all finite-dimensional distributions of the two processes coincide in the limit. In other words: two organisms with the same genetic code---whether continuous or discrete---produce macroscopically indistinguishable trajectories.

\begin{theorem}[Hille-Yosida for Feller Processes]
Let $E$ be a locally compact separable metric space. A linear operator $\calL$ on $C_{0}(E)$ generates a Feller semigroup if and only if it satisfies the Hille-Yosida conditions and additionally satisfies the \textbf{positive maximum principle}: for any $f \in \calD(\calL)$ and $x_{0} \in E$ such that $f(x_{0}) = \sup_{x \in E} f(x) \geq 0$, one has $\calL f(x_{0}) \leq 0$.
\end{theorem}

\medskip
\noindent\textbf{Why this matters.} The positive maximum principle is a deceptively simple but profound condition: if a function attains a global nonnegative maximum at a point, the generator must be non-positive there---meaning the process cannot continue climbing from a position that has already reached the peak. This captures a natural constraint on probabilistic processes: once a stock hits an all-time high, its expected infinitesimal change cannot remain positive, for that would violate the boundedness of probability measures.

\begin{theorem}[Trotter-Kato Semigroup Convergence]
Let $\{\calL^{(n)}\}$ be a sequence of generators of $C_{0}$-semigroups $\{T_{t}^{(n)}\}$ on a Banach space $B$, and let $\calL$ generate $\{T_{t}\}$. Assume there exists a core $\calD_{0} \subseteq \calD(\calL)$ such that:
\begin{enumerate}[label=(\arabic*)]
    \item For each $f \in \calD_{0}$, there exist $f_{n} \in \calD(\calL^{(n)})$ with $f_{n} \to f$ and $\calL^{(n)} f_{n} \to \calL f$ in $B$.
    \item The semigroups $\{T_{t}^{(n)}\}$ are uniformly bounded in $n$.
\end{enumerate}
Then for each $f \in B$ and $T > 0$,
\[
\lim_{n \to \infty} \sup_{t \in [0,T]} \|T_{t}^{(n)} f - T_{t} f\| = 0.
\]
\end{theorem}

\medskip
\noindent\textbf{Why this matters.} The Trotter-Kato theorem is the convergence guarantee of CTMC approximation theory. It answers a critical technical question: does convergence of generators imply convergence of semigroups? The answer is yes---as long as the generators converge pointwise on a dense core, the semigroups converge uniformly over any finite time horizon. This is the final link in the logical chain that takes us from matching generators to matching prices.

\begin{corollary}[CTMC Convergence]
If one constructs a sequence of CTMC generators $Q^{(N)}$ such that $Q^{(N)} f \to \calL f$ for all $f$ in a core of $\calL$, then the Trotter-Kato theorem guarantees that the CTMC transition matrices $P_{t}^{(N)} = e^{t Q^{(N)}}$ converge, in the appropriate operator topology, to the semigroup $\{P_{t}\}$ of the target process. This is the fundamental justification for the CTMC approximation methodology.
\end{corollary}

\medskip
\noindent\textbf{Why this matters.} The key idea is teaching a frog to mimic a swimming fish. The fish (diffusion process) glides smoothly and continuously; the frog (CTMC) can only hop between lily pads. But as long as the average direction and the randomness of each hop match those of the fish, from a distance the frog's trajectory looks nearly identical to the fish's. This is the essence of moment matching---not replicating every step, but matching the statistical signature.

\begin{theorem}[Semigroup Error Bound]
Let $\calL$ and $\calL^{(N)}$ generate $C_{0}$-semigroups $\{P_{t}\}$ and $\{P_{t}^{(N)}\}$ respectively. If there exists $M < \infty$ and a dense set $\calD_{0}$ such that $\|\calL^{(N)} f - \calL f\| \leq M \omega(N) \|f\|_{\calD}$ for all $f \in \calD_{0}$ (where $\|\cdot\|_{\calD}$ is a graph norm), then for all $t \geq 0$,
\[
\|P_{t}^{(N)} f - P_{t} f\| \leq t \cdot M \omega(N) \|f\|_{\calD}.
\]
In the CTMC context, $\omega(N)$ is typically $\mathcal{O}(h^{2})$ where $h$ is the grid spacing, yielding a global error estimate of $\mathcal{O}(T h^{2})$ over a finite time horizon $[0, T]$.
\end{theorem}

\medskip
\noindent\textbf{Why this matters.} This error bound provides the second piece of critical information---not just that convergence happens, but how fast. The local generator error of $\mathcal{O}(h^2)$ propagates linearly over time to a global error of $\mathcal{O}(T h^2)$. How small must $h$ be to price an option to four decimal places? How long is $T$? This bound answers those questions quantitatively. Note that the error grows linearly with $T$---for long-dated instruments such as LEAPS options ($T > 2$ years), one must either refine the grid or adopt higher-order schemes.

\begin{corollary}[Uniqueness]
The generator $\calL$ uniquely determines the semigroup $\{T_{t}\}$, and hence the finite-dimensional distributions of the associated Markov process. Two Markov processes whose generators agree on a common core have the same law.
\end{corollary}

\section{Application to Feller Processes}

\begin{theorem}[Hille-Yosida for Feller Processes]
Let $E$ be a locally compact separable metric space. A linear operator $\calL$ on $C_{0}(E)$ generates a Feller semigroup if and only if it satisfies the Hille-Yosida conditions and additionally satisfies the \textbf{positive maximum principle}: for any $f \in \calD(\calL)$ and $x_{0} \in E$ such that $f(x_{0}) = \sup_{x \in E} f(x) \geq 0$, one has $\calL f(x_{0}) \leq 0$.
\end{theorem}

\section{Trotter-Kato Theorem: Convergence of Semigroups}

\begin{theorem}[Weak Convergence of CTMC]
\label{thm:weak-convergence-ctmc}
Under the assumptions of Theorem~\ref{thm:generator-approximation}, let $X^{(n)}$ be the CTMC with initial distribution converging to $X_{0}$.

Then the sequence of processes converges in distribution:
\[
X^{(n)} \Rightarrow X \quad \text{in } D([0,T],E),
\]
where $X$ is the diffusion process generated by $(\mathcal{L}, D(\mathcal{L}))$.

The convergence holds in the Skorokhod topology.
\end{theorem}

\begin{lemma}[Tightness via Compact Containment and Aldous Condition]
\label{lem:tightness-ctmc}
Let $X^{(n)}$ be the CTMCs in Theorem~\ref{thm:weak-convergence-ctmc}. Suppose that for every $T>0$ and $\eta>0$ there exists a compact set $K\subset E$ such that
\[
\inf_n \mathbb{P}\left(X_t^{(n)}\in K \text{ for all } 0\le t\le T\right) \ge 1-\eta,
\]
and that Aldous' condition holds: for every $T>0$, $\eta>0$, and $\varepsilon>0$,
\[
\lim_{\delta\downarrow0}\sup_n\sup_{\tau\le T}\sup_{0\le\theta\le\delta}
\mathbb{P}\left(d(X_{\tau+\theta}^{(n)},X_{\tau}^{(n)})>\varepsilon\right)=0,
\]
where the middle supremum is over stopping times bounded by $T$. Then $\{X^{(n)}\}$ is tight in $D([0,T],E)$.
\end{lemma}

\begin{proof}[Verification principle]
For the diffusion and finite-activity jump models considered here, compact containment follows from the Lyapunov moment bounds in Chapter~\ref{app:lyapunov}; bounded jump intensity and local moment matching control the probability of a large displacement over $[\tau,\tau+\theta]$. These two estimates imply Aldous' condition, hence tightness. Semigroup convergence identifies every subsequential limit as the martingale-problem solution for $(\mathcal{L},D(\mathcal{L}))$; uniqueness of that martingale problem yields the stated weak convergence.
\end{proof}

\begin{corollary}[CTMC Convergence]
If one constructs a sequence of CTMC generators $Q^{(N)}$ such that $Q^{(N)} f \to \calL f$ for all $f$ in a core of $\calL$, then the Trotter-Kato theorem guarantees that the CTMC transition matrices $P_{t}^{(N)} = e^{t Q^{(N)}}$ converge, in the appropriate operator topology, to the semigroup $\{P_{t}\}$ of the target process. This is the fundamental justification for the CTMC approximation methodology.
\end{corollary}

\section{Error Estimates via the Hille-Yosida Framework}

\begin{theorem}[Semigroup Error Bound]
Let $\calL$ and $\calL^{(N)}$ generate $C_{0}$-semigroups $\{P_{t}\}$ and $\{P_{t}^{(N)}\}$ respectively. If there exists $M < \infty$ and a dense set $\calD_{0}$ such that $\|\calL^{(N)} f - \calL f\| \leq M \omega(N) \|f\|_{\calD}$ for all $f \in \calD_{0}$ (where $\|\cdot\|_{\calD}$ is a graph norm), then for all $t \geq 0$,
\[
\|P_{t}^{(N)} f - P_{t} f\| \leq t \cdot M \omega(N) \|f\|_{\calD}.
\]
In the CTMC context, $\omega(N)$ is typically $\mathcal{O}(h^{2})$ where $h$ is the grid spacing, yielding a global error estimate of $\mathcal{O}(T h^{2})$ over a finite time horizon $[0, T]$.
\end{theorem}

% ═══════════════════════════════════════════════════════════════
% PART IV — CTMC APPROXIMATION
% ═══════════════════════════════════════════════════════════════
\part{CTMC Approximation}
\label{part:ctmc}

\begin{quote}
\itshape This part contains the constructive core of the approximation theory. The central idea, running through every chapter that follows, is the mapping from generator to generator:

\[
\boxed{\mathcal{L} \;\longrightarrow\; Q^{(N)}}
\]

Given the specification of a Markov process via its infinitesimal generator $\mathcal{L}$ (Parts~II--III), we construct a finite-state CTMC whose generator matrix $Q^{(N)}$ approximates $\mathcal{L}$ in the sense required by the Trotter--Kato theorem. The method is local moment matching via Taylor expansion. Every step is formulated as a theorem or a constructive proposition; no heuristic justification is offered. This mapping $\mathcal{L} \mapsto Q^{(N)}$ is the computational pivot of the entire monograph: once $Q^{(N)}$ is built, the pricing pipeline (Parts~V--VI) operates on it without any further reference to the continuous process.
\end{quote}

% ──────────────────────────────────────────────────────────────
\chapter{Finite-State CTMCs and Generator Matrices}

\noindent\textbf{Purpose.} This chapter defines the mathematical objects that constitute the CTMC approximation: the finite-state chain, its generator matrix $Q$, and its transition matrix $P_{t} = e^{tQ}$. These definitions translate the continuous generator $\mathcal{L}$ (Part~III) into a concrete finite-dimensional linear-algebraic framework. Every subsequent construction, convergence theorem, and pricing algorithm operates on the objects defined here.

\textbf{Prerequisites.} Part~I (probability), Part~II (Markov property).

\textbf{Used in.} Chapters~12--14 (moment matching and model-specific constructions), Chapter~16 (Bermudan DP), Chapter~18 (main pricing theorem).

\section{Definition of Finite-State CTMCs}

\begin{definition}[Finite-State CTMC]
Let $\calS^{(N)} = \{x_{1}, x_{2}, \ldots, x_{N}\} \subset \R$ be a finite set. A continuous-time Markov chain with state space $\calS^{(N)}$ is a stochastic process $\{\widetilde{X}_{t}^{(N)}\}_{t \geq 0}$ such that for all $i, j \in \{1, \ldots, N\}$, $i \neq j$,
\[
\Pbb\left(\widetilde{X}_{t+h}^{(N)} = x_{j} \mid \widetilde{X}_{t}^{(N)} = x_{i}\right) = q_{ij} h + o(h), \qquad h \to 0,
\]
where $q_{ij} \geq 0$ is the instantaneous transition rate from $x_{i}$ to $x_{j}$.
\end{definition}

\medskip
\noindent\textbf{Why this matters.} A finite-state CTMC is a continuous-time random walk on a discrete space. Picture a frog hopping among $N$ lily pads; the timing of each hop is random, but each pair of adjacent pads has a fixed instantaneous rate $q_{ij}$. The quantity $q_{ij} h$ is the probability of jumping from pad $i$ to pad $j$ in a very short interval $h$. The only difference from a diffusion is that the frog can stand only on the discrete pads, never on the water between them.

\begin{definition}[Generator Matrix]
The $N \times N$ matrix $Q^{(N)} = (q_{ij})_{i,j=1}^{N}$ is a valid CTMC generator matrix if
\begin{align}
\text{(C1)} &\quad q_{ij} \geq 0 \quad \forall i \neq j, \\
\text{(C2)} &\quad q_{ii} = -\sum_{j \neq i} q_{ij} \quad \forall i \quad (\text{equivalently}, \sum_{j} q_{ij} = 0).
\end{align}
\end{definition}

\medskip
\noindent\textbf{Why this matters.} Conditions (C1) and (C2) are the ID card of a CTMC generator matrix. Nonnegative off-diagonal entries ($q_{ij} \ge 0$) guarantee that probabilities cannot be negative; zero row sums guarantee total probability conservation---from any state, the sum of probability flows to all possible destinations is zero (incoming equals outgoing, adjusted by the diagonal). These two conditions are the continuous-time analogue of "each row of a transition matrix sums to 1."

\begin{theorem}[Transition Matrix as Matrix Exponential]
For a finite-state CTMC with generator $Q^{(N)}$, the transition probability matrix over a time interval $t \geq 0$ is
\[
P_{t}^{(N)} := \left(\Pbb(\widetilde{X}_{t}^{(N)} = x_{j} \mid \widetilde{X}_{0}^{(N)} = x_{i})\right)_{i,j} = e^{t Q^{(N)}} := \sum_{k=0}^{\infty} \frac{(t Q^{(N)})^{k}}{k!}.
\]
\end{theorem}

\medskip
\noindent\textbf{Why this matters.} The key point is that the deceptively simple formula $P_t = e^{tQ}$ is the engine of the entire CTMC pricing framework. For continuous processes we are not so lucky---the transition density of a diffusion rarely has a closed form. But for a CTMC, the transition matrix is just the matrix exponential, computable efficiently via Pad\'e approximation or uniformization. This means that once $Q$ is constructed, pricing reduces to a linear-algebra exercise.

\begin{remark}[Numerical Stability of the Matrix Exponential]
Direct computation of $e^{tQ}$ via the Taylor series or Pad\'e approximation becomes ill-conditioned for large $t$ or large grid size $N$, because $\|Q\| \sim 1/h^{2} \sim N^{2}$. The condition number $\kappa(e^{tQ}) \leq e^{t \max_{i}|q_{ii}|}$ grows exponentially with $t$ and quadratically with $N$. For practical pricing with $N \geq 500$ and $T \geq 1$, uniformization (see the Spectral Analysis appendix) is numerically mandatory: it replaces $e^{tQ}$ with a Poisson-weighted sum of discrete-step stochastic matrices $S^{m}$, where $S = I + Q/\Lambda$ has entries in $[0,1]$, avoiding catastrophic cancellation.
\end{remark}

\begin{definition}
For a function $f: \calS^{(N)} \to \R$, identified with the column vector $\mathbf{f} = (f(x_{1}), \ldots, f(x_{N}))^{\top}$, the action of $Q^{(N)}$ on $\mathbf{f}$ is
\[
(Q^{(N)} \mathbf{f})_{i} = \sum_{j=1}^{N} q_{ij} f(x_{j}) = \sum_{j \neq i} q_{ij} [f(x_{j}) - f(x_{i})].
\]
\end{definition}

\medskip
\noindent\textbf{Why this matters.} The second form of this definition, $Q \mathbf{f} = \sum_{j \neq i} q_{ij}[f(x_j) - f(x_i)]$, is the core modeling tool. It expresses the generator as the function change caused by a jump from $i$ to $j$, weighted by the jump rate. When we later Taylor-expand $f(x_j) - f(x_i) = f'(x_i)(x_j-x_i) + \frac{1}{2}f''(x_i)(x_j-x_i)^2 + \cdots$, the first moment $\sum q_{ij}(x_j-x_i)$ and second moment $\sum q_{ij}(x_j-x_i)^2$ emerge directly as the central quantities that determine how well $Q$ approximates $\mathcal{L}$.

\begin{definition}[Grid]
Choose a bounded interval $[a, b] \subset \R$ and partition it into $N$ subintervals. Let $x_{i}$ be a representative point (e.g., the midpoint) of the $i$-th subinterval. Define the local spacings
\[
h_{i} := x_{i+1} - x_{i}, \qquad h_{i-1} := x_{i} - x_{i-1}.
\]
For a uniform grid, $h_{i} = h$ for all $i$. The grid fineness parameter is $h := \max_{i} h_{i}$, and the approximation is studied in the limit $h \to 0$ (equivalently, $N \to \infty$).
\end{definition}

\medskip
\noindent\textbf{Why this matters.} The grid is the bridge that discretizes continuous space. Partitioning $[a,b]$ into $N$ segments is like turning a continuous landscape painting into a pixel image---the larger $N$ (the smaller $h$), the closer the pixels match the original. The grid spacing $h_i$ need not be uniform: near the strike price, where price action is most intense, a denser grid captures more detail; far out in the tails, a coarser grid suffices.

\begin{definition}[Generator Matrix]
The $N \times N$ matrix $Q^{(N)} = (q_{ij})_{i,j=1}^{N}$ is a valid CTMC generator matrix if
\begin{align}
\text{(C1)} &\quad q_{ij} \geq 0 \quad \forall i \neq j, \\
\text{(C2)} &\quad q_{ii} = -\sum_{j \neq i} q_{ij} \quad \forall i \quad (\text{equivalently}, \sum_{j} q_{ij} = 0).
\end{align}
\end{definition}

\begin{theorem}[Transition Matrix as Matrix Exponential]
For a finite-state CTMC with generator $Q^{(N)}$, the transition probability matrix over a time interval $t \geq 0$ is
\[
P_{t}^{(N)} := \left(\Pbb(\widetilde{X}_{t}^{(N)} = x_{j} \mid \widetilde{X}_{0}^{(N)} = x_{i})\right)_{i,j} = e^{t Q^{(N)}} := \sum_{k=0}^{\infty} \frac{(t Q^{(N)})^{k}}{k!}.
\]
\end{theorem}

\begin{proof}
The Kolmogorov forward equation $dP_{t}/dt = P_{t} Q$, $P_{0} = I$, is a linear matrix ODE whose unique solution is $P_{t} = e^{tQ}$. Convergence of the series is absolute for all $t$ since $Q$ is a finite matrix.
\end{proof}

\section{Generator as an Operator on Functions}

\begin{definition}
For a function $f: \calS^{(N)} \to \R$, identified with the column vector $\mathbf{f} = (f(x_{1}), \ldots, f(x_{N}))^{\top}$, the action of $Q^{(N)}$ on $\mathbf{f}$ is
\[
(Q^{(N)} \mathbf{f})_{i} = \sum_{j=1}^{N} q_{ij} f(x_{j}) = \sum_{j \neq i} q_{ij} [f(x_{j}) - f(x_{i})].
\]
\end{definition}

\section{State-Space Discretization}

\begin{definition}[Grid]
Choose a bounded interval $[a, b] \subset \R$ and partition it into $N$ subintervals. Let $x_{i}$ be a representative point (e.g., the midpoint) of the $i$-th subinterval. Define the local spacings
\[
h_{i} := x_{i+1} - x_{i}, \qquad h_{i-1} := x_{i} - x_{i-1}.
\]
For a uniform grid, $h_{i} = h$ for all $i$. The grid fineness parameter is $h := \max_{i} h_{i}$, and the approximation is studied in the limit $h \to 0$ (equivalently, $N \to \infty$).
\end{definition}

% ──────────────────────────────────────────────────────────────
\chapter{Local Consistency and Moment Matching}

\noindent\textbf{Purpose.} Translates the generator approximation problem into a concrete algebraic condition: equate the first two instantaneous moments of the CTMC to $\mu$ and $\sigma^{2}$. Derives the local truncation error and the convergence chain from moment matching to weak convergence.
\textbf{Prerequisites.} Chapters~9 (generator formulas),~11 (finite-state CTMCs).
\textbf{Used in.} Chapters~13 (diffusion construction),~14 (jump/RS construction).

\section{The Generator Matching Principle}

Let $\{X_{t}\}$ be a Feller process with generator $\calL$, and let $\{\widetilde{X}_{t}^{(N)}\}$ be a CTMC on $\calS^{(N)}$ with generator $Q^{(N)}$. The CTMC is a consistent approximation to $X$ if
\[
(Q^{(N)} f)(x_{i}) \to (\calL f)(x_{i})
\]
as $h \to 0$, for all $f$ in a core of $\calL$ and for all grid points $x_{i}$.

\begin{lemma}[Taylor Expansion with Remainder]
Let $f \in C^{3}(\R)$ and $x_{i} \in \calS^{(N)}$. For any $x_{j} \in \calS^{(N)}$,
\[
f(x_{j}) - f(x_{i}) = f'(x_{i})(x_{j} - x_{i}) + \frac{1}{2} f''(x_{i})(x_{j} - x_{i})^{2} + \frac{1}{6} f'''(\xi_{ij})(x_{j} - x_{i})^{3},
\]
for some $\xi_{ij}$ between $x_{i}$ and $x_{j}$.
\end{lemma}

Inserting this into the CTMC generator action:
\[
(Q^{(N)} f)(x_{i}) = f'(x_{i}) \sum_{j \neq i} q_{ij} (x_{j} - x_{i}) + \frac{1}{2} f''(x_{i}) \sum_{j \neq i} q_{ij} (x_{j} - x_{i})^{2} + R_{i}^{(3)}(f),
\]
where $R_{i}^{(3)}(f)$ is the remainder involving third derivatives.

\section{First and Second Moment Matching}

\begin{definition}[First Moment Matching --- Drift]
\[
\boxed{\sum_{j=1}^{N} q_{ij} (x_{j} - x_{i}) = \mu(x_{i})}.
\]
\end{definition}

\medskip
\noindent\textbf{Why this matters.} The first-moment matching condition ensures that the average drift direction of the CTMC agrees with the target diffusion. Think of a price process as a leaf drifting in the wind: $\mu(x_i)$ is the instantaneous wind speed at position $x_i$. The CTMC adjusts the left and right jump rates $q_{i,i\pm 1}$ so that the weighted average displacement equals the wind speed---each step is a discrete hop, but the average trend matches the continuous drift.

\begin{definition}[Second Moment Matching --- Diffusion]
\[
\boxed{\sum_{j=1}^{N} q_{ij} (x_{j} - x_{i})^{2} = \sigma^{2}(x_{i})}.
\]
\end{definition}

\medskip
\noindent\textbf{Why this matters.} The second-moment matching condition ensures that the random jitter of the CTMC matches the target diffusion. $\sigma^2(x_i)$ is the diffusion coefficient---it measures how much uncertainty each step carries. The CTMC adjusts the absolute magnitude of the jump rates so that the weighted average of squared displacements equals the continuous-diffusion variance. Together, the first and second moments equip the CTMC with two control knobs: a drift steering wheel and a volatility gas pedal.

The second condition matches the instantaneous variance of the CTMC to the diffusion coefficient of the continuous process. These two conditions are sufficient for generator convergence because the error in the Taylor approximation is of order $\mathcal{O}(h^{3})$ multiplied by $q_{ij}$, which scales as $\mathcal{O}(1/h^{2})$, yielding an overall local truncation error of $\mathcal{O}(h)$ for the third-order term, with the first- and second-order terms canceling exactly.

\section{From Moment Matching to Weak Convergence}

It is essential to distinguish the \emph{construction} of the CTMC from the \emph{convergence} guarantee. The following chain must be established; no single step follows automatically from the previous:

\[
\boxed{\begin{array}{ccccccc}
\text{Moment Matching} & \longrightarrow & \text{Local Consistency} & \longrightarrow & \text{Generator Convergence} & \longrightarrow & \text{Weak Convergence} \\[4pt]
\text{(construction)} & & \text{(pointwise } Q^{(N)}f \to \mathcal{L}f\text{)} & & \text{(Trotter--Kato)} & & \text{(process law)}
\end{array}}
\]

\subsection{Moment Matching $\Rightarrow$ Local Consistency}

Moment matching is the \emph{construction principle}, not a convergence result. By equating the first two instantaneous moments of the CTMC to the drift and diffusion coefficients,
\[
\sum_{j} q_{ij}(x_{j} - x_{i}) = \mu(x_{i}), \qquad
\sum_{j} q_{ij}(x_{j} - x_{i})^{2} = \sigma^{2}(x_{i}),
\]
and inserting a Taylor expansion of $f$, one obtains
\[
(Q^{(N)} f)(x_{i}) = \mu(x_{i}) f'(x_{i}) + \tfrac{1}{2} \sigma^{2}(x_{i}) f''(x_{i}) + \mathcal{O}(h^{2}) = (\mathcal{L} f)(x_{i}) + \mathcal{O}(h^{2}),
\]
for every smooth test function $f$. This is \textbf{local consistency}: at each grid point, the discrete generator approximates the continuous generator pointwise with a local error of order $h^{2}$.

\begin{remark}
Moment matching alone does \emph{not} imply weak convergence. It guarantees that the drift and diffusion of the CTMC match those of the target process at the infinitesimal level, but the CTMC could still produce qualitatively different sample-path behavior (e.g., it could be explosive, or it could visit regions of the state space that the continuous process never reaches). Additional conditions---uniform boundedness of the semigroup, non-negativity of rates, appropriate boundary treatment---are required.
\end{remark}

\subsection{Local Consistency $\Rightarrow$ Generator Convergence}

Local consistency (pointwise $Q^{(N)} f \to \mathcal{L} f$) on a dense core of $\mathcal{L}$, together with a uniform bound on the CTMC semigroups $P_{t}^{(N)} = e^{t Q^{(N)}}$, implies convergence of the generators to $\mathcal{L}$ in the extended limit sense. This is the input required by the Trotter--Kato theorem (Chapter~\ref{ch:hille-yosida}).

\subsection{Generator Convergence $\Rightarrow$ Semigroup Convergence}

The Trotter--Kato theorem guarantees that generator convergence (on a core, with uniform boundedness) implies strong convergence of the corresponding semigroups:
\[
\|P_{t}^{(N)} f - P_{t} f\| \to 0, \qquad \forall f \in C_{0}(\R),
\]
uniformly on compact time intervals. This is a statement about the transition operators.

\subsection{Semigroup Convergence $\Rightarrow$ Weak Convergence}

Semigroup convergence implies convergence of finite-dimensional distributions. Specifically, for any $0 \leq t_{1} < \cdots < t_{k}$ and any bounded continuous $g: \R^{k} \to \R$,
\[
\E[g(\widetilde{X}_{t_{1}}^{(N)}, \ldots, \widetilde{X}_{t_{k}}^{(N)})] \to \E[g(X_{t_{1}}, \ldots, X_{t_{k}})].
\]
This is precisely \textbf{weak convergence} of the CTMC process to the target process (Ethier \& Kurtz, 1986, Theorem~4.2.5). Together with the uniform boundedness of the Bermudan dynamic programming operator, this finally yields convergence of option prices.

\subsection{Why This Chain Is Necessary}

Attempting to claim ``moment matching implies pricing convergence'' directly is a category error. Moment matching is a constructive tool; it provides the matrix $Q^{(N)}$. The logical chain from $Q^{(N)}$ to the American option price traverses several mathematical domains:
\begin{itemize}
    \item \textbf{Operator theory:} Generator convergence (Trotter--Kato) is an analytic statement about linear operators on Banach spaces. It does not involve probability measures directly.
    \item \textbf{Probability theory:} Weak convergence is a statement about laws of stochastic processes. It follows from operator convergence via the martingale problem characterization or the Ethier--Kurtz theory.
    \item \textbf{Optimal stopping:} Even with weak convergence, convergence of American option prices requires additional verification---the Snell envelope must converge, and the exercise boundary must be well-approximated. This is addressed in Part~V.
\end{itemize}
No single step in the chain is automatic; each is verified in the corresponding chapter of this monograph.

\subsection{Summary of Convergence Arguments by Model Class}

\begin{center}
\begin{tabular}{>{\raggedright\arraybackslash}p{3cm}>{\raggedright\arraybackslash}p{3cm}>{\raggedright\arraybackslash}p{3cm}>{\raggedright\arraybackslash}p{3cm}}
\toprule
\textbf{Model} & \textbf{Local Consistency} & \textbf{Generator Conv.} & \textbf{Weak Conv.} \\
\midrule
Diffusion (D) & Taylor $\mathcal{O}(h^{2})$ & Trotter--Kato & Ethier--Kurtz \\
Jump-Diffusion (J) & Taylor $+$ jump quadrature & Trotter--Kato (combined) & Ethier--Kurtz \\
Regime-Switching (R) & Block Taylor & Trotter--Kato (block) & Ethier--Kurtz \\
\bottomrule
\end{tabular}
\end{center}

% ──────────────────────────────────────────────────────────────
\chapter{Diffusion Approximation}

\noindent\textbf{Purpose.} This chapter translates the abstract moment matching conditions of Chapter~12 into an explicit construction: a tridiagonal generator matrix $Q^{(N)}$ whose entries are closed-form algebraic expressions in the drift $\mu$ and diffusion coefficient $\sigma^{2}$. This is the computational workhorse of the entire CTMC framework. The formulas derived here are directly implementable and produce a valid (non-negative, row-sum-zero) generator matrix.

\textbf{Prerequisites.} Chapters~10--12.

\textbf{Used in.} Chapters~14 (jump-diffusion extension), 15--16 (optimal stopping/pricing). The uniform-grid formulas are the starting point for all numerical implementations.

\section{Tridiagonal Generator Structure}

For a pure diffusion process, the CTMC is a birth-death process:
\[
Q^{(N)} = \begin{pmatrix}
-q_{1}^{+} & q_{1}^{+} & 0 & \cdots & 0 \\
q_{2}^{-} & -(q_{2}^{+} + q_{2}^{-}) & q_{2}^{+} & \cdots & 0 \\
0 & q_{3}^{-} & -(q_{3}^{+} + q_{3}^{-}) & \cdots & 0 \\
\vdots & & \ddots & \ddots & \vdots \\
0 & 0 & \cdots & q_{N}^{-} & -q_{N}^{-}
\end{pmatrix},
\]
where $q_{i}^{+} := q_{i,i+1}$ and $q_{i}^{-} := q_{i,i-1}$.

\section{The Moment Matching Linear System}

For a nonuniform grid, the moment matching conditions reduce to the $2 \times 2$ system:
\[
\begin{cases}
q_{i}^{+} h_{i} - q_{i}^{-} h_{i-1} = \mu(x_{i}), \\[4pt]
q_{i}^{+} h_{i}^{2} + q_{i}^{-} h_{i-1}^{2} = \sigma^{2}(x_{i}).
\end{cases}
\]

\begin{proposition}[Unique Solution of the Moment Matching System]
For each interior grid point $i = 2, \ldots, N-1$, the system has the unique solution
\[
\boxed{\begin{aligned}
q_{i}^{+} &= \frac{h_{i-1} \mu(x_{i}) + \sigma^{2}(x_{i})}{h_{i}(h_{i} + h_{i-1})}, \\[8pt]
q_{i}^{-} &= \frac{-h_{i} \mu(x_{i}) + \sigma^{2}(x_{i})}{h_{i-1}(h_{i} + h_{i-1})}.
\end{aligned}}
\]
\end{proposition}

\medskip
\noindent\textbf{Why this matters.} This formula is the production shop of CTMC modeling. Given any $\mu(x)$ and $\sigma^2(x)$, plug them into these two formulas and the jump rates at every grid point emerge directly. Note that $\sigma^2$ always appears positively in the numerator---diffusion always contributes positive rate; the sign of $\mu$ determines direction: positive drift boosts the rightward rate, negative drift boosts the leftward rate. If $|\mu|$ is so large that a rate turns negative, the upwind scheme must be invoked as a correction.

\begin{corollary}[Uniform Grid Simplification]
For a uniform grid $h_{i} = h_{i-1} = h$,
\[
\boxed{q_{i}^{+} = \frac{h \mu(x_{i}) + \sigma^{2}(x_{i})}{2h^{2}}, \qquad
q_{i}^{-} = \frac{-h \mu(x_{i}) + \sigma^{2}(x_{i})}{2h^{2}}}.
\]
\end{corollary}

\medskip
\noindent\textbf{Why this matters.} The uniform-grid formula reveals the relationship between rates and $h$ more clearly: the diffusion term scales as $1/h^2$, while the drift term scales as $1/h$. This means the diffusion effect explodes on a fine grid (tiny steps require enormous jump frequency to generate the same total variance), while the drift effect grows only moderately. This is why very small $h$ produces enormous rate values---the CTMC compensates for tiny step sizes with extremely high frequency.

\begin{remark}[Positivity of Transition Rates]
The standard formulas produce valid (non-negative) transition rates if and only if $\sigma^{2}(x_{i}) \geq h |\mu(x_{i})|$ at every grid point. When this condition is violated, the upwind scheme (Appendix~\ref{app:upwind}) must be invoked to restore non-negativity. For driftless processes ($\mu \equiv 0$, as with CEV futures), positivity is guaranteed for all $h$. The condition also imposes a lower bound on grid refinement: one cannot take $h \to 0$ independently of verifying $\sigma^{2} \geq h|\mu|$ at each step.
\end{remark}

\begin{theorem}[Generator Approximation Theorem]
\label{thm:generator-approximation}
Let $(X_t)_{t\ge 0}$ be a Feller process on $E \subset \mathbb{R}^d$ with infinitesimal generator $(\mathcal{L}, D(\mathcal{L}))$.

Let $(X^{(n)}_t)_{t\ge 0}$ be a sequence of CTMCs with generator matrices $Q^{(n)}$ defined on grids $E_n$ with mesh size $h_n \to 0$.

Assume:
\begin{enumerate}[label=(A\arabic*)]
    \item \textbf{Consistency of domain.} There exists a core $\mathcal{D} \subset D(\mathcal{L}) \cap C_0^2(E)$ such that $\mathcal{D}$ is dense in $C_0(E)$.
    \item \textbf{Local operator consistency.} For every $f \in \mathcal{D}$,
    \[
    \sup_{x \in E_n} \left| Q^{(n)} f(x) - \mathcal{L} f(x) \right| \le C_f \, \varepsilon_n, \quad \varepsilon_n \to 0.
    \]
    \item \textbf{Conservativity.}
    \[
    \sum_{y \in E_n} Q^{(n)}(x,y) = 0.
    \]
\end{enumerate}
Then for every $f \in C_0(E)$,
\[
\lim_{n \to \infty} \| P^{(n)}_t f - P_t f \|_\infty = 0, \quad \forall t \ge 0,
\]
where $P_t^{(n)} = e^{tQ^{(n)}}$ is the finite-dimensional CTMC semigroup and $P_t$ is the strongly continuous contraction semigroup generated by $(\mathcal{L},D(\mathcal{L}))$ under the Hille--Yosida assumptions stated in the Operator Framework.
\end{theorem}

\medskip
\noindent\textbf{Why this matters.} This theorem upgrades the argument from moment matching to operator convergence. Moment matching remains the constructive mechanism for verifying (A2), but the convergence claim itself is now phrased at the generator and semigroup level, which is the appropriate Kurtz--Ethier framework.

\begin{proof}[Proof sketch]
Assumption~(A1) identifies a dense core on which the limiting generator is determined, so convergence on $\mathcal{D}$ extends to the closed operator by graph-norm closure. Assumption~(A2) gives convergence of the discrete generators to $\mathcal{L}$ on that core, and Assumption~(A3) ensures that each $Q^{(n)}$ is conservative and hence generates a contraction Markov semigroup on $C_0(E_n)$. The Hille--Yosida theorem gives the limiting contraction semigroup, and the Trotter--Kato theorem lifts local operator convergence to strong convergence of semigroups on $C_0(E)$.
\end{proof}

\begin{corollary}[Uniform Grid Simplification]
For a uniform grid $h_{i} = h_{i-1} = h$,
\[
\boxed{q_{i}^{+} = \frac{h \mu(x_{i}) + \sigma^{2}(x_{i})}{2h^{2}}, \qquad
q_{i}^{-} = \frac{-h \mu(x_{i}) + \sigma^{2}(x_{i})}{2h^{2}}}.
\]
\end{corollary}

\section{Local Truncation Error Analysis}

\begin{proposition}[Local Consistency Order for Diffusion CTMC]
Let $\mu, \sigma^{2} \in C^{2}(\R)$ and $f \in C^{4}(\R)$. For the tridiagonal generator constructed above,
\[
|(Q^{(N)} f)(x_{i}) - (\calL f)(x_{i})| \leq C_{f}(x_{i}) \cdot h^{2},
\]
where $C_{f}(x_{i})$ depends on $\max(|f'''|, |f^{(4)}|)$, $|\mu|$, $|\sigma^{2}|$, and the grid geometry.
\end{proposition}

\begin{proof}
We carry out the Taylor expansion explicitly on a uniform grid ($h_{i} = h_{i-1} = h$).

\textbf{Step 1 --- Taylor expansions.} Expand $f(x_{i} \pm h)$ to fourth order:
\[
\begin{aligned}
f(x_{i} + h) &= f(x_{i}) + h f'(x_{i}) + \frac{h^{2}}{2!} f''(x_{i}) + \frac{h^{3}}{3!} f'''(x_{i}) + \frac{h^{4}}{4!} f^{(4)}(\xi_{+}), \\
f(x_{i} - h) &= f(x_{i}) - h f'(x_{i}) + \frac{h^{2}}{2!} f''(x_{i}) - \frac{h^{3}}{3!} f'''(x_{i}) + \frac{h^{4}}{4!} f^{(4)}(\xi_{-}),
\end{aligned}
\]
for some $\xi_{+} \in (x_{i}, x_{i} + h)$ and $\xi_{-} \in (x_{i} - h, x_{i})$.

\textbf{Step 2 --- Compute the CTMC generator action.}
\[
\begin{aligned}
(Q^{(N)}f)(x_{i}) &= q_{i}^{+}[f(x_{i}+h) - f(x_{i})] + q_{i}^{-}[f(x_{i}-h) - f(x_{i})] \\
&= f'(x_{i}) \cdot (q_{i}^{+} h - q_{i}^{-} h)
   + f''(x_{i}) \cdot \tfrac{1}{2}(q_{i}^{+} h^{2} + q_{i}^{-} h^{2}) \\
&\quad + f'''(x_{i}) \cdot \tfrac{1}{6}(q_{i}^{+} h^{3} - q_{i}^{-} h^{3})
   + \tfrac{h^{4}}{24}\left[q_{i}^{+} f^{(4)}(\xi_{+}) + q_{i}^{-} f^{(4)}(\xi_{-})\right].
\end{aligned}
\]

\textbf{Step 3 --- Apply moment matching.} By construction, $q_{i}^{+} h - q_{i}^{-} h = \mu(x_{i})$ and $q_{i}^{+} h^{2} + q_{i}^{-} h^{2} = \sigma^{2}(x_{i})$. Substituting:
\[
(Q^{(N)}f)(x_{i}) = \mu(x_{i}) f'(x_{i}) + \tfrac{1}{2} \sigma^{2}(x_{i}) f''(x_{i}) + \underbrace{\tfrac{1}{6} f'''(x_{i}) h^{3}(q_{i}^{+} - q_{i}^{-})}_{\text{third-order}} + \underbrace{\tfrac{h^{4}}{24}[q_{i}^{+} f^{(4)}(\xi_{+}) + q_{i}^{-} f^{(4)}(\xi_{-})]}_{\text{fourth-order}}.
\]

The first two terms are $(\calL f)(x_{i})$.

\textbf{Step 4 --- Bound the remainder.} From the moment-matching formulas, $|q_{i}^{\pm}| \leq C / h^{2}$. Moreover, $q_{i}^{+} - q_{i}^{-} = \mu(x_{i})/h$. Thus the third-order term is $\frac{1}{6} f'''(x_{i}) \mu(x_{i}) h^{2} = \mathcal{O}(h^{2})$. The fourth-order term is bounded by $\frac{h^{4}}{24} \cdot \frac{2C}{h^{2}} \cdot \|f^{(4)}\|_{\infty} = \mathcal{O}(h^{2})$.

Hence $|(Q^{(N)} f)(x_{i}) - (\calL f)(x_{i})| \leq C_{f}(x_{i}) \cdot h^{2}$.
\end{proof}

% ═══════════════════════════════════════════════════════════════
% PART IV — PROCESS CLASSES
% ═══════════════════════════════════════════════════════════════
\part{Process Classes}

\begin{quote}
\itshape This part extends the CTMC construction to the three canonical model classes encountered in financial applications: diffusions, jump-diffusions, and regime-switching processes. Each class is presented with a uniform structure---model definition, generator, CTMC construction, consistency analysis, and pricing interface---so that the reader can directly compare the approximation strategies and implement them in code.
\end{quote}

% ──────────────────────────────────────────────────────────────
\chapter{Jump-Diffusion and Regime-Switching Approximation}

\noindent\textbf{Purpose.} Extends the CTMC construction to jump-diffusions (general L\'{e}vy framework) and regime-switching models (block generator). Covers additive decomposition, compensated diffusion rates, and block matrix exponential.
\textbf{Prerequisites.} Chapters~9 (jump/RS generators),~13 (diffusion CTMC construction).
\textbf{Used in.} Chapters~18 (main theorem for J/R cases), Part~VII (parameter estimation with jumps/regimes).

\section{Generator Decomposition: The General L\'{e}vy Case}

Let $X$ be a jump-diffusion satisfying Assumption~J. Its generator admits the general decomposition into a diffusion part and a jump part governed by a L\'{e}vy measure $\nu$:

\[
\boxed{\mathcal{L} f(x) = \mu(x) f'(x) + \frac{1}{2} \sigma^{2}(x) f''(x) + \int_{\R \setminus \{0\}} \bigl[f(x + h(x,y)) - f(x)\bigr] \, \nu(dy)}.
\]

The integral term is the \textbf{jump generator} $\mathcal{L}_{J}$. It captures the contributions of discontinuities of arbitrary size and frequency. The L\'{e}vy measure $\nu$ on $\R \setminus \{0\}$ satisfies the standard integrability condition $\int_{\R \setminus \{0\}} (1 \wedge y^{2}) \, \nu(dy) < \infty$ (Assumption~J3).

This formulation encompasses the full range of jump specifications used in financial modeling:
\begin{itemize}
    \item \textbf{Compound Poisson:} $\nu(dy) = \lambda F_{Y}(dy)$, where $\lambda$ is the jump intensity and $F_{Y}$ is the jump-size distribution. This is the Merton (1976) framework.
    \item \textbf{State-dependent intensity:} $\nu(dy; x) = \lambda(x) F_{Y}(dy; x)$, generalizing to intensity and jump-size distribution that depend on the current state. The CEV-DEJD model falls into this category.
    \item \textbf{Infinite-activity L\'{e}vy processes:} $\nu(\R) = \infty$ but $\int (1 \wedge y^{2}) \nu(dy) < \infty$. Examples include the Variance Gamma (VG), Normal Inverse Gaussian (NIG), and CGMY processes. The CTMC framework can approximate these, though the infinite-activity case requires truncation of small jumps (see below).
    \item \textbf{Double-Exponential Jump-Diffusion (DEJD):} A special case of compound Poisson where $F_{Y}$ is the double-exponential distribution. This is treated in detail in Section~\ref{sec:dejd} as the running example.
\end{itemize}

\begin{remark}[Infinite-activity L\'{e}vy processes]
For processes with infinite jump activity ($\nu(\R) = \infty$), the jump generator integral $\int [f(x+y) - f(x)] \nu(dy)$ must be interpreted with care: the integral over small jumps may diverge if the first-order term $f'(x)y$ is not compensated. In such cases, one uses the compensated form
\[
\mathcal{L}_{J} f(x) = \int_{|y| \leq 1} [f(x+y) - f(x) - y f'(x)] \nu(dy) + \int_{|y| > 1} [f(x+y) - f(x)] \nu(dy).
\]
The CTMC approximation for infinite-activity processes requires either (i) truncation of jumps smaller than a threshold $\varepsilon$ (with $\varepsilon \to 0$ jointly with $h \to 0$), or (ii) a Gaussian approximation to the small-jump component justified by a Central Limit Theorem argument. This monograph focuses on finite-activity jump-diffusions (compound Poisson type), for which the simpler formulation without small-jump compensation suffices.
\end{remark}

\subsection{The Additive Decomposition Principle}

The additive decomposition $\mathcal{L} = \mathcal{L}_{D} + \mathcal{L}_{J}$ is the theoretical basis for the layered-construction strategy. The diffusion contribution $\mathcal{L}_{D}$ and the jump contribution $\mathcal{L}_{J}$ each handle their own domain without interfering. This additivity carries over to the CTMC construction, splitting it into the independent matrix sums $Q_{D}$ and $Q_{J}$.

\begin{proposition}[CTMC Generator for General Jump-Diffusions]
For a jump-diffusion with L\'{e}vy measure $\nu(dy; x)$ (state-dependent), the CTMC jump generator entries are
\[
q_{ij}^{J} := \int_{\R \setminus \{0\}} \1_{\{x_{i} + h(x_{i}, y) \in B_{j}\}} \, \nu(dy; x_{i}), \qquad i \neq j,
\]
with $B_{j}$ the Voronoi cell associated with grid point $x_{j}$. For the common additive jump specification $h(x, y) = y$, this reduces to
\[
q_{ij}^{J} = \int_{\R \setminus \{0\}} \1_{\{x_{i} + y \in B_{j}\}} \, \nu(dy; x_{i}) = \nu\bigl(\{y : x_{i} + y \in B_{j}\}; x_{i}\bigr).
\]
\end{proposition}

\medskip
\noindent\textbf{Why this matters.} The definition of the jump transition rate is intuitive: if the process jumps from $x_i$ by size $y$, and the landing point $x_i + y$ falls in the Voronoi cell of grid point $j$, that counts as a transition from $i$ to $j$. For finite-activity processes, the integral over the L\'{e}vy measure is a finite sum (over discretized jump sizes) or a one-dimensional integral over the jump-size density---both computationally manageable.

\begin{definition}[Jump Moment Contributions]
For each state $i$, compute:
\[
M_{1}^{J}(x_{i}) := \sum_{j=1}^{N} q_{ij}^{J} (x_{j} - x_{i}), \qquad
M_{2}^{J}(x_{i}) := \sum_{j=1}^{N} q_{ij}^{J} (x_{j} - x_{i})^{2}.
\]
\end{definition}

\medskip
\noindent\textbf{Why this matters.} $M_1^J$ and $M_2^J$ are the moment contributions from the jump component. They quantify how much average displacement and variance the jumps contribute. When constructing $Q_D$, these two quantities must be subtracted from the total targets $\mu$ and $\sigma^2$; otherwise, the diffusion and jump parts would double-count.

\begin{theorem}[Compensated Diffusion Rates --- General Form]
For a jump-diffusion where the target generator is $\mathcal{L} = \mathcal{L}_{D} + \mathcal{L}_{J}$, construct $Q_{J}^{(N)}$ to approximate $\mathcal{L}_{J}$, and compute $M_{1}^{J}, M_{2}^{J}$. Then construct the diffusion part $Q_{D}^{(N)}$ using the \emph{compensated} moments:
\[
\mu^{\text{net}}(x_{i}) := \mu(x_{i}) - M_{1}^{J}(x_{i}), \qquad
(\sigma^{2})^{\text{net}}(x_{i}) := \sigma^{2}(x_{i}) - M_{2}^{J}(x_{i}).
\]
On a uniform grid of spacing $h$,
\[
\boxed{\begin{aligned}
q_{i}^{+} &= \frac{(\sigma^{2})^{\text{net}}(x_{i}) + h \mu^{\text{net}}(x_{i})}{2h^{2}}, \\[8pt]
q_{i}^{-} &= \frac{(\sigma^{2})^{\text{net}}(x_{i}) - h \mu^{\text{net}}(x_{i})}{2h^{2}}.
\end{aligned}}
\]
The combined generator $Q^{(N)} = Q_{D}^{(N)} + Q_{J}^{(N)}$ then satisfies the full moment matching conditions.
\end{theorem}

\section{Running Example: Double-Exponential Jump-Diffusion (DEJD)}
\label{sec:dejd}

To make the general construction concrete, we specialize to the DEJD model---a compound Poisson jump-diffusion where the jump-size distribution is double-exponential. This is the most widely used parametric jump model in equity and commodity derivatives pricing.

\subsection{Model Specification}

The L\'{e}vy measure for the DEJD model is
\[
\nu(dy) = \lambda \, f_{Y}(y) \, dy,
\]
where $\lambda > 0$ is the jump intensity and $f_{Y}$ is the double-exponential density:
\[
f_{Y}(y) = p \eta_{1} e^{-\eta_{1} y} \1_{\{y \geq 0\}} + (1-p) \eta_{2} e^{\eta_{2} y} \1_{\{y < 0\}}.
\]
The parameter $p \in [0,1]$ controls the mixture of upward and downward jumps; $\eta_{1} > 0, \eta_{2} > 0$ are the exponential decay rates.

The jump generator then takes the concrete form:
\[
\mathcal{L}_{J} f(x) = \lambda \int_{-\infty}^{\infty} [f(x + h(x,y)) - f(x)] \, f_{Y}(y) \, dy.
\]

For multiplicative jumps (the standard in equity modeling), $h(x,y) = x(e^{y} - 1)$, and the generator becomes
\[
\mathcal{L}_{J} f(x) = \lambda \int_{-\infty}^{\infty} [f(x e^{y}) - f(x)] \, f_{Y}(y) \, dy.
\]

The detailed derivation of jump-transition rates, moment-generating functions, and the case-by-case integration formulas for this model are provided in Appendix~\ref{app:dejd-mgf} and Appendix~\ref{app:case-jump}. The reader may substitute any other jump-size distribution (e.g., Gaussian as in Merton, or a mixture) by replacing $f_{Y}$ and the corresponding primitive functions.

\subsection{Change of Variables for Multiplicative Jumps}

For the common multiplicative jump specification where the post-jump price is $x_{i} \cdot e^{Y}$, define $Z := e^{Y} - 1$ (the relative jump size). By the change-of-variables theorem,
\[
f_{Z}(z) = f_{Y}(\ln(z+1)) \cdot \frac{1}{z+1}, \qquad z > -1.
\]

\begin{theorem}[PDF of Relative Jump Size for Double-Exponential Jumps]
Suppose $Y$ follows a double-exponential distribution with density
\[
f_{Y}(y) = p \eta_{1} e^{-\eta_{1} y} \1_{\{y \geq 0\}} + (1-p) \eta_{2} e^{\eta_{2} y} \1_{\{y < 0\}}.
\]
Then the density of $Z = e^{Y} - 1$ is
\[
\boxed{f_{Z}(z) = p \eta_{1} (z+1)^{-1-\eta_{1}} \1_{\{z \geq 0\}} + (1-p) \eta_{2} (z+1)^{\eta_{2}-1} \1_{\{-1 < z < 0\}}}.
\]
\end{theorem}

\begin{proof}
For $z \geq 0$: $g^{-1}(z) = \ln(z+1) \geq 0$. The positive branch of $f_{Y}$ applies:
\[
f_{Z}(z) = p \eta_{1} e^{-\eta_{1} \ln(z+1)} \cdot \frac{1}{z+1} = p \eta_{1} (z+1)^{-1-\eta_{1}}.
\]
For $-1 < z < 0$: $g^{-1}(z) = \ln(z+1) < 0$. The negative branch applies:
\[
f_{Z}(z) = (1-p) \eta_{2} e^{\eta_{2} \ln(z+1)} \cdot \frac{1}{z+1} = (1-p) \eta_{2} (z+1)^{\eta_{2}-1}.
\]
For $z \leq -1$, $f_{Z}(z) = 0$.
\end{proof}

The jump transition rates are then computed via case-by-case integration, using the primitive functions $\Phi_{+}(z) = -p (u+1)^{-\eta_{1}}$ and $\Phi_{-}(z) = (1-p)(u+1)^{\eta_{2}}$, with the cases distinguished by the signs of the integration endpoints $\alpha_{i}(j-1)$ and $\alpha_{i}(j)$.

\section{Construction of $Q_{D}^{(N)}$ with Full Jump Compensation}

\begin{definition}[Jump Moment Contributions]
For each state $i$, compute:
\[
M_{1}^{J}(x_{i}) := \sum_{j=1}^{N} q_{ij}^{J} (x_{j} - x_{i}), \qquad
M_{2}^{J}(x_{i}) := \sum_{j=1}^{N} q_{ij}^{J} (x_{j} - x_{i})^{2}.
\]
\end{definition}

\begin{theorem}[Compensated Diffusion Rates]
On a uniform grid of spacing $h$,
\[
\boxed{\begin{aligned}
q_{i}^{+} &= \frac{\sigma_{\text{model}}^{2}(x_{i}) - M_{2}^{J}(x_{i}) + h(\mu(x_{i}) - M_{1}^{J}(x_{i}))}{2h^{2}}, \\[8pt]
q_{i}^{-} &= \frac{\sigma_{\text{model}}^{2}(x_{i}) - M_{2}^{J}(x_{i}) - h(\mu(x_{i}) - M_{1}^{J}(x_{i}))}{2h^{2}}.
\end{aligned}}
\]
\end{theorem}

\section{Regime-Switching: Block-Matrix Construction}

\subsection{Model Definition}

Under Assumption~R, the joint process $(X_{t}, R_{t})$ on $\R \times \{1,\ldots,M\}$ is a Feller process. The regime process $\{R_{t}\}$ is a finite-state CTMC with generator $Q^{R} = (\lambda_{mn})$ where $\lambda_{mn} \geq 0$ for $m \neq n$ and $\sum_{n} \lambda_{mn} = 0$. Conditionally on $R_{t} = m$, the spatial component follows
\[
dX_{t} = \mu_{m}(X_{t}) dt + \sigma_{m}(X_{t}) dW_{t},
\]
with $W_{t}$ independent of $\{R_{t}\}$.

\subsection{Generator}

The generator of the joint process on $C_{c}^{\infty}(\R \times \{1,\ldots,M\})$ is
\[
\boxed{\mathcal{L} f(x, m) = \mu_{m}(x) \frac{\partial f}{\partial x}(x,m) + \frac{1}{2} \sigma_{m}^{2}(x) \frac{\partial^{2} f}{\partial x^{2}}(x,m) + \sum_{n \neq m} \lambda_{mn} \bigl[f(x,n) - f(x,m)\bigr]}.
\]
This operator is the sum of three components: (i) the drift term within regime $m$, (ii) the diffusion term within regime $m$, and (iii) the regime-switching term, which captures jumps between regimes at rates $\lambda_{mn}$ with no change in the spatial coordinate.

\subsection{CTMC Construction: The Block Generator}

The CTMC approximation for a regime-switching process relies on the observation that the spatial diffusion within each regime can be discretized independently using the moment-matching formulas of Chapter~13. The regime transitions are exact (the regime process is already a finite-state CTMC). The combined generator is assembled as a block matrix.

\begin{theorem}[Block Generator for Regime-Switching]
For a regime-switching process with $M$ regimes and $N$ spatial grid points, the $K \times K$ generator ($K = N \cdot M$) has the block structure:
\[
\boxed{Q = \begin{pmatrix}
Q^{11} & Q^{12} & \cdots & Q^{1M} \\
Q^{21} & Q^{22} & \cdots & Q^{2M} \\
\vdots & \vdots & \ddots & \vdots \\
Q^{M1} & Q^{M2} & \cdots & Q^{MM}
\end{pmatrix}},
\]
where the blocks are $N \times N$ matrices defined by:
\[
Q^{mm} = Q_{D}^{(m)} - \Bigl(\sum_{n \neq m} \lambda_{mn}\Bigr) I_{N}, \qquad
Q^{mn} = \lambda_{mn} I_{N} \;\; (m \neq n).
\]
Here $Q_{D}^{(m)}$ is the tridiagonal diffusion generator for regime $m$, constructed using $\mu_{m}$ and $\sigma_{m}^{2}$ via the formulas of Chapter~13. The term $-(\sum_{n \neq m} \lambda_{mn}) I_{N}$ in the diagonal blocks accounts for the probability mass that leaves the current regime.
\end{theorem}

\medskip
\noindent\textbf{Why this matters.} The block-matrix structure decouples the construction into three steps: (1) build $M$ independent tridiagonal matrices $Q_{D}^{(m)}$ using the per-regime parameters; (2) subtract the regime-switching rates from the diagonal blocks; (3) place the regime-switching rates on the off-diagonal blocks. Each off-diagonal block $Q^{mn}$ is a scalar multiple of the identity, reflecting the fact that regime transitions occur at a common rate for all spatial states and leave the spatial coordinate unchanged.

\begin{theorem}[Generator Consistency for Regime-Switching]
For any $f \in C_{c}^{\infty}(\R \times \{1,\ldots,M\})$, restricted to the grid as $\mathbf{f}_{m} = (f(x_{1},m), \ldots, f(x_{N},m))^{\top}$, the block CTMC generator satisfies
\[
(Q \mathbf{f})_{i,m} = \mu_{m}(x_{i}) f_{x}(x_{i},m) + \frac{1}{2} \sigma_{m}^{2}(x_{i}) f_{xx}(x_{i},m) + \sum_{n \neq m} \lambda_{mn} [f(x_{i},n) - f(x_{i},m)] + \mathcal{O}(h^{2}),
\]
uniformly over the grid. The spatial error $\mathcal{O}(h^{2})$ is inherited from the single-regime construction; the regime-switching terms are exact.
\end{theorem}

\begin{proof}
For regime $m$ at grid point $i$, the action of the block generator decomposes as:
\[
\begin{aligned}
(Q \mathbf{f})_{i,m} &= \underbrace{(Q_{D}^{(m)} \mathbf{f}_{m})_{i}}_{\text{spatial diffusion}} \;-\; \underbrace{\Bigl(\sum_{n \neq m} \lambda_{mn}\Bigr) f(x_{i},m)}_{\text{exit from regime } m} \;+\; \underbrace{\sum_{n \neq m} \lambda_{mn} f(x_{i},n)}_{\text{entry into regime } m} \\
&= (Q_{D}^{(m)} \mathbf{f}_{m})_{i} + \sum_{n \neq m} \lambda_{mn} [f(x_{i},n) - f(x_{i},m)].
\end{aligned}
\]
The spatial term $(Q_{D}^{(m)} \mathbf{f}_{m})_{i} = \mu_{m}(x_{i}) f_{x}(x_{i},m) + \frac{1}{2} \sigma_{m}^{2}(x_{i}) f_{xx}(x_{i},m) + \mathcal{O}(h^{2})$ by Theorem~13.2.5. The regime-switching terms are exact (no approximation error). Summing yields the stated bound.
\end{proof}

\subsection{Block Transition Matrix}

For the Bermudan dynamic programming recursion, we require the block transition matrix $P_{\Delta t} = e^{\Delta t Q}$. When the block matrix has the special structure $Q^{mm} = A_{m} - \alpha_{m} I$ and $Q^{mn} = \lambda_{mn} I$ ($m \neq n$), the matrix exponential can be decomposed:

\begin{proposition}[Block Exponential for Regime-Switching]
Let $Q$ be the $K \times K$ block generator defined above. Write it as
\[
Q = \underbrace{\begin{pmatrix}
Q_{D}^{(1)} & 0 & \cdots & 0 \\
0 & Q_{D}^{(2)} & \cdots & 0 \\
\vdots & \vdots & \ddots & \vdots \\
0 & 0 & \cdots & Q_{D}^{(M)}
\end{pmatrix}}_{\text{spatial}} \;+\; \underbrace{(Q^{R} \otimes I_{N})}_{\text{regime}},
\]
where $\otimes$ denotes the Kronecker product. Since the spatial and regime generators commute (they act on different coordinates), we have
\[
e^{\Delta t Q} = e^{\Delta t (Q^{R} \otimes I_{N})} \cdot \operatorname{diag}\bigl(e^{\Delta t Q_{D}^{(1)}}, \ldots, e^{\Delta t Q_{D}^{(M)}}\bigr).
\]
In particular, for $M=2$, $e^{\Delta t (Q^{R} \otimes I_{N})} = P^{R}(\Delta t) \otimes I_{N}$, where $P^{R}(\Delta t)$ is the $2 \times 2$ regime transition matrix.
\end{proposition}

\subsection{Two-State Closed-Form Transition Probabilities}

For the important special case $M=2$, the regime transition matrix has an analytical closed form:

\begin{theorem}[Two-State Regime Transition Probabilities]
For $M = 2$ with generator $Q^{R} = \bigl(\begin{smallmatrix} -\lambda_{12} & \lambda_{12} \\ \lambda_{21} & -\lambda_{21} \end{smallmatrix}\bigr)$,
\[
\boxed{P^{R}(\Delta t) = \frac{1}{\lambda_{12} + \lambda_{21}} \begin{pmatrix}
\lambda_{21} + \lambda_{12} e^{-(\lambda_{12}+\lambda_{21})\Delta t} & \lambda_{12}(1 - e^{-(\lambda_{12}+\lambda_{21})\Delta t}) \\[4pt]
\lambda_{21}(1 - e^{-(\lambda_{12}+\lambda_{21})\Delta t}) & \lambda_{12} + \lambda_{21} e^{-(\lambda_{12}+\lambda_{21})\Delta t}
\end{pmatrix}}.
\]
\end{theorem}

\medskip
\noindent\textbf{Why this matters.} This closed-form solution eliminates the need to compute the matrix exponential for the regime component. $\lambda_{12} + \lambda_{21}$ is the total rate of regime switching; the larger it is, the faster regime probabilities converge to the stationary distribution $\pi = (\frac{\lambda_{21}}{\lambda_{12}+\lambda_{21}}, \frac{\lambda_{12}}{\lambda_{12}+\lambda_{21}})$. In parameter estimation (Hamilton filter), this closed form saves one $2 \times 2$ exponential computation per observation.

\subsection{Pricing Interface}

For the regime-switching CTMC, the Bermudan dynamic programming recursion operates on the stacked value vector $\mathbf{V} \in \R^{K}$, where the first $N$ entries are values in regime~1 and the next $N$ entries are values in regime~2. The recursion is:
\[
\mathbf{V}^{(\ell)} = \max\bigl(\boldsymbol{\phi}_{\text{stack}}, \; e^{-r \Delta t} \,\mathbf{P}_{\Delta t}^{(N)} \,\mathbf{V}^{(\ell+1)}\bigr),
\]
where $\mathbf{P}_{\Delta t}^{(N)} = e^{\Delta t Q}$ is the $K \times K$ block transition matrix and $\boldsymbol{\phi}_{\text{stack}}$ is the payoff vector repeated for each regime (since the payoff depends only on the spatial coordinate, not the regime). The computational cost per step is $\mathcal{O}(K^{2}) = \mathcal{O}(M^{2} N^{2})$ for dense matrix-vector multiplication, reduced to $\mathcal{O}(M^{2} N)$ when spatial blocks are tridiagonal.

\begin{proof}
Solve the Kolmogorov ODE $\frac{d}{dt}p_{11}(t) = -\lambda_{12}p_{11}(t) + \lambda_{21}(1 - p_{11}(t))$ with initial condition $p_{11}(0) = 1$. The solution is $p_{11}(t) = (\lambda_{21} + \lambda_{12} e^{-(\lambda_{12}+\lambda_{21})t}) / (\lambda_{12}+\lambda_{21})$. The remaining entries follow from probability conservation and symmetry.
\end{proof}

% ═══════════════════════════════════════════════════════════════
% PART V — OPTIMAL STOPPING
% ═══════════════════════════════════════════════════════════════
\part{Optimal Stopping}

% ──────────────────────────────────────────────────────────────
\chapter{Optimal Stopping Theory}

\noindent\textbf{Purpose.} This chapter formulates the American option pricing problem as an optimal stopping problem. It introduces the Snell envelope, the continuation and exercise regions, and the free-boundary PDE characterization. This is the mathematical bridge from process theory (Parts~I--III) to the numerical pricing scheme (Chapter~16).

\textbf{Prerequisites.} Part~II (Markov processes), Part~III (generator theory, Feynman--Kac).

\textbf{Used in.} Chapter~16 (Bermudan approximation and dynamic programming).

\section{The American Option as an Optimal Stopping Problem}

Let $\{X_{t}\}$ be a Markov process under the risk-neutral measure $\Q$, $r \geq 0$ the risk-free rate, and $\phi: \R \to \R$ the payoff function. The American option value at time $t$ is:
\[
V(t,x) = \sup_{\tau \in \calT_{[t,T]}} \E^{\Q}\left[e^{-r(\tau-t)} \phi(X_{\tau}) \mid X_{t} = x\right],
\]
where $\calT_{[t,T]}$ is the set of all $\{\F_{s}\}$-stopping times taking values in $[t,T]$.

\section{The Snell Envelope}

\begin{definition}[Snell Envelope]
Let $Z_{t} := e^{-rt} \phi(X_{t})$. The Snell envelope of $\{Z_{t}\}$ is the smallest supermartingale that dominates $\{Z_{t}\}$:
\[
U_{t} := \operatorname*{ess\,sup}_{\tau \in \calT_{[t,T]}} \E[Z_{\tau} \mid \F_{t}].
\]
\end{definition}

\medskip
\noindent\textbf{Why this matters.} The Snell envelope is the central tool for American option pricing. Think of the discounted option payoff $Z_t$ as a curve fluctuating over time; the Snell envelope $U_t$ is a supermartingale that stays at or above this curve at all times while hugging it as closely as possible. It must be a supermartingale because in a fair game (martingale) you cannot profit by timing---the supermartingale property guarantees that discounted prices under the optimal exercise strategy admit no arbitrage.

\begin{theorem}[Properties of the Snell Envelope]
\begin{enumerate}[label=(\arabic*)]
    \item $U$ is a supermartingale.
    \item $U_{t} \geq Z_{t}$ for all $t$.
    \item $U_{T} = Z_{T}$.
    \item An optimal stopping time is $\tau^{*} := \inf\{t \in [0,T] : U_{t} = Z_{t}\}$.
\end{enumerate}
\end{theorem}

\medskip
\noindent\textbf{Why this matters.} The fourth property gives an operational definition of the optimal stopping time: exercise the option at the first moment the Snell envelope touches the payoff curve ($U_t = Z_t$). If $U_t > Z_t$, holding on is worth more than exercising immediately---you should wait. This is like bidding at an auction: as long as the current bid is below your reserve, you stay silent; the instant it hits your reserve, you act.

\begin{theorem}[Continuation and Exercise Regions]
There exists an optimal exercise boundary $b: [0,T] \to \R$ such that:
\begin{itemize}
    \item \textbf{Continuation region} $\calC := \{(t,x) : V(t,x) > \phi(x)\}$,
    \item \textbf{Exercise region} $\calE := \{(t,x) : V(t,x) = \phi(x)\}$.
\end{itemize}
In $\calC$, the value function satisfies $\frac{\partial V}{\partial t} + \calL V - rV = 0$, where $\calL$ is the generator of $X$.
\end{theorem}

\medskip
\noindent\textbf{Why this matters.} The dividing line $b(t)$ between the continuation and exercise regions is the critical boundary of American option pricing. In the continuation region, the option is not yet worth exercising and its price satisfies the standard PDE ($\frac{\partial V}{\partial t} + \mathcal{L}V - rV = 0$); once $b(t)$ is crossed into the exercise region, $V$ collapses to the payoff $\phi$. Solving this free-boundary problem is the central numerical challenge---the location of the boundary is unknown in advance and must be solved jointly with the price function.

\section{The Free-Boundary Problem}

\begin{theorem}[Continuation and Exercise Regions]
There exists an optimal exercise boundary $b: [0,T] \to \R$ such that:
\begin{itemize}
    \item \textbf{Continuation region} $\calC := \{(t,x) : V(t,x) > \phi(x)\}$,
    \item \textbf{Exercise region} $\calE := \{(t,x) : V(t,x) = \phi(x)\}$.
\end{itemize}
In $\calC$, the value function satisfies $\frac{\partial V}{\partial t} + \calL V - rV = 0$, where $\calL$ is the generator of $X$.
\end{theorem}

% ──────────────────────────────────────────────────────────────
\chapter{Bermudan Approximation and Dynamic Programming}

\noindent\textbf{Purpose.} This chapter bridges optimal stopping theory (Chapter~15) to the CTMC pricing algorithm. It shows how to discretize the continuous exercise opportunity set into finitely many Bermudan exercise dates, derives the backward dynamic programming recursion, and applies it to the CTMC transition matrix $P_{\Delta t}^{(N)}$. The $\max(\phi, e^{-r\Delta t} P_{\Delta t} V)$ recursion is the core numerical routine for American option pricing.

\textbf{Prerequisites.} Chapters~13 (CTMC generator), 15 (optimal stopping theory).

\textbf{Used in.} Chapters~17--18 (pricing pipeline and main theorem). This chapter provides the temporal discretization component of the error decomposition.

\section{Discretization of the Exercise Opportunity Set}

\begin{definition}[Bermudan Option]
For a partition $\Pi_{L} := \{0 = t_{0} < t_{1} < \cdots < t_{L} = T\}$ with $\Delta t = T/L$, the Bermudan value is
\[
V_{B}^{(L)}(t_{\ell}, x) := \sup_{\tau \in \calT_{\Pi_{L}}^{t_{\ell}}} \E\left[e^{-r(\tau - t_{\ell})} \phi(X_{\tau}) \mid X_{t_{\ell}} = x\right],
\]
where $\calT_{\Pi_{L}}^{t_{\ell}}$ is the set of stopping times taking values in $\{t_{\ell}, \ldots, t_{L}\}$.
\end{definition}

\medskip

\noindent\textbf{Why this matters.} A Bermudan option is the restricted cousin of the American option---you may exercise only on a pre-specified list of fixed dates (e.g., each quarter-end), not at any moment as with an American option. This discretizes the continuous anytime-exercise right into finitely many exercise opportunities, reducing the problem from continuous optimal stopping to discrete dynamic programming, which can be solved exactly by backward induction.

\begin{definition}[Snell Envelope Operator]
For a payoff $g\in C_b(E)$ and discount rate $r\ge0$, define the one-step stopping operator associated with a semigroup $P_t$ by
\[
\mathcal{S}_{\Delta t}u(x) := \max\left\{g(x), e^{-r\Delta t}P_{\Delta t}u(x)\right\}.
\]
For the CTMC approximation, replace $P_{\Delta t}$ by $P_{\Delta t}^{(n)}=e^{\Delta t Q^{(n)}}$ and define
\[
\mathcal{S}_{\Delta t}^{(n)}u(x) := \max\left\{g(x), e^{-r\Delta t}P_{\Delta t}^{(n)}u(x)\right\}, \qquad x\in E_n.
\]
The Bermudan value is the backward iteration of this nonlinear operator, not merely an ad hoc dynamic-programming formula.
\end{definition}

\begin{lemma}[Non-Expansiveness of the Stopping Operator]
For $u,v\in C_b(E)$,
\[
\|\mathcal{S}_{\Delta t}u-\mathcal{S}_{\Delta t}v\|_\infty
\le e^{-r\Delta t}\|u-v\|_\infty.
\]
The same bound holds for $\mathcal{S}_{\Delta t}^{(n)}$ on $E_n$.
\end{lemma}

\begin{proof}
The map $a\mapsto \max\{g(x),a\}$ is one-Lipschitz for each $x$, and $P_{\Delta t}$ and $P_{\Delta t}^{(n)}$ are contraction semigroups in the supremum norm.
\end{proof}

\begin{theorem}[Convergence of Bermudan to American Value]
For a Feller process $X$ and Lipschitz $\phi$,
\[
\lim_{L \to \infty} V_{B}^{(L)}(0, X_{0}) = V(0, X_{0}), \quad \text{a.s.\ and in } L^{1}.
\]
\end{theorem}

\medskip
\noindent\textbf{Why this matters.} This limit theorem guarantees that by adding more exercise dates ($L \to \infty$, i.e., $\Delta t \to 0$), the Bermudan price converges to the true American price. Intuitively, if you can exercise almost every day, the difference from continuous exercise becomes negligible. This legitimizes the CTMC dynamic-programming method---we only need to compute on discrete time points, choosing $L$ large enough.

\begin{theorem}[Bermudan DP Recursion]
\[
\begin{aligned}
V(t_{L}, x) &= \phi(x), \\
V(t_{\ell}, x) &= \max\left(\phi(x), \; e^{-r \Delta t} \E\left[V(t_{\ell+1}, X_{t_{\ell+1}}) \mid X_{t_{\ell}} = x\right]\right), \quad \ell = L-1, \ldots, 0.
\end{aligned}
\]
In operator notation: $\mathbf{V}_{\ell} = \max(\boldsymbol{\phi}, \; e^{-r \Delta t} P_{\Delta t} \mathbf{V}_{\ell+1})$.
\end{theorem}

\medskip
\noindent\textbf{Why this matters.} The dynamic-programming recursion is the computational heart of American option pricing. Working backward from terminal time $T$: at each time step, decide whether immediate exercise or continued holding is worth more, and take the larger. The continuation value is the discounted expectation of the next period's value ($\mathbb{E}[V_{\ell+1} \mid X_{\ell}]$). This $\max$ operation is the discrete version of the optimal stopping decision---at every step, it is a binary choice: exercise or wait.

\begin{theorem}[CTMC Pricing Recursion]
Let $\mathbf{P}_{\Delta t}^{(N)} = e^{\Delta t \cdot Q^{(N)}}$ and $\boldsymbol{\phi} = (\phi(x_{1}), \ldots, \phi(x_{N}))^{\top}$. Then:
\begin{enumerate}[label=(\arabic*)]
    \item \textbf{Terminal Condition:} $\mathbf{V}^{(L)} = \boldsymbol{\phi}$.
    \item \textbf{Backward Recursion} for $\ell = L-1, \ldots, 0$:
    \[
    \mathbf{V}^{(\ell)} = \max\left(\boldsymbol{\phi}, \; e^{-r \Delta t} \mathbf{P}_{\Delta t}^{(N)} \mathbf{V}^{(\ell+1)}\right).
    \]
    \item \textbf{Price Interpolation:} $V_{0}^{(N)} = \sum_{j=1}^{N} [\mathbf{P}_{\Delta t}^{(N)}]_{i_{0}, j} \cdot V_{j}^{(1)}$, where $i_{0}$ is the grid state corresponding to the initial price $x_{0}$.
\end{enumerate}
\end{theorem}

\medskip
\noindent\textbf{Why this matters.} The CTMC pricing recursion is the last mile from theory to implementation. It reduces the entire pricing problem to matrix-vector multiplication and elementwise $\max$ operations---both efficiently handled by any linear-algebra library. Each row of $\mathbf{P}_{\Delta t}^{(N)}$ is a probability distribution; $e^{-r\Delta t}\mathbf{P}_{\Delta t} \mathbf{V}^{(\ell+1)}$ is the risk-neutral-probability-weighted average of next-period values, discounted. The whole process runs $L$ steps at $\mathcal{O}(N^2)$ per step (dense) or $\mathcal{O}(N)$ (tridiagonal)---a numerical scheme of moderate algorithmic complexity.

\begin{proposition}[Error Accumulation]
\[
|V_{0}^{(N)} - V(0, x_{0})| \leq \underbrace{\mathcal{O}(1/\sqrt{L})}_{\text{Bermudan}} + \underbrace{\mathcal{O}(T h^{2})}_{\text{CTMC discretization}} + \underbrace{\mathcal{O}(h)}_{\text{exercise boundary}}.
\]
\end{proposition}

\medskip
\noindent\textbf{Why this matters.} This error decomposition identifies three independent error sources and their respective convergence rates. The Bermudan error decays at $1/\sqrt{L}$ (the slowest), the CTMC discretization error decays at $h^2$ (faster), and the exercise-boundary discretization error decays at $h$. The boundary error is often the dominant term because it is only first-order---this suggests grid refinement near the strike price. Richardson extrapolation can further accelerate Bermudan convergence.

% ═══════════════════════════════════════════════════════════════
% PART VI — UNIFIED FRAMEWORK
% ═══════════════════════════════════════════════════════════════
\part{Pricing Framework}

\chapter{The CTMC Pricing Algorithm}

\noindent\textbf{Purpose.} Presents the unified pricing pipeline: Model Specification $\to$ Generator $\mathcal{L}$ $\to$ CTMC Generator $Q^{(N)}$ $\to$ Transition Matrix $P_{\Delta t}$ $\to$ Dynamic Programming $\to$ Option Price.
\textbf{Prerequisites.} Chapters~13--14 (CTMC construction),~16 (Bermudan DP).
\textbf{Used in.} Chapter~18 (main theorem brings all pieces together).

The entire methodology is captured by a single pipeline:
\[
\boxed{
\begin{array}{c}
\textbf{Model Specification} \;\longrightarrow\; \textbf{Generator } \calL \;\longrightarrow\; \textbf{CTMC Generator } Q^{(N)} \\
\downarrow \\
\textbf{Transition Matrix } \mathbf{P}_{\Delta t} = e^{\Delta t \cdot Q^{(N)}} \;\longrightarrow\; \textbf{Dynamic Programming} \\
\downarrow \\
\textbf{American Option Price}
\end{array}
}
\]

\chapter{Model Class Unification}

\noindent\textbf{Purpose.} States the main convergence theorems (Diffusion, Jump-Diffusion, Regime-Switching cases), the invariance proposition (separation of model and pricing), and the unified error decomposition---the capstone results of the monograph.
\textbf{Prerequisites.} Chapters~13--14 (CTMC construction),~16 (Bermudan DP),~17 (pricing pipeline).
\textbf{Used in.} This is the terminal chapter of the mainline; it provides the theoretical guarantee for the entire CTMC pricing methodology.

\section{The Separation Principle}

\begin{theorem}[Invariance Proposition --- Separation of Model and Pricing]
The mapping $\text{Model Parameters} \mapsto Q^{(N)} \mapsto \mathbf{P}_{\Delta t} \mapsto \mathbf{V}^{(0)}$ is a composition of two decoupled stages:
\begin{enumerate}[label=(\arabic*)]
    \item \textbf{Model-specific stage:} $\text{Parameters} \mapsto Q^{(N)}$. This is the only stage that depends on the model class (diffusion, jump-diffusion, or regime-switching). The construction follows the same structural principle---moment matching of generator coefficients---regardless of the model.
    \item \textbf{Model-agnostic stage:} $Q^{(N)} \mapsto \mathbf{P}_{\Delta t} \mapsto \mathbf{V}^{(0)}$. This stage consists solely of the matrix exponential $e^{\Delta t Q^{(N)}}$ and the Bermudan dynamic programming recursion. It depends only on the transition matrix and the payoff vector, not on how the transition matrix was obtained.
\end{enumerate}
\end{theorem}

\medskip
\noindent\textbf{Why this matters.} The separation principle reveals the decoupled-design philosophy of the framework: the pricing engine is completely oblivious to which model was used upstream. Any model---GBM, CEV, Merton, regime-switching---that can produce a valid generator matrix $Q^{(N)}$ plugs directly into the same pricing pipeline without modification.

\section{Main Theorem C: Pricing Convergence}

This section assembles the three main theorems of the monograph. Theorem~A (Chapters~12--13) provides generator approximation. Theorem~B (Chapters~9,~12) provides weak convergence. Theorem~C (stated below) provides the final pricing guarantee.

\begin{theorem}[American Option Pricing Convergence]
\label{thm:main}
Let $V^{(n)}(t,x)$ be the Bermudan option value computed under the CTMC with exercise grid size $\Delta t_n$ and state grid $E_n$.

Let $V(t,x)$ be the value of the corresponding American option:
\[
V(t,x) = \sup_{\tau \in \mathcal{T}} \mathbb{E}_{t,x}
\left[e^{-r(\tau-t)} g(X_\tau)\right].
\]

Assume:
\begin{enumerate}[label=(B\arabic*)]
    \item \textbf{CTMC convergence.} $X^{(n)} \Rightarrow X$.
    \item \textbf{Payoff regularity.} The payoff function satisfies $g \in C_b(E)$.
    \item \textbf{Uniform integrability.} The family of discounted payoffs is uniformly integrable.
    \item \textbf{Snell-envelope stability.} The Bermudan values are generated by the stopping operators $\mathcal{S}_{\Delta t_n}^{(n)}$ and $\mathcal{S}_{\Delta t_n}$, whose non-expansiveness controls error propagation.
\end{enumerate}
Then
\[
\lim_{n \to \infty} V^{(n)}(t,x) = V(t,x) \quad \forall (t,x).
\]
Moreover, the Bermudan approximation satisfies
\[
|V^{(n)} - V| \le C (\Delta t_n + \varepsilon_n).
\]
\end{theorem}

\begin{proof}[Proof Architecture --- How Theorems A, B, C Connect]
\textbf{Theorem~A (Generator Approximation, Chapters~12--13):} Local operator consistency on a dense core, together with conservativity of the CTMC generators, gives semigroup convergence $P_t^{(n)} f \to P_t f$ for every $f \in C_0(E)$.

\textbf{Theorem~B (Weak Convergence, Trotter--Kato/Kurtz):} Semigroup convergence plus convergence of initial laws yields process-level convergence $X^{(n)} \Rightarrow X$ in the Skorokhod space $D([0,T],E)$.

\textbf{Theorem~C (Pricing Convergence, below):} Optimal-stopping stability, weak convergence, payoff regularity, and uniform integrability yield convergence of CTMC Bermudan prices to the American value.
\end{proof}

\begin{theorem}[Main Convergence Theorem --- Jump-Diffusion Case]
Under Assumption~J and the construction of Chapter~14 with combined generator $Q^{(N)} = Q_{D}^{(N)} + Q_{J}^{(N)}$, the same convergence statement holds with the generator error rate modified to
\[
\|Q^{(N)} f - \mathcal{L} f\|_{\infty} \leq C_{f} (h^{2} + \varepsilon(\nu)),
\]
where $\varepsilon(\nu)$ is the discretization error of the L\'{e}vy measure (identically zero for finite-activity compound Poisson processes, and $\mathcal{O}(h)$ for infinite-activity processes with small-jump truncation at threshold $\asymp h$).
\end{theorem}

\begin{theorem}[Main Convergence Theorem --- Regime-Switching Case]
Under Assumption~R and the block-matrix construction of Chapter~14, the same convergence statement holds. The generator error inherits the diffusion rate $\mathcal{O}(h^{2})$ within each regime; the regime transition probabilities are exact (the finite-state regime CTMC is not an approximation).
\end{theorem}

\section{Unified Error Decomposition}

For all three model classes, the pricing error admits the additive decomposition:

\begin{center}
\begin{tabular}{>{\raggedright\arraybackslash}p{3cm}>{\raggedright\arraybackslash}p{3cm}>{\raggedright\arraybackslash}p{3cm}>{\raggedright\arraybackslash}p{3cm}}
\toprule
\textbf{Error Source} & \textbf{Origin} & \textbf{Rate} & \textbf{Chapter} \\
\midrule
Generator (spatial) & $Q^{(N)} f - \mathcal{L} f$ & $\mathcal{O}(h^{2})$ & 12--13 \\
Bermudan (temporal) & Discrete exercise dates & $\mathcal{O}(L^{-1/2})$ & 16 \\
Boundary & Grid truncation & $\mathcal{O}(h)$ & \ref{ch:boundary} \\
Truncation & Probability outside $[a,b]$ & $\mathcal{O}(e^{-\gamma/h})$ & \ref{ch:boundary} \\
Jump discretization (if applicable) & L\'{e}vy measure approximation & $\mathcal{O}(h)$ or $0$ & 14 \\
\bottomrule
\end{tabular}
\end{center}

\medskip
\noindent\textbf{Why this matters.} The error decomposition identifies five independent sources of error, each with its own convergence rate. The dominant terms in practice are (i) the Bermudan error $\mathcal{O}(L^{-1/2})$, which decays slowly and typically requires $L \sim 1000$--$5000$ for basis-point accuracy, and (ii) the boundary error $\mathcal{O}(h)$, which motivates local grid refinement near important thresholds (strike prices, barriers). The generator error $\mathcal{O}(h^{2})$ is usually smaller than the boundary error for moderate $h$.

% ═══════════════════════════════════════════════════════════════


\chapter{Error Analysis for the CTMC Pricing Framework}

\noindent\textbf{Purpose.} Provides a systematic five-layer error decomposition ($\varepsilon_{\text{gen}}$, $\varepsilon_{\text{semi}}$, $\varepsilon_{\text{berm}}$, $\varepsilon_{\text{bound}}$, $\varepsilon_{\text{price}}$) and the error propagation chain. Supplies quantitative rates and practical guidance for grid resolution and time-stepping.
\textbf{Prerequisites.} Chapters~13--14 (consistency),~16 (Bermudan convergence),~18 (main theorem).
\textbf{Used in.} Reference for numerical implementation; complements Chapter~18 with detailed estimates.

This appendix provides a systematic decomposition of the error between the CTMC-Bermudan option price $V_{0}^{(N,L)}$ and the true American option value $V(0,X_{0})$. The analysis is organized in two dimensions: (i) a \emph{layer} decomposition identifying the independent sources of error, and (ii) a \emph{propagation} analysis tracking how error at one stage of the pipeline amplifies or attenuates in subsequent stages.

\section{Error Decomposition by Source (Layered Analysis)}

The total pricing error can be partitioned into five independent sources, each with a distinct origin in the approximation pipeline and its own convergence rate:

\begin{definition}[Five-Layer Error Decomposition]
\[
\boxed{\varepsilon_{\text{total}} := |V_{0}^{(N,L)} - V(0,X_{0})| \leq \varepsilon_{\text{gen}} + \varepsilon_{\text{semi}} + \varepsilon_{\text{berm}} + \varepsilon_{\text{bound}} + \varepsilon_{\text{price}}},
\]
where each term is defined below.
\end{definition}

\subsection{$\varepsilon_{\text{gen}}$: Generator (Spatial) Error}

\textbf{Source.} The discrete generator $Q^{(N)}$ approximates the continuous generator $\mathcal{L}$ on a core of test functions:
\[
\varepsilon_{\text{gen}} \sim \sup_{f \in \mathcal{D}_{0}} \|Q^{(N)} f - \mathcal{L} f\|_{\infty}.
\]

\textbf{Rate.} For the tridiagonal moment-matching construction on a uniform grid of spacing $h$:
\[
\varepsilon_{\text{gen}} \leq C_{1} T h^{2},
\]
where $C_{1}$ depends on $\max(|\mu|, |\mu'|, |\mu''|, |\sigma^{2}|, |(\sigma^{2})'|, |(\sigma^{2})''|)$ and on the fourth derivatives of the test functions in the core.

\textbf{Origin.} The error $\mathcal{O}(h^{2})$ arises from the third-derivative term in the Taylor expansion that is not canceled by the first two moment matching conditions. The constant $C_{1}$ is derived explicitly in Chapter~13 (Theorem~13.3.1).

\subsection{$\varepsilon_{\text{semi}}$: Semigroup (Transition) Error}

\textbf{Source.} Generator convergence at rate $\mathcal{O}(h^{2})$ implies semigroup convergence via the Trotter--Kato estimate:
\[
\varepsilon_{\text{semi}} \sim \sup_{t \in [0,T]} \|e^{t Q^{(N)}} - P_{t}\|.
\]

\textbf{Rate.} Under uniform boundedness of the semigroups,
\[
\varepsilon_{\text{semi}} \leq T \cdot \varepsilon_{\text{gen}} \leq C_{2} T^{2} h^{2}.
\]
The error grows linearly with the time horizon $T$ because the discrepancy between $e^{t Q^{(N)}}$ and $P_{t}$ accumulates over time.

\textbf{Origin.} The semigroup error is derived from the variation-of-parameters formula:
\[
e^{t Q^{(N)}} f - P_{t} f = \int_{0}^{t} e^{(t-s) Q^{(N)}} (Q^{(N)} - \mathcal{L}) P_{s} f \, ds.
\]
Taking norms and using $\|e^{(t-s) Q^{(N)}}\| \leq 1$ and $\|P_{s}\| \leq 1$ yields the linear-in-$T$ bound.

\subsection{$\varepsilon_{\text{berm}}$: Bermudan (Temporal Discretization) Error}

\textbf{Source.} Restricting exercise to $L$ equally spaced dates $\{t_{0}, \ldots, t_{L}\}$ instead of the continuous interval $[0,T]$:
\[
\varepsilon_{\text{berm}} = V(0,X_{0}) - V_{B}^{(L)}(0,X_{0}) \geq 0.
\]

\textbf{Rate.} For a diffusion with Lipschitz payoff and bounded coefficients:
\[
\varepsilon_{\text{berm}} \leq C_{3} \frac{\sigma_{\max} \sqrt{T}}{\sqrt{L}},
\]
where $\sigma_{\max} = \sup_{x} |\sigma(x)|$ and $C_{3}$ depends on the Lipschitz constant of the payoff and the drift bound.

\textbf{Origin.} The ${1}/{\sqrt{L}}$ rate reflects the Brownian scaling: between exercise dates, the process diffuses by $\mathcal{O}(\sqrt{\Delta t}) = \mathcal{O}(\sqrt{T/L})$, limiting the potential gain from more frequent exercise. This rate is sharp (cannot be improved to $1/L$ without additional smoothness).

\subsection{$\varepsilon_{\text{bound}}$: Boundary (Truncation) Error}

\textbf{Source.} Truncating the unbounded (or semi-bounded) state space to a finite interval $[a,b]$ and applying reflecting or absorbing boundary conditions:
\[
\varepsilon_{\text{bound}} \leq \varepsilon_{\text{boundary layer}} + \varepsilon_{\text{truncation}}.
\]

\textbf{Rate.}
\begin{itemize}
    \item Boundary layer error (from the incorrect boundary condition at $a$ and $b$): $\mathcal{O}(h)$.
    \item Truncation error (from probability mass outside $[a,b]$): $\mathcal{O}(e^{-\gamma / h})$ for diffusions with a stationary distribution, where $\gamma \propto \min(|a - X_{0}|, |b - X_{0}|)$.
\end{itemize}

\textbf{Origin.} The $\mathcal{O}(h)$ boundary layer error dominates for moderate $h$. It arises because the moment-matching formulas at interior grid points are $\mathcal{O}(h^{2})$ accurate, but at the two boundary points only one-sided rates are available, reducing the local consistency order to $\mathcal{O}(h)$.

\subsection{$\varepsilon_{\text{price}}$: Pricing (Value Function) Error}

\textbf{Source.} The cumulative effect of all upstream errors on the value function, propagated through the backward dynamic programming recursion:
\[
\varepsilon_{\text{price}} \sim \|\mathbf{V}^{(0)} - \mathbf{V}_{\text{true}}^{(0)}\|_{\infty}.
\]

\textbf{Rate.} For the Bermudan-CTMC recursion, the value function error is bounded by the transition error propagated through $L$ steps:
\[
\varepsilon_{\text{price}} \leq \varepsilon_{\text{semi}} \cdot \sum_{\ell=0}^{L-1} e^{-r \ell \Delta t} \leq \varepsilon_{\text{semi}} \cdot \frac{1 - e^{-r T}}{1 - e^{-r \Delta t}} \approx \varepsilon_{\text{semi}} \cdot \frac{L}{r T} \cdot (1 - e^{-r T}).
\]
For $r=0$ (undiscounted case, as with futures options), the sum simplifies to $L$, giving $\varepsilon_{\text{price}} \leq L \cdot \varepsilon_{\text{semi}}$.

\textbf{Origin.} The dynamic programming operator is non-expansive in the supremum norm (it is a contraction under discounting). Error at each time step is amplified by at most the number of remaining steps.

\section{Error Propagation Chain}

\begin{lemma}[Semigroup Error Propagation]
Assume that for all $f$ in a core $\mathcal{D}$,
\[
\|Q^{(n)}f-\mathcal{L}f\|_\infty \le C_f\varepsilon_n.
\]
If the associated semigroups are contractions, then for $t\in[0,T]$ and sufficiently regular $f$,
\[
\|P_t^{(n)}f-P_tf\|_\infty \le C_{f,T}\varepsilon_n.
\]
Equivalently, local generator error is propagated through the Duhamel formula
\[
P_t^{(n)}f-P_tf = \int_0^t P_{t-s}^{(n)}(Q^{(n)}-\mathcal{L})P_sf\,ds,
\]
with the integral understood first on the common core and then by closure.
\end{lemma}

\begin{proof}[Gronwall/Duhamel argument]
The contraction property gives $\|P_{t-s}^{(n)}\|\le1$. Bounding the integrand by the core consistency estimate and integrating over $[0,t]$ yields a global $O(t\varepsilon_n)$ estimate. For time-dependent value iterations, the non-expansiveness of the Snell operator gives the corresponding discrete Gronwall accumulation over exercise dates.
\end{proof}

The five error sources do not act independently; they form a cascade from construction to pricing:

\[
\boxed{\begin{array}{ccccccc}
\underbrace{\varepsilon_{\text{gen}}}_{\mathcal{O}(h^{2})}
& \longrightarrow &
\underbrace{\varepsilon_{\text{semi}}}_{\mathcal{O}(T h^{2})}
& \longrightarrow &
\underbrace{\varepsilon_{\text{price}}}_{\mathcal{O}(L T h^{2})}
& \longrightarrow &
V_{0}^{(N,L)} \\[6pt]
& & & & \uparrow & & \uparrow \\[6pt]
& & & & \varepsilon_{\text{berm}} & + & \varepsilon_{\text{bound}} \\[6pt]
& & & & \mathcal{O}(1/\sqrt{L}) & & \mathcal{O}(h)
\end{array}}
\]

The chain reads as follows:
\begin{enumerate}[label=(\arabic*)]
    \item \textbf{Generator error $\varepsilon_{\text{gen}} = \mathcal{O}(h^{2})$:} The CTMC generator $Q^{(N)}$ approximates $\mathcal{L}$ pointwise on a core, with local truncation error $\mathcal{O}(h^{2})$.
    \item \textbf{Semigroup error $\varepsilon_{\text{semi}} = \mathcal{O}(T h^{2})$:} Via Trotter--Kato, the semigroups converge at the same rate, with error growing linearly in $T$.
    \item \textbf{Pricing error $\varepsilon_{\text{price}}$:} The transition error propagates through $L$ backward DP steps, accumulating to $\mathcal{O}(L T h^{2})$. For $r>0$, discounting attenuates this by the annuity factor.
    \item \textbf{Bermudan error $\varepsilon_{\text{berm}} = \mathcal{O}(1/\sqrt{L})$:} Independent of spatial discretization; determined solely by the number of exercise dates.
    \item \textbf{Boundary error $\varepsilon_{\text{bound}} = \mathcal{O}(h)$:} Independent of temporal discretization; determined by the grid truncation and boundary treatment.
\end{enumerate}

\section{Combined Error and Optimal Resource Allocation}

To leading order, the total error is
\[
\boxed{\varepsilon_{\text{total}} \leq \underbrace{C_{1} T h^{2}}_{\varepsilon_{\text{gen}}} \;+\; \underbrace{C_{2} T^{2} h^{2}}_{\varepsilon_{\text{semi}}} \;+\; \underbrace{C_{3} \frac{\sqrt{T}}{\sqrt{L}}}_{\varepsilon_{\text{berm}}} \;+\; \underbrace{C_{4} h}_{\varepsilon_{\text{bound}}} \;+\; \underbrace{C_{5} L T h^{2}}_{\varepsilon_{\text{price}}}}.
\]

For practical parameter values ($T \sim 0.5$, $\sigma \sim 0.3$, $h \sim 10^{-2}$, $L \sim 10^{3}$):
\begin{itemize}
    \item $\varepsilon_{\text{gen}} \sim 10^{-4}$ (small, second-order in $h$),
    \item $\varepsilon_{\text{semi}} \sim 5 \times 10^{-5}$ (small, but grows with $T$),
    \item $\varepsilon_{\text{berm}} \sim 3 \times 10^{-3}$ (the dominant term in most cases),
    \item $\varepsilon_{\text{bound}} \sim 10^{-2}$ (the largest term when the grid is truncated aggressively),
    \item $\varepsilon_{\text{price}} \sim 10^{-4}$ (the propagation amplifies $\varepsilon_{\text{semi}}$ by $L$).
\end{itemize}

\textbf{Optimal resource allocation:} To achieve a target accuracy $\varepsilon$, balance the three main terms:
\[
h \sim \varepsilon^{1/4} \quad (\text{from } \varepsilon_{\text{gen}} + \varepsilon_{\text{bound}} \sim h^{2} + h), \qquad
L \sim \varepsilon^{-2} \quad (\text{from } \varepsilon_{\text{berm}} \sim 1/\sqrt{L}).
\]
With this scaling, the total computational cost of the tridiagonal DP scheme is $\mathcal{O}(N L) = \mathcal{O}(\varepsilon^{-1/4} \cdot \varepsilon^{-2}) = \mathcal{O}(\varepsilon^{-9/4})$, which is sub-quadratic in $1/\varepsilon$.

\section{Summary: Error Layers, Rates, and Mitigations}

\begin{center}
\begin{tabular}{>{\raggedright\arraybackslash}p{2.2cm}>{\raggedright\arraybackslash}p{2.2cm}>{\raggedright\arraybackslash}p{2cm}>{\raggedright\arraybackslash}p{3cm}>{\raggedright\arraybackslash}p{2.5cm}}
\toprule
\textbf{Error Layer} & \textbf{Rate} & \textbf{Dominant?} & \textbf{Origin} & \textbf{Mitigation} \\
\midrule
Generator & $\mathcal{O}(h^{2})$ & Usually not & Taylor remainder & Refine grid (larger $N$) \\
Semigroup & $\mathcal{O}(T h^{2})$ & For large $T$ & Trotter--Kato bound & Higher-order schemes \\
Bermudan & $\mathcal{O}(1/\sqrt{L})$ & Often yes & Discretized exercise & Richardson extrapolation \\
Boundary & $\mathcal{O}(h)$ & Often yes & Grid truncation & Local refinement near barriers \\
Pricing (propagation) & $\mathcal{O}(L T h^{2})$ & For large $L$ & DP error accumulation & Discounting, uniformization \\
\bottomrule
\end{tabular}
\end{center}

\begin{remark}
The Bermudan error $\mathcal{O}(1/\sqrt{L})$ is typically the bottleneck for high-precision pricing (basis-point accuracy). Richardson extrapolation with two grid resolutions $L$ and $2L$ can improve the rate to $\mathcal{O}(1/L)$ at the cost of an additional DP evaluation.
\end{remark}

\begin{remark}
The boundary error $\mathcal{O}(h)$ motivates non-uniform grids with higher resolution near economically important points (strike price, initial price). A smoothly varying grid with $\rho_{i} \ll 1$ (see Appendix~\ref{app:consistency}) preserves the interior $\mathcal{O}(h^{2})$ rate while reducing the boundary layer thickness.
\end{remark}

% ═══════════════════════════════════════════════════════════════
% PART VII — ADVANCED TOPICS
% ═══════════════════════════════════════════════════════════════
\part{Advanced Topics}

\begin{quote}
\itshape This part collects topics that deepen the theoretical understanding of the CTMC approximation but are not required on a first reading. The chapters herein cover rigorous operator theory (Hille--Yosida, Trotter--Kato, martingale problem), spectral analysis, numerical stability (upwind scheme, uniformization), and advanced PDE theory (viscosity solutions, Feynman--Kac derivation). Each chapter is self-contained and can be read independently.
\end{quote}

\chapter{CEV Model: Complete Generator}

\begin{theorem}[CEV Generator]
For $dF_{t}/F_{t} = \sigma (F_{t}/F_{0})^{\beta} dW_{t}$,
\[
\calL_{\text{CEV}} f(F) = \frac{1}{2} \sigma^{2} \frac{F^{2\beta+2}}{F_{0}^{2\beta}} f''(F).
\]
\end{theorem}

\medskip
\noindent\textbf{Why this matters.} The CEV (Constant Elasticity of Variance) generator contains only a second-derivative term, no first-derivative term---because the futures price is a martingale under the risk-neutral measure (zero drift). The parameter $\beta$ controls the elasticity of volatility with respect to the price level: $\beta = 0$ recovers GBM (volatility independent of price), $\beta < 0$ means volatility falls as the price rises (leverage effect), and $\beta > 0$ means volatility rises with the price. The factor $F^{2\beta+2} / F_0^{2\beta}$ ensures that at $F = F_0$ the volatility reverts to $\sigma^2$.

\begin{theorem}[CEV CTMC Generator]
On a grid $\{F_{i}\}$ with $D_{i} = \sigma^{2} F_{i}^{2\beta+2} / F_{0}^{2\beta}$,
\[
q_{i,i+1} = \frac{D_{i}}{h_{i}(h_{i} + h_{i-1})}, \quad
q_{i,i-1} = \frac{D_{i}}{h_{i-1}(h_{i} + h_{i-1})}, \quad
q_{i,i} = -\frac{D_{i}}{h_{i} h_{i-1}}.
\]
\end{theorem}

\medskip
\noindent\textbf{Why this matters.} The CTMC generator for the CEV model is especially clean---since there is no drift ($\mu = 0$), both $q_{i,i+1}$ and $q_{i,i-1}$ depend only on $D_i$ (the diffusion variance). The difference between the up and down jump rates is zero, guaranteeing that the CTMC is also a martingale. Note that the diagonal $q_{i,i}$ equals $-(D_i/h_i h_{i-1})$, which exceeds the sum of the left and right rates---it simultaneously cancels the total probability flow from both directions.

\chapter{Two-State Regime-Switching: Complete Construction}

For a two-regime CEV model with parameters $\theta = (\sigma_{1}, \beta_{1}, \sigma_{2}, \beta_{2}, \lambda_{12}, \lambda_{21})$:
\begin{enumerate}[label=(\arabic*)]
    \item Compute per-regime diffusion variance $D_{i}^{(m)} = \sigma_{m}^{2} F_{i}^{2\beta_{m}+2} / F_{0}^{2\beta_{m}}$.
    \item Construct tridiagonal matrices $\Lambda^{(1)}, \Lambda^{(2)}$ as in Chapter~13.
    \item Assemble the $2N \times 2N$ block generator:
    \[
    Q = \begin{pmatrix}
    \Lambda^{(1)} - \lambda_{12} I_{N} & \lambda_{12} I_{N} \\[4pt]
    \lambda_{21} I_{N} & \Lambda^{(2)} - \lambda_{21} I_{N}
    \end{pmatrix}.
    \]
\end{enumerate}

\chapter{Poisson Random Measures: A Rigorous Construction}

\section{Motivation and Definition}

The jump component of a semimartingale is modeled via a point process on a product space. We require a measure-theoretic object that simultaneously encodes the timing and magnitude of jumps.

\begin{definition}[Point Process]
Let $(E, \calE)$ be a measurable space. A point process on $E$ is a random counting measure: a measurable mapping $\xi: \Omega \to \M(E)$, where $\M(E)$ is the space of $\N_{0} \cup \{\infty\}$-valued measures on $E$ that are expressible as countable sums of Dirac masses.
\end{definition}

\medskip
\noindent\textbf{Why this matters.} A point process is a counter of randomly scattered points. Picture randomly dropping beans on a two-dimensional map---the point process $\xi(B)$ tells you how many beans landed in region $B$. Dirac masses mean each bean occupies an exact point location. In jump-diffusion models, point processes describe the timing and size of jumps---each bean corresponds to a jump, with its $x$-coordinate marking the time and its $y$-coordinate marking the jump magnitude.

\begin{definition}[Poisson Random Measure]
Let $(E, \calE, \nu)$ be a $\sigma$-finite measure space. A Poisson random measure $N$ on $[0, \infty) \times E$ with intensity measure $dt \otimes \nu(dy)$ is a point process satisfying:
\begin{enumerate}[label=(\arabic*)]
    \item For any $A \in \calE$ with $\nu(A) < \infty$, the process $\{N([0,t] \times A)\}_{t \geq 0}$ is a Poisson process with intensity $\nu(A)$.
    \item For any disjoint sets $A_{1}, \ldots, A_{k} \in \calE$ with $\nu(A_{i}) < \infty$, the processes $\{N([0,t] \times A_{1})\}, \ldots, \{N([0,t] \times A_{k})\}$ are independent.
\end{enumerate}
\end{definition}

\medskip
\noindent\textbf{Why this matters.} A Poisson random measure encodes time and space together into a single measure. The intuition: jumps are generated along the time axis at rate $dt$, and their sizes are assigned along the space axis according to the distribution $\nu(dy)$, independently of each other. The first condition says that for a fixed-size set $A$, the count of jumps landing in it is a Poisson process; the second says that jump counts in disjoint size sets are independent. This is like separate thunderclouds producing lightning independently, each at its own characteristic frequency.

\begin{proposition}[Expectation of a Poisson Random Measure]
For any nonnegative measurable $g: [0,\infty) \times E \times \Omega \to \R$ that is predictable,
\[
\E\left[\int_{0}^{t} \int_{E} g(s, y) N(ds, dy)\right] = \E\left[\int_{0}^{t} \int_{E} g(s, y) \nu(dy) ds\right].
\]
In differential notation: $\E[N(dt, dy)] = \nu(dy) dt$.
\end{proposition}

\medskip
\noindent\textbf{Why this matters.} This expectation formula is the cornerstone of the compensator concept. The expectation of the Poisson random measure exactly equals its intensity measure---the average frequency and size of random jumps are governed by $\nu(dy)dt$. This is crucial for CTMC construction: the contributions of jumps to the first and second moments are computed precisely through this expectation formula.

\begin{definition}[Compensated Poisson Random Measure]
The compensated Poisson random measure is
\[
\widetilde{N}(dt, dy) := N(dt, dy) - \nu(dy) dt.
\]
\end{definition}

\medskip
\noindent\textbf{Why this matters.} The compensated measure $\widetilde{N}$ is the zero-centered jump noise. The raw measure $N$ has a positive expectation (jump counts are nonnegative), but after subtracting the compensator $\nu(dy)dt$, the expectation of $\widetilde{N}$ becomes zero---it becomes a martingale. This is analogous to accounting profit after deducting expected costs: the pure unexpected gain or loss.

\begin{proposition}[Martingale Property of the Compensated Measure]
For any predictable $g$ with $\E[\int_{0}^{t} \int_{E} |g(s,y)| \nu(dy) ds] < \infty$, the process
\[
M_{t} := \int_{0}^{t} \int_{E} g(s, y) \widetilde{N}(ds, dy)
\]
is a martingale. In particular, $\E[M_{t}] = 0$ for all $t$, and $\E[\widetilde{N}(dt, dy)] = 0$.
\end{proposition}

\medskip
\noindent\textbf{Why this matters.} The martingale property of the compensated measure guarantees that stochastic integrals built with $\widetilde{N}$ have zero expectation under the risk-neutral measure---a prerequisite for no-arbitrage pricing. In the jump-diffusion generator, the compensated-form integral $\int [f(x+h) - f(x)] \widetilde{N}(dt,dy)$ vanishes upon taking conditional expectation, leaving only the predictable $\nu(dy)dt$ integral, which is exactly the jump-integral term in the generator.

\begin{proof}
First consider $g(s, y) = \1_{[0,t]}(s) \1_{A}(y)$ with $A \in \calE$, $\nu(A) < \infty$. Then $\E[\int_{0}^{t} \int_{E} \1_{A}(y) N(ds, dy)] = \E[N([0,t] \times A)] = \nu(A) t = \int_{0}^{t} \int_{E} \1_{A}(y) \nu(dy) ds$. Extend to simple functions, then to general predictable integrands by monotone convergence and linearity.
\end{proof}

\begin{definition}[Compensated Poisson Random Measure]
The compensated Poisson random measure is
\[
\widetilde{N}(dt, dy) := N(dt, dy) - \nu(dy) dt.
\]
\end{definition}

\begin{proposition}[Martingale Property of the Compensated Measure]
For any predictable $g$ with $\E[\int_{0}^{t} \int_{E} |g(s,y)| \nu(dy) ds] < \infty$, the process
\[
M_{t} := \int_{0}^{t} \int_{E} g(s, y) \widetilde{N}(ds, dy)
\]
is a martingale. In particular, $\E[M_{t}] = 0$ for all $t$, and $\E[\widetilde{N}(dt, dy)] = 0$.
\end{proposition}

\begin{proof}
Write $\int g \widetilde{N} = \int g N - \int g \nu(dy) ds$. The expectation of the first term equals the expectation of the second (by the previous proposition), so $\E[M_{t}] = 0$. The martingale property follows from the independent-increments property of the Poisson random measure and the predictability of the integrand.
\end{proof}

\begin{proposition}[It\^{o} Isometry for Compensated Poisson Integrals]
For a predictable $g$ with $\E[\int_{0}^{t} \int_{E} g(s,y)^{2} \nu(dy) ds] < \infty$,
\[
\E\left[\left(\int_{0}^{t} \int_{E} g(s,y) \widetilde{N}(ds, dy)\right)^{2}\right] = \E\left[\int_{0}^{t} \int_{E} g(s,y)^{2} \nu(dy) ds\right].
\]
\end{proposition}

\medskip
\noindent\textbf{Why this matters.} The It\^{o} isometry formula tells us that the second moment of a stochastic integral with respect to the compensated measure $\widetilde{N}$ equals the integral of $g^2$ against the intensity measure $\nu$. This parallels the ordinary-calculus fact that the variance of $\int g dW_t$ equals $\int g^2 dt$. This result is directly usable for computing the variance contribution of the jump component---it supplies the theoretical value of $M_2^J$.

\begin{proposition}[Moments of a Compound Poisson Process]
For $J_{t}$ as above:
\[
\E[J_{t}] = \lambda t \, \E[Y_{1}], \qquad \Var(J_{t}) = \lambda t \, \E[Y_{1}^{2}].
\]
In infinitesimal form: $\E[dJ_{t} \mid \F_{t}] = \lambda \E[Y_{1}] dt$, and $\E[(dJ_{t})^{2} \mid \F_{t}] = \lambda \E[Y_{1}^{2}] dt$.
\end{proposition}

\medskip
\noindent\textbf{Why this matters.} The moment formulas for a compound Poisson process give the direct contribution of jumps to drift and volatility. The first moment $\lambda \mathbb{E}[Y_1]$ is the average jump trend (jump size times jump frequency), and the second moment $\lambda \mathbb{E}[Y_1^2]$ is the volatility contribution of the jumps. These are precisely the continuous-limit counterparts of $M_1^J$ and $M_2^J$ defined earlier---the closed-form expressions for the double-exponential case are substituted directly in the derivations that follow.

\section{Compound Poisson Processes as Special Cases}

\begin{example}[Compound Poisson Process]
Let $\{N_{t}\}$ be a Poisson process with intensity $\lambda$, and let $\{Y_{k}\}$ be i.i.d.\ random variables with distribution $F_{Y}$, independent of $\{N_{t}\}$. The compound Poisson process $J_{t} = \sum_{k=1}^{N_{t}} Y_{k}$ can be represented as
\[
J_{t} = \int_{0}^{t} \int_{\R} y \, N(ds, dy),
\]
where $N$ is a Poisson random measure on $[0,\infty) \times \R$ with intensity measure $\lambda dt \otimes F_{Y}(dy)$. The corresponding L\'{e}vy measure is $\nu(dy) = \lambda F_{Y}(dy)$.
\end{example}

\begin{proposition}[Moments of a Compound Poisson Process]
For $J_{t}$ as above:
\[
\E[J_{t}] = \lambda t \, \E[Y_{1}], \qquad \Var(J_{t}) = \lambda t \, \E[Y_{1}^{2}].
\]
In infinitesimal form: $\E[dJ_{t} \mid \F_{t}] = \lambda \E[Y_{1}] dt$, and $\E[(dJ_{t})^{2} \mid \F_{t}] = \lambda \E[Y_{1}^{2}] dt$.
\end{proposition}

\begin{proof}
By the Wald identities for compound Poisson processes. For the first moment: $\E[J_{t}] = \E[N_{t}] \E[Y_{1}] = \lambda t \E[Y_{1}]$. The variance uses $\E[J_{t}^{2}] = \E[N_{t}] \E[Y_{1}^{2}] + \E[N_{t}(N_{t}-1)] (\E[Y_{1}])^{2} = \lambda t \E[Y_{1}^{2}] + \lambda^{2} t^{2} (\E[Y_{1}])^{2}$, hence $\Var(J_{t}) = \E[J_{t}^{2}] - (\E[J_{t}])^{2} = \lambda t \E[Y_{1}^{2}]$.
\end{proof}

\chapter{Complete Derivation of the Moment-Generating Function for Double-Exponential Jumps}
\label{app:dejd-mgf}

\section{The Double-Exponential Distribution}

\begin{definition}[Double-Exponential Distribution]
A random variable $Y$ follows a double-exponential (or Laplace) distribution with parameters $p \in [0,1]$, $\eta_{1} > 0$, $\eta_{2} > 0$ if its probability density function is
\[
f_{Y}(y) = p \eta_{1} e^{-\eta_{1} y} \1_{\{y \geq 0\}} + (1-p) \eta_{2} e^{\eta_{2} y} \1_{\{y < 0\}}.
\]
\end{definition}

\medskip
\noindent\textbf{Why this matters.} The double-exponential distribution is the standard choice for financial jump models---the positive half-axis decays exponentially at rate $\eta_1$, the negative half-axis at rate $\eta_2$, and the probability $p$ controls the mix of up and down jumps. This distribution has two key advantages: (1) its tails are heavier than the normal distribution, capturing the fat tails observed in financial data; (2) its moment-generating function has a closed form, allowing the jump moments $M_1^J$ and $M_2^J$ to be computed analytically.

\begin{proposition}[Normalization]
$\int_{-\infty}^{\infty} f_{Y}(y) dy = 1$.
\end{proposition}

\medskip
\noindent\textbf{Why this matters.} Normalization guarantees that $f_Y$ is indeed a probability density---total probability equals 1. The verification is simple: the positive half-axis contributes $p$, the negative half-axis contributes $1-p$, and they sum to 1. Though this appears trivial, it ensures that all subsequent expectations and moment computations based on $f_Y$ carry the correct probabilistic interpretation.

\begin{proof}
\[
\int_{-\infty}^{\infty} f_{Y}(y) dy = p \eta_{1} \int_{0}^{\infty} e^{-\eta_{1} y} dy + (1-p) \eta_{2} \int_{-\infty}^{0} e^{\eta_{2} y} dy = p \eta_{1} \cdot \frac{1}{\eta_{1}} + (1-p) \eta_{2} \cdot \frac{1}{\eta_{2}} = p + (1-p) = 1.
\]
\end{proof}

\section{Moment-Generating Function}

\begin{theorem}[MGF of the Double-Exponential Distribution]
For $-\eta_{2} < k < \eta_{1}$,
\[
\boxed{\E[e^{kY}] = p \frac{\eta_{1}}{\eta_{1} - k} + (1-p) \frac{\eta_{2}}{\eta_{2} + k}}.
\]
\end{theorem}

\medskip
\noindent\textbf{Why this matters.} The domain of convergence $(-\eta_2, \eta_1)$ is the defining feature of the double-exponential distribution---outside this interval, the exponential moments diverge (e.g., $\mathbb{E}[e^{\eta_1 Y}] = \infty$). This closed form directly yields jump moments of any order: differentiate with respect to $k$ and evaluate at $k=0$ to obtain the corresponding raw moment. In CTMC pricing, we need $\mathbb{E}[e^Y]$ and $\mathbb{E}[e^{2Y}]$ to compute the jump compensator and variance.

\begin{corollary}[First Moment of the Double-Exponential]
\[
\E[Y] = \frac{p}{\eta_{1}} - \frac{1-p}{\eta_{2}}.
\]
\end{corollary}

\medskip
\noindent\textbf{Why this matters.} The first moment vividly exposes the asymmetry of the double-exponential distribution: $p/\eta_1$ is the average contribution of upward jumps, $(1-p)/\eta_2$ is the average contribution of downward jumps. If $p > \eta_1/(\eta_1 + \eta_2)$, the average jump is positive (upside bias); otherwise it is negative. The constraints $\eta_1 > 1$ and $\eta_2 > 0$ guarantee that the mean exists and is finite.

\begin{corollary}[Second Moment and Variance]
\[
\E[Y^{2}] = \frac{2p}{\eta_{1}^{2}} + \frac{2(1-p)}{\eta_{2}^{2}}, \qquad
\Var(Y) = \frac{2p}{\eta_{1}^{2}} + \frac{2(1-p)}{\eta_{2}^{2}} - \left(\frac{p}{\eta_{1}} - \frac{1-p}{\eta_{2}}\right)^{2}.
\]
\end{corollary}

\medskip
\noindent\textbf{Why this matters.} The second moment (and the variance derived from it) measures the dispersion of jump magnitudes. The smaller $\eta_1$, the thicker the positive tail and the larger the second moment; similarly for $\eta_2$ on the negative side. In CTMC construction, $\mathbb{E}[(e^Y - 1)^2]$ is the total contribution of jumps to price variance---a critical input for both parameter estimation and risk measurement.

\begin{corollary}[Jump Compensator for Martingale Pricing]
\[
\zeta := \E[e^{Y}] - 1 = \frac{p\eta_{1}}{\eta_{1} - 1} + \frac{(1-p)\eta_{2}}{\eta_{2} + 1} - 1.
\]
\end{corollary}

\medskip
\noindent\textbf{Why this matters.} $\zeta$ is the jump compensator---it tells us how much to subtract from the drift term of the SDE to keep the futures price a martingale. Intuitively: if the average jump is a positive gain ($\mathbb{E}[e^Y] > 1$), we must deduct $\lambda\zeta dt$ from the drift to offset this positive bias; otherwise the futures price acquires a systematic tilt. This is the concrete meaning of "compensated" in "compensated Poisson process."

\begin{theorem}[Simplified Jump Variance]
For $Y$ following a double-exponential distribution with $\eta_{1} > 2$,
\[
\boxed{\E[(e^{Y} - 1)^{2}] = 2\left(\frac{p}{(\eta_{1} - 1)(\eta_{1} - 2)} + \frac{1-p}{(\eta_{2} + 1)(\eta_{2} + 2)}\right)}.
\]
\end{theorem}

\medskip
\noindent\textbf{Why this matters.} $\mathbb{E}[(e^Y - 1)^2]$ is the pure contribution of jumps to price variance---it is the compact form of $\mathbb{E}[e^{2Y}] - 2\mathbb{E}[e^Y] + 1$. The condition $\eta_1 > 2$ is rigid: if violated, this second moment diverges, the volatility contribution of jumps becomes infinite, and the model loses mathematical meaning. In practical calibration, $\eta_1$ must be bounded below by at least 2.001 for numerical stability.

\begin{proof}
\[
\begin{aligned}
\E[e^{kY}] &= \int_{-\infty}^{\infty} e^{ky} f_{Y}(y) dy \\
&= p \eta_{1} \int_{0}^{\infty} e^{ky} e^{-\eta_{1} y} dy + (1-p) \eta_{2} \int_{-\infty}^{0} e^{ky} e^{\eta_{2} y} dy \\
&= p \eta_{1} \int_{0}^{\infty} e^{-(\eta_{1} - k) y} dy + (1-p) \eta_{2} \int_{-\infty}^{0} e^{(\eta_{2} + k) y} dy.
\end{aligned}
\]

\textbf{First integral (positive branch).} For $\eta_{1} - k > 0$ (i.e., $k < \eta_{1}$):
\[
\int_{0}^{\infty} e^{-(\eta_{1} - k) y} dy = \left[-\frac{1}{\eta_{1} - k} e^{-(\eta_{1} - k) y}\right]_{0}^{\infty} = 0 - \left(-\frac{1}{\eta_{1} - k}\right) = \frac{1}{\eta_{1} - k}.
\]

\textbf{Second integral (negative branch).} For $\eta_{2} + k > 0$ (i.e., $k > -\eta_{2}$):
\[
\int_{-\infty}^{0} e^{(\eta_{2} + k) y} dy = \left[\frac{1}{\eta_{2} + k} e^{(\eta_{2} + k) y}\right]_{-\infty}^{0} = \frac{1}{\eta_{2} + k} - 0 = \frac{1}{\eta_{2} + k}.
\]

Multiplying by the prefactors:
\[
\E[e^{kY}] = p \eta_{1} \cdot \frac{1}{\eta_{1} - k} + (1-p) \eta_{2} \cdot \frac{1}{\eta_{2} + k} = p \frac{\eta_{1}}{\eta_{1} - k} + (1-p) \frac{\eta_{2}}{\eta_{2} + k}.
\]

The domain of convergence is determined by the simultaneous validity of both integrals: $-\eta_{2} < k < \eta_{1}$.
\end{proof}

\begin{corollary}[First Moment of the Double-Exponential]
\[
\E[Y] = \frac{p}{\eta_{1}} - \frac{1-p}{\eta_{2}}.
\]
\end{corollary}

\begin{proof}
Differentiate the MGF and evaluate at $k = 0$, or compute directly:
\[
\E[Y] = p \eta_{1} \int_{0}^{\infty} y e^{-\eta_{1} y} dy + (1-p) \eta_{2} \int_{-\infty}^{0} y e^{\eta_{2} y} dy.
\]
For the positive branch: $\int_{0}^{\infty} y e^{-\eta_{1} y} dy = 1/\eta_{1}^{2}$ (integration by parts). For the negative branch: $\int_{-\infty}^{0} y e^{\eta_{2} y} dy = -1/\eta_{2}^{2}$. Hence $\E[Y] = p \eta_{1} \cdot 1/\eta_{1}^{2} + (1-p) \eta_{2} \cdot (-1/\eta_{2}^{2}) = p/\eta_{1} - (1-p)/\eta_{2}$.
\end{proof}

\begin{corollary}[Second Moment and Variance]
\[
\E[Y^{2}] = \frac{2p}{\eta_{1}^{2}} + \frac{2(1-p)}{\eta_{2}^{2}}, \qquad
\Var(Y) = \frac{2p}{\eta_{1}^{2}} + \frac{2(1-p)}{\eta_{2}^{2}} - \left(\frac{p}{\eta_{1}} - \frac{1-p}{\eta_{2}}\right)^{2}.
\]
\end{corollary}

\begin{proof}
$\E[Y^{2}] = p \eta_{1} \int_{0}^{\infty} y^{2} e^{-\eta_{1} y} dy + (1-p) \eta_{2} \int_{-\infty}^{0} y^{2} e^{\eta_{2} y} dy$. For the positive branch: $\int_{0}^{\infty} y^{2} e^{-\eta_{1} y} dy = 2/\eta_{1}^{3}$. For the negative branch: $\int_{-\infty}^{0} y^{2} e^{\eta_{2} y} dy = 2/\eta_{2}^{3}$. Hence $\E[Y^{2}] = 2p/\eta_{1}^{2} + 2(1-p)/\eta_{2}^{2}$.
\end{proof}

\section{Complete Simplification of $\E[(e^{Y} - 1)^{2}]$}

\begin{theorem}[Simplified Jump Variance]
For $Y$ following a double-exponential distribution with $\eta_{1} > 2$,
\[
\boxed{\E[(e^{Y} - 1)^{2}] = 2\left(\frac{p}{(\eta_{1} - 1)(\eta_{1} - 2)} + \frac{1-p}{(\eta_{2} + 1)(\eta_{2} + 2)}\right)}.
\]
\end{theorem}

\begin{proof}
We provide the complete algebraic derivation with every intermediate step shown.

\textbf{Step 1 --- Expand the square.}
\[
\E[(e^{Y} - 1)^{2}] = \E[e^{2Y} - 2e^{Y} + 1] = \E[e^{2Y}] - 2\E[e^{Y}] + 1.
\]

\textbf{Step 2 --- Evaluate using the MGF.} Applying Theorem~H.2 with $k=2$ (requires $\eta_{1} > 2$) and $k=1$:
\[
\begin{aligned}
\E[e^{2Y}] &= p \frac{\eta_{1}}{\eta_{1} - 2} + (1-p) \frac{\eta_{2}}{\eta_{2} + 2}, \\
\E[e^{Y}] &= p \frac{\eta_{1}}{\eta_{1} - 1} + (1-p) \frac{\eta_{2}}{\eta_{2} + 1}.
\end{aligned}
\]

\textbf{Step 3 --- Substitute and group by $p$ and $(1-p)$.}
\[
\begin{aligned}
\E[(e^{Y} - 1)^{2}] &= \left[p \frac{\eta_{1}}{\eta_{1} - 2} + (1-p) \frac{\eta_{2}}{\eta_{2} + 2}\right] - 2\left[p \frac{\eta_{1}}{\eta_{1} - 1} + (1-p) \frac{\eta_{2}}{\eta_{2} + 1}\right] + 1 \\
&= p\left[ \frac{\eta_{1}}{\eta_{1} - 2} - \frac{2\eta_{1}}{\eta_{1} - 1} \right] + (1-p)\left[ \frac{\eta_{2}}{\eta_{2} + 2} - \frac{2\eta_{2}}{\eta_{2} + 1} \right] + 1.
\end{aligned}
\]

\textbf{Step 4 --- Simplify the $p$-bracket.} Combine over a common denominator:
\[
\begin{aligned}
\frac{\eta_{1}}{\eta_{1} - 2} - \frac{2\eta_{1}}{\eta_{1} - 1}
&= \frac{\eta_{1}(\eta_{1} - 1) - 2\eta_{1}(\eta_{1} - 2)}{(\eta_{1} - 1)(\eta_{1} - 2)} \\
&= \frac{\eta_{1}^{2} - \eta_{1} - 2\eta_{1}^{2} + 4\eta_{1}}{(\eta_{1} - 1)(\eta_{1} - 2)} \\
&= \frac{-\eta_{1}^{2} + 3\eta_{1}}{(\eta_{1} - 1)(\eta_{1} - 2)}.
\end{aligned}
\]

Now add the $+1 = +p \cdot 1$ that appears as part of the $+p$ in the expansion (since $1 = p + (1-p)$):
\[
\begin{aligned}
p\left[ \frac{-\eta_{1}^{2} + 3\eta_{1}}{(\eta_{1} - 1)(\eta_{1} - 2)} + 1 \right]
&= p\left[ \frac{-\eta_{1}^{2} + 3\eta_{1} + (\eta_{1} - 1)(\eta_{1} - 2)}{(\eta_{1} - 1)(\eta_{1} - 2)} \right] \\
&= p\left[ \frac{-\eta_{1}^{2} + 3\eta_{1} + \eta_{1}^{2} - 3\eta_{1} + 2}{(\eta_{1} - 1)(\eta_{1} - 2)} \right] \\
&= p \cdot \frac{2}{(\eta_{1} - 1)(\eta_{1} - 2)}.
\end{aligned}
\]

\textbf{Step 5 --- Simplify the $(1-p)$-bracket.} Analogously:
\[
\begin{aligned}
\frac{\eta_{2}}{\eta_{2} + 2} - \frac{2\eta_{2}}{\eta_{2} + 1}
&= \frac{\eta_{2}(\eta_{2} + 1) - 2\eta_{2}(\eta_{2} + 2)}{(\eta_{2} + 1)(\eta_{2} + 2)} \\
&= \frac{\eta_{2}^{2} + \eta_{2} - 2\eta_{2}^{2} - 4\eta_{2}}{(\eta_{2} + 1)(\eta_{2} + 2)} \\
&= \frac{-\eta_{2}^{2} - 3\eta_{2}}{(\eta_{2} + 1)(\eta_{2} + 2)}.
\end{aligned}
\]

Adding the $+1 = +(1-p)$ part:
\[
\begin{aligned}
(1-p)\left[ \frac{-\eta_{2}^{2} - 3\eta_{2}}{(\eta_{2} + 1)(\eta_{2} + 2)} + 1 \right]
&= (1-p)\left[ \frac{-\eta_{2}^{2} - 3\eta_{2} + (\eta_{2} + 1)(\eta_{2} + 2)}{(\eta_{2} + 1)(\eta_{2} + 2)} \right] \\
&= (1-p)\left[ \frac{-\eta_{2}^{2} - 3\eta_{2} + \eta_{2}^{2} + 3\eta_{2} + 2}{(\eta_{2} + 1)(\eta_{2} + 2)} \right] \\
&= (1-p) \cdot \frac{2}{(\eta_{2} + 1)(\eta_{2} + 2)}.
\end{aligned}
\]

\textbf{Step 6 --- Sum the contributions.}
\[
\E[(e^{Y} - 1)^{2}] = \frac{2p}{(\eta_{1} - 1)(\eta_{1} - 2)} + \frac{2(1-p)}{(\eta_{2} + 1)(\eta_{2} + 2)}.
\]

Factoring out the common factor $2$ yields the stated result.
\end{proof}

\begin{remark}[Required Domain Condition]
The expression for $\E[e^{2Y}]$ requires $2 < \eta_{1}$ for the MGF to be finite at $k=2$. If $\eta_{1} \leq 2$, the second moment of $e^{Y}$ diverges, and the model is not suitable for pricing applications (infinite quadratic variation). The constraint $\eta_{1} > 2$ (strengthening $\eta_{1} > 1$) must be enforced in parameter estimation.
\end{remark}

\section{The Compensator $\zeta$}

\begin{corollary}[Jump Compensator for Martingale Pricing]
\[
\zeta := \E[e^{Y}] - 1 = \frac{p\eta_{1}}{\eta_{1} - 1} + \frac{(1-p)\eta_{2}}{\eta_{2} + 1} - 1.
\]
\end{corollary}

The drift term $-\lambda \zeta dt$ in the CEV-DEJD SDE ensures $\E^{\Q}[dF_{t} \mid \F_{t}] = 0$, i.e., the futures price process is a martingale under the risk-neutral measure. This is the no-arbitrage condition for futures.

\begin{proof}[Verification of the Martingale Condition]
\[
\begin{aligned}
\E^{\Q}[dF_{t} \mid \F_{t}] &= \E^{\Q}\left[F_{t}\left(-\lambda \zeta dt + \sigma\left(\frac{F_{t}}{F_{0}}\right)^{\beta} dW_{t} + d\left(\sum_{i=1}^{N_{t}}(e^{Y_{i}}-1)\right)\right) \mid \F_{t}\right] \\
&= F_{t}\left(-\lambda \zeta dt + 0 + \lambda \E[e^{Y} - 1] dt\right) \\
&= F_{t}\left(-\lambda(\E[e^{Y}] - 1) + \lambda(\E[e^{Y}] - 1)\right) dt = 0.
\end{aligned}
\]
The stochastic integral $\int \sigma(\cdot) dW$ is a martingale (expectation zero), and the compensated jump term has expectation $\lambda \E[e^{Y}-1] dt = \lambda \zeta dt$, which is exactly canceled by the explicit drift correction $-\lambda \zeta dt$.
\end{proof}

\chapter{Complete Case-by-Case Jump Integration for Double-Exponential Jumps}
\label{app:case-jump}

We provide the exhaustive derivation of each case in the jump transition rate computation, with every integration step shown explicitly.

\section{Notation and Setup}

Recall from Chapter~14: for a jump from state $i$ (price $x_{i}$) to state $j$, the relative jump size $Z = e^{Y} - 1$ must fall in the interval $(\alpha_{i}(j-1), \alpha_{i}(j)]$, where
\[
\alpha_{i}(k) = \frac{\frac{1}{2}(x_{k} + x_{k+1})}{x_{i}} - 1, \qquad k = 0, \ldots, N,
\]
with $\alpha_{i}(0) = -1$ and $\alpha_{i}(N) = \infty$. The jump transition rate is
\[
q_{ij}^{J} = \lambda(x_{i}) \int_{\alpha_{i}(j-1)}^{\alpha_{i}(j)} f_{Z}(z) dz,
\]
with $f_{Z}(z) = p \eta_{1} (z+1)^{-1-\eta_{1}} \1_{\{z \geq 0\}} + (1-p) \eta_{2} (z+1)^{\eta_{2}-1} \1_{\{-1 < z < 0\}}$.

\section{Primitive Functions}

\begin{lemma}[Primitives for Jump Integration]
\[
\begin{aligned}
\int p \eta_{1} (z+1)^{-1-\eta_{1}} dz &= -p (z+1)^{-\eta_{1}} + C, \\
\int (1-p) \eta_{2} (z+1)^{\eta_{2}-1} dz &= (1-p) (z+1)^{\eta_{2}} + C.
\end{aligned}
\]
\end{lemma}

\medskip
\noindent\textbf{Why this matters.} These two primitive functions are the foundation of jump-transition-rate computation. $\int (z+1)^{-1-\eta_1} dz = -(z+1)^{-\eta_1}/\eta_1$ is a power-law integral; the prefactor $p\eta_1$ cancels the $\eta_1$ in the denominator. Note that the antiderivative of a negative power is decreasing---this means that on an interval $[a,b]$, the integral $F(b) - F(a) = -p(b+1)^{-\eta_1} + p(a+1)^{-\eta_1} > 0$ (since $b > a$ implies $(b+1)^{-\eta_1} < (a+1)^{-\eta_1}$), guaranteeing positive jump rates.

\begin{theorem}[Case 1 Integration]
\[
\boxed{q_{ij}^{J} = \lambda(x_{i}) \, p \left[(\alpha_{i}(j-1)+1)^{-\eta_{1}} - (\alpha_{i}(j)+1)^{-\eta_{1}}\right]}.
\]
\end{theorem}

\medskip
\noindent\textbf{Why this matters.} Case 1 handles pure upward jumps---both the starting and landing points lie above the current price. Only the positive branch of the double-exponential ($p\eta_1$ part) is used, yielding a simple power-law difference. The key observation is that $\alpha_i(j-1)+1$ and $\alpha_i(j)+1$ are the relative price ratios $x_j/x_i$ and $x_{j+1}/x_i$---the farther the grid cell, the smaller the power-law value and the smaller its contribution.

\begin{theorem}[Case 2 Integration]
\[
\boxed{q_{ij}^{J} = \lambda(x_{i}) (1-p) \left[(\alpha_{i}(j)+1)^{\eta_{2}} - (\alpha_{i}(j-1)+1)^{\eta_{2}}\right]}.
\]
\end{theorem}

\medskip
\noindent\textbf{Why this matters.} Case 2 is the symmetric pure downward jump---both endpoints lie below the current price. Here the negative branch ($(1-p)\eta_2$ part) is used, producing a positive-power difference. Since $\eta_2 > 0$, $(z+1)^{\eta_2}$ is an increasing function, so $q_{ij}^J > 0$ holds automatically.

\begin{theorem}[Case 3 Integration]
\[
\boxed{\begin{aligned}
q_{ij}^{J} = \lambda(x_{i}) \Big\{ &(1-p)\left[1 - (\alpha_{i}(j-1)+1)^{\eta_{2}}\right] \\
&+ p\left[1 - (\alpha_{i}(j)+1)^{-\eta_{1}}\right]\Big\}.
\end{aligned}}
\]
\end{theorem}

\medskip
\noindent\textbf{Why this matters.} Case 3 is the cross-zero jump---the integration interval straddles $z=0$ and must be split into negative and positive halves. The first segment (negative half) runs from $\alpha_i(j-1)$ to 0, contributing $(1-p)[1 - (\alpha_i(j-1)+1)^{\eta_2}]$; the second segment (positive half) runs from 0 to $\alpha_i(j)$, contributing $p[1 - (\alpha_i(j)+1)^{-\eta_1}]$. The origin of the $1$s is $(0+1)^{\eta_2} = 1$ and $(0+1)^{-\eta_1} = 1$---both power functions evaluate to 1 at zero.

\begin{corollary}[Diagonal Completion]
For completeness, the diagonal entries of the jump generator are set to
\[
q_{ii}^{J} = -\sum_{j \neq i} q_{ij}^{J},
\]
ensuring the row-sum-zero condition (C2) for the jump generator $Q_{J}^{(N)}$.
\end{corollary}

\medskip
\noindent\textbf{Why this matters.} The diagonal entry $q_{ii}^J$ is not defined by "jumping to oneself" but as the negative row sum that guarantees total probability-flow conservation. If $q_{ii}^J$ is $-10$, it means the total exit rate from state $i$ is 10 units---every decay in the probability of staying put is balanced by the sum of the rates of jumping to other states.

\chapter{Complete Derivation of the Piccioni Formula for Diffusion CTMC Construction}

\section{Motivation}

\textit{Piccioni (1987)} introduced a systematic method for constructing CTMC approximations of diffusion processes by decomposing the transition rates into components that separately control the drift and diffusion moments. This appendix provides the complete derivation.

\section{Decomposition Ansatz}

\begin{definition}[Piccioni Decomposition]
For a diffusion process with drift $\mu(x)$ and diffusion coefficient $\sigma^{2}(x)$, define auxiliary functions $\delta_{+}(x)$, $\delta_{-}(x)$, and $\eta(x)$ such that
\[
w_{n}(x) = \delta_{+}(x) + \eta(x), \qquad v_{n}(x) = \delta_{-}(x) + \eta(x),
\]
and the transition rates are
\[
q_{i,i+1} = \frac{w_{n}(x_{i})}{x_{i+1} - x_{i}}, \qquad q_{i,i-1} = \frac{v_{n}(x_{i})}{x_{i} - x_{i-1}}.
\]
\end{definition}

\medskip
\noindent\textbf{Why this matters.} The Piccioni decomposition assigns the duties of drift and diffusion to separate components. $\delta_+$ and $\delta_-$ handle the drift direction (first moment), and $\eta$ handles the volatility intensity (second moment). The quantities $w_n$ and $v_n$ are composite force coefficients, divided by the grid spacing to yield the final jump rates. This separation allows independent adjustment of drift and diffusion contributions---especially useful when a large drift threatens to turn one direction's rate negative, in which case $\eta$ can be tuned as a remedy.

\begin{theorem}[Two-State Transition Probabilities --- Full Derivation]
For a CTMC on $\{1,2\}$ with generator $Q = \bigl(\begin{smallmatrix} -\alpha & \alpha \\ \beta & -\beta \end{smallmatrix}\bigr)$, the transition matrix is
\[
P(t) = e^{tQ} = \frac{1}{\alpha+\beta} \begin{pmatrix}
\beta + \alpha e^{-(\alpha+\beta)t} & \alpha(1-e^{-(\alpha+\beta)t}) \\
\beta(1-e^{-(\alpha+\beta)t}) & \alpha + \beta e^{-(\alpha+\beta)t}
\end{pmatrix}.
\]
\end{theorem}

\medskip
\noindent\textbf{Why this matters.} The two-state CTMC is the atom of regime-switching models---every multi-regime model can be viewed as a composite of two-state switches. The closed form tells us that the probability of staying in regime 1 after time $t$ is $(\beta + \alpha e^{-(\alpha+\beta)t})/(\alpha+\beta)$---it starts at $\alpha/(\alpha+\beta)$ (from the stationary distribution) and decays exponentially toward the stationary probability $\beta/(\alpha+\beta)$. This formula eliminates the need to recompute the matrix exponential $e^{tQ}$ each time.

\begin{proof}
We solve the Kolmogorov forward equation $\frac{d}{dt}P(t) = P(t)Q$, $P(0) = I$, element by element.

\textbf{Step 1 --- ODE for $p_{11}(t)$.} From the matrix multiplication:
\[
\frac{d}{dt} \begin{pmatrix} p_{11} & p_{12} \\ p_{21} & p_{22} \end{pmatrix} = \begin{pmatrix} p_{11} & p_{12} \\ p_{21} & p_{22} \end{pmatrix} \begin{pmatrix} -\alpha & \alpha \\ \beta & -\beta \end{pmatrix}.
\]

For the $(1,1)$ entry: $\frac{d}{dt} p_{11}(t) = -\alpha p_{11}(t) + \beta p_{12}(t)$.

Using the stochasticity constraint $p_{11}(t) + p_{12}(t) = 1$ (row sum equals 1), we eliminate $p_{12}$:
\[
\frac{d}{dt} p_{11}(t) = -\alpha p_{11}(t) + \beta(1 - p_{11}(t)) = -(\alpha+\beta) p_{11}(t) + \beta.
\]

\textbf{Step 2 --- Solve the first-order linear ODE.} The homogeneous equation $\frac{d}{dt}p_{11}^{h} = -(\alpha+\beta) p_{11}^{h}$ has solution $p_{11}^{h}(t) = C e^{-(\alpha+\beta)t}$.

A particular constant solution: set $\frac{d}{dt}p_{11}^{p} = 0$, giving $0 = -(\alpha+\beta)p_{11}^{p} + \beta$, so $p_{11}^{p} = \frac{\beta}{\alpha+\beta}$.

The general solution is $p_{11}(t) = C e^{-(\alpha+\beta)t} + \frac{\beta}{\alpha+\beta}$.

\textbf{Step 3 --- Apply the initial condition.} $p_{11}(0) = 1$ (the process starts in state 1 with probability 1):
\[
C + \frac{\beta}{\alpha+\beta} = 1 \;\Longrightarrow\; C = 1 - \frac{\beta}{\alpha+\beta} = \frac{\alpha}{\alpha+\beta}.
\]

Thus:
\[
p_{11}(t) = \frac{\alpha}{\alpha+\beta} e^{-(\alpha+\beta)t} + \frac{\beta}{\alpha+\beta} = \frac{\beta + \alpha e^{-(\alpha+\beta)t}}{\alpha+\beta}.
\]

\textbf{Step 4 --- Remaining entries by probability conservation.}
\[
p_{12}(t) = 1 - p_{11}(t) = 1 - \frac{\beta + \alpha e^{-(\alpha+\beta)t}}{\alpha+\beta} = \frac{\alpha(1 - e^{-(\alpha+\beta)t})}{\alpha+\beta}.
\]

By symmetry (exchange regimes 1 and 2, $\alpha \leftrightarrow \beta$):
\[
p_{22}(t) = \frac{\alpha + \beta e^{-(\alpha+\beta)t}}{\alpha+\beta}, \qquad
p_{21}(t) = \frac{\beta(1 - e^{-(\alpha+\beta)t})}{\alpha+\beta}.
\]

\textbf{Step 5 --- Stationary distribution.} Taking $t \to \infty$:
\[
\lim_{t \to \infty} P(t) = \frac{1}{\alpha+\beta} \begin{pmatrix} \beta & \alpha \\ \beta & \alpha \end{pmatrix},
\]
so the stationary distribution is $\pi = (\frac{\beta}{\alpha+\beta}, \frac{\alpha}{\alpha+\beta})$. Each row converges to $\pi$, reflecting ergodicity of the two-state CTMC.
\end{proof}

\chapter{The Martingale Problem: Existence, Uniqueness, and the Markov Property}

\section{Formulation}

\begin{definition}[Martingale Problem --- Ethier-Kurtz]
Let $E$ be a complete separable metric space, $\calL: \calD(\calL) \subseteq B(E) \to B(E)$ a linear operator, and $\nu$ a probability measure on $E$. A progressively measurable process $\{X_{t}\}$ with initial distribution $\nu$ is a \textbf{solution to the martingale problem for $(\calL, \nu)$} if for every $f \in \calD(\calL)$,
\[
M_{t}^{f} := f(X_{t}) - f(X_{0}) - \int_{0}^{t} \calL f(X_{s}) ds
\]
is a martingale with respect to the natural filtration $\{\F_{t}^{X}\}$.
\end{definition}

\medskip
\noindent\textbf{Why this matters.} The martingale problem provides a path that bypasses SDEs and constructs a process directly from its generator. The traditional route is to write down an SDE and then use It\^{o}'s formula to derive the generator. The martingale problem reverses the direction: given the generator $\mathcal{L}$, find a process such that $f(X_t) - \int_0^t \mathcal{L}f(X_s)ds$ is a martingale. For CTMC approximation, this is exactly what we need---we only know the generator $\mathcal{L}$; we never construct an SDE. Instead, we approximate $\mathcal{L}$ directly with $Q^{(N)}$ and rely on the well-posedness of the martingale problem to guarantee that the approximation is valid.

\begin{theorem}[Uniqueness and the Markov Property --- Ethier-Kurtz, Theorem 4.4.2]
If the martingale problem for $(\calL, \nu)$ is well-posed (i.e., has exactly one solution in law for each initial distribution $\nu$), then any solution $\{X_{t}\}$ is a Markov process, and its transition semigroup $\{P_{t}\}$ is uniquely determined by $\calL$. Furthermore, $\calL$ is the infinitesimal generator of $\{P_{t}\}$.
\end{theorem}

\medskip
\noindent\textbf{Why this matters.} Well-posedness is the first domino in the chain of CTMC approximation logic. If the martingale problem for the original process is ill-posed (multiple solutions), then even if the CTMC approximates the generator, there is no guarantee it approximates the right solution. Fortunately, for Lipschitz diffusions and jump-diffusions with bounded jump intensity, the martingale problem is known to be well-posed---so the target process of the CTMC approximation is unique.

\begin{proposition}[Well-Posedness for Lipschitz Diffusions]
If $\mu, \sigma$ are globally Lipschitz, the martingale problem for
\[
\calL f(x) = \mu(x) f'(x) + \frac{1}{2} \sigma^{2}(x) f''(x)
\]
on $C_{c}^{\infty}(\R)$ is well-posed. The unique solution coincides with the strong solution of the SDE $dX_{t} = \mu(X_{t}) dt + \sigma(X_{t}) dW_{t}$.
\end{proposition}

\medskip
\noindent\textbf{Why this matters.} This result confirms the equivalence of the martingale-problem approach and the SDE approach under Lipschitz conditions. For CTMC approximation, it means that if our $Q^{(N)}$ approximates $\mathcal{L}$, then it approximates the strong solution of the SDE---there is no ambiguity. The intuitive meaning of the Lipschitz condition is that the drift and volatility cannot vary faster than linearly---this prevents pathological explosion behavior in finite time.

\begin{proposition}[Well-Posedness for Jump-Diffusions with Bounded Intensity]
If $\mu, \sigma$ are Lipschitz, $\lambda(\cdot)$ is bounded and Lipschitz, and $F_{Y}$ has finite second moment, the martingale problem for the jump-diffusion generator is well-posed.
\end{proposition}

\medskip
\noindent\textbf{Why this matters.} This proposition extends well-posedness to the jump-diffusion setting. The critical extra conditions are that $\lambda$ is bounded and $F_Y$ has a finite second moment---the former prevents jumps from occurring infinitely often per unit time, and the latter keeps the jump contribution to variance finite. If the second moment does not exist (as with certain extreme stable distributions), the volatility contribution of jumps becomes infinite and well-posedness fails---this is the deep reason why $\eta_1 > 2$ is required in earlier theorems.

\begin{proposition}[Well-Posedness for Jump-Diffusions with Bounded Intensity]
If $\mu, \sigma$ are Lipschitz, $\lambda(\cdot)$ is bounded and Lipschitz, and $F_{Y}$ has finite second moment, the martingale problem for the jump-diffusion generator is well-posed.
\end{proposition}

\chapter{Feller Property of the Approximating CTMC}

\begin{theorem}[CTMC Generators Generate Feller Semigroups]
For a finite-state space $\calS^{(N)} = \{x_{1}, \ldots, x_{N}\}$, any valid generator matrix $Q^{(N)}$ (satisfying (C1)--(C2)) generates a Feller semigroup $P_{t}^{(N)} = e^{t Q^{(N)}}$ on $C_{0}(\calS^{(N)}) \cong \R^{N}$, where $C_{0}$ is equipped with the supremum norm.
\end{theorem}

\begin{proof}
For a finite set, $C_{0}(\calS^{(N)}) = \R^{N}$ (all functions on a finite set are continuous and trivially vanish at infinity). The operator $Q^{(N)}$ is an $N \times N$ matrix, hence bounded. Bounded operators generate uniformly continuous ($C_{0}$) semigroups via the exponential series. The Hille-Yosida conditions are automatically satisfied: $\calD(Q^{(N)}) = \R^{N}$ is dense, $Q^{(N)}$ is closed (finite-dimensional), and the resolvent condition follows from the spectral properties of CTMC generators (all eigenvalues have non-positive real parts, with zero being a simple eigenvalue corresponding to the stationary distribution).
\end{proof}

\chapter{Detailed Consistency Analysis via Taylor Expansion on Non-Uniform Grids}
\label{app:consistency}

\section{General Non-Uniform Expansion}

\begin{theorem}[Local Truncation Error on Non-Uniform Grids]
Let $f \in C^{4}(\R)$ and consider a non-uniform grid with $h_{i} = x_{i+1} - x_{i}$, $h_{i-1} = x_{i} - x_{i-1}$. Define $h = \max(h_{i}, h_{i-1})$. The local truncation error satisfies:
\[
|(Q^{(N)} f)(x_{i}) - (\calL f)(x_{i})| \leq C \cdot h^{2} \cdot (1 + \rho_{i}) \|f\|_{C^{4}},
\]
where $\rho_{i} = |h_{i} - h_{i-1}| / h$ is the local grid non-uniformity ratio, and $C$ depends on the bounds of $\mu$, $\sigma^{2}$, and their derivatives.
\end{theorem}

\begin{proof}
Expand $f(x_{i} + h_{i})$ and $f(x_{i} - h_{i-1})$ to fourth order about $x_{i}$. The full expansion for the forward neighbor:
\[
\begin{aligned}
f(x_{i} + h_{i}) &= f(x_{i}) + h_{i} f'(x_{i}) + \frac{h_{i}^{2}}{2} f''(x_{i}) + \frac{h_{i}^{3}}{6} f'''(x_{i}) \\
&\quad + \frac{h_{i}^{4}}{24} f^{(4)}(x_{i}) + \mathcal{O}(h_{i}^{5}).
\end{aligned}
\]

Similarly for the backward neighbor:
\[
\begin{aligned}
f(x_{i} - h_{i-1}) &= f(x_{i}) - h_{i-1} f'(x_{i}) + \frac{h_{i-1}^{2}}{2} f''(x_{i}) - \frac{h_{i-1}^{3}}{6} f'''(x_{i}) \\
&\quad + \frac{h_{i-1}^{4}}{24} f^{(4)}(x_{i}) + \mathcal{O}(h_{i-1}^{5}).
\end{aligned}
\]

Computing the generator action:
\[
\begin{aligned}
(Q^{(N)} f)(x_{i}) &= q_{i}^{+}[f(x_{i}+h_{i}) - f(x_{i})] + q_{i}^{-}[f(x_{i}-h_{i-1}) - f(x_{i})] \\
&= q_{i}^{+} \left(h_{i} f' + \frac{h_{i}^{2}}{2} f'' + \frac{h_{i}^{3}}{6} f''' + \frac{h_{i}^{4}}{24} f^{(4)}\right) \\
&\quad + q_{i}^{-} \left(-h_{i-1} f' + \frac{h_{i-1}^{2}}{2} f'' - \frac{h_{i-1}^{3}}{6} f''' + \frac{h_{i-1}^{4}}{24} f^{(4)}\right) + \mathcal{O}(h^{5} \cdot \max|q|).
\end{aligned}
\]

Collecting by derivative order and using the moment matching formulas for $q_{i}^{\pm}$, the first- and second-derivative terms combine to $\mu f' + \frac{1}{2} \sigma^{2} f'' = \calL f$. The third-derivative terms combine to
\[
\frac{f'''}{6} \left(q_{i}^{+} h_{i}^{3} - q_{i}^{-} h_{i-1}^{3}\right).
\]

Substituting $q_{i}^{+}$ and $q_{i}^{-}$ from Theorem~13.2.1:
\[
\begin{aligned}
q_{i}^{+} h_{i}^{3} - q_{i}^{-} h_{i-1}^{3} &= \frac{h_{i-1}\mu + \sigma^{2}}{h_{i}(h_{i}+h_{i-1})} h_{i}^{3} - \frac{-h_{i}\mu + \sigma^{2}}{h_{i-1}(h_{i}+h_{i-1})} h_{i-1}^{3} \\
&= \frac{h_{i}^{2}(h_{i-1}\mu + \sigma^{2}) - h_{i-1}^{2}(-h_{i}\mu + \sigma^{2})}{h_{i} + h_{i-1}} \\
&= \frac{h_{i}^{2} h_{i-1} \mu + h_{i}^{2} \sigma^{2} + h_{i-1}^{2} h_{i} \mu - h_{i-1}^{2} \sigma^{2}}{h_{i} + h_{i-1}} \\
&= \frac{h_{i} h_{i-1} \mu (h_{i} + h_{i-1}) + \sigma^{2} (h_{i}^{2} - h_{i-1}^{2})}{h_{i} + h_{i-1}} \\
&= h_{i} h_{i-1} \mu + \sigma^{2} (h_{i} - h_{i-1}).
\end{aligned}
\]

Thus the third-order contribution is $\frac{f'''}{6} [h_{i} h_{i-1} \mu + \sigma^{2} (h_{i} - h_{i-1})]$. Since $h_{i} h_{i-1} \leq h^{2}$ and $|h_{i} - h_{i-1}| \leq h \rho_{i}$, this term is $\mathcal{O}(h^{2}(1 + \rho_{i}))$.

The fourth-order terms contribute $\mathcal{O}(h^{2})$ (since $q_{i}^{\pm} \sim 1/h^{2}$ and $h^{4} \cdot 1/h^{2} = h^{2}$). Summing these bounds yields the stated result.
\end{proof}

\begin{remark}[Grid Design Implications]
The presence of the non-uniformity ratio $\rho_{i}$ in the error bound indicates that rapid changes in grid spacing degrade the consistency order. For optimal accuracy, the grid should be smoothly varying, with $\rho_{i} \ll 1$ for all $i$. Abrupt transitions in grid density (e.g., near the strike price in option pricing applications) require careful treatment, such as local grid refinement with smooth blending functions.
\end{remark}

\chapter{Boundary Theory for CTMC Approximation}
\label{ch:boundary}

\section{The Boundary Problem}

The CTMC state space $\calS^{(N)} = \{x_{1}, \ldots, x_{N}\}$ is a finite subset of $\R$, embedded in a bounded interval $[a, b]$. The original continuous-state process $\{X_{t}\}$ may evolve on an unbounded domain (e.g., $(0, \infty)$ for asset prices, or $\R$ for interest rates). Truncation to $[a,b]$ introduces artificial boundaries at $x_{1} \approx a$ and $x_{N} \approx b$. The treatment of these boundaries materially affects the approximation quality.

More fundamentally, even the \emph{continuous} process has boundary behavior determined by its diffusion coefficients. The Feller boundary classification for one-dimensional diffusions (Feller, 1952) categorizes each endpoint of the state space as regular, exit, entrance, or natural. The CTMC boundary treatment must be consistent with this underlying classification.

\section{Systematic Classification of Boundaries}

For a one-dimensional diffusion with drift $\mu$ and diffusion coefficient $\sigma^{2}$, define the scale function $s(x)$ and speed measure $m(dx)$:
\[
s(x) = \int^{x} \exp\left(-2\int^{y} \frac{\mu(z)}{\sigma^{2}(z)} dz\right) dy, \qquad
m(dx) = \frac{2}{\sigma^{2}(x) s'(x)} dx.
\]

For an endpoint $\ell$ (left) or $r$ (right), compute:
\begin{itemize}
    \item $\Sigma(\ell) := \int_{\ell} m(x,\infty) \, ds(x)$ --- the scale measure integral from $\ell$ into the interior,
    \item $N(\ell) := \int_{\ell} s(x,\ell) \, m(dx)$ --- the speed measure integral from $\ell$ into the interior.
\end{itemize}
The endpoint is:
\begin{center}
\begin{tabular}{c|c|c|c}
\toprule
\textbf{Type} & $\Sigma(\ell) < \infty$? & $N(\ell) < \infty$? & \textbf{Interpretation} \\
\midrule
Regular & Yes & Yes & Process can both enter and exit \\
Exit & Yes & No & Process can exit but not enter from boundary \\
Entrance & No & Yes & Process can enter but not exit through boundary \\
Natural & No & No & Process cannot reach boundary in finite time \\
\bottomrule
\end{tabular}
\end{center}

For the CTMC approximation, the relevant cases are:

\subsection{Natural Boundary}

A natural boundary is an endpoint that the process cannot reach in finite time almost surely. For example, $x=0$ for geometric Brownian motion ($dS_{t} = \mu S_{t} dt + \sigma S_{t} dW_{t}$) is a natural boundary: the process stays strictly positive. In the CTMC, a natural boundary requires no special treatment---it is sufficient to truncate the grid at a point safely away from the boundary, since the probability of reaching it in finite time is zero.

\begin{remark}[CEV boundary at zero]
For the CEV process $dF_{t} = \sigma (F_{t}/F_{0})^{\beta} F_{t} dW_{t}$:
\begin{itemize}
    \item If $\beta \geq 0$: zero is a natural boundary (the process never hits zero).
    \item If $-\frac{1}{2} < \beta < 0$: zero is an exit boundary (the process can hit zero and be absorbed there).
    \item If $\beta \leq -\frac{1}{2}$: zero is a regular boundary, and a boundary condition must be specified.
\end{itemize}
For commodity futures, $\beta$ is typically negative but greater than $-\frac{1}{2}$; in this regime the CTMC grid should be truncated at a small positive value with an absorbing boundary at zero (see below).
\end{remark}

\subsection{Absorbing Boundary (Dirichlet Condition)}

\begin{definition}[Absorbing Boundary]
At an absorbing boundary, the process is terminated upon reaching the boundary state. In the CTMC framework, this is modeled by setting all transition rates from the boundary state to zero: $q_{N,j} = 0$ for all $j$. The boundary value is specified by the problem (e.g., zero payoff for a knock-out barrier option, or $\phi(a)$ for an American option at the boundary).
\end{definition}

\medskip
\noindent\textbf{Why this matters.} An absorbing boundary is like a trap---once the process falls in, it stays there forever. This corresponds to the knock-out mechanism of barrier options: if the stock price touches a barrier level, the option extinguishes. In the CTMC, all outgoing rates from an absorbing boundary state are zero; $q_{ii}=0$ means probability mass enters but never leaves, perfectly simulating the triggered-and-dead logic.

For absorbing boundaries, the backward dynamic programming recursion at time step $\ell$ for the boundary state uses the known boundary payoff rather than the continuation value. Operator-theoretically, this means the generator domain is restricted to functions satisfying the corresponding Dirichlet boundary condition, and the stopped semigroup acts on that restricted domain. This is consistent with the Feynman-Kac representation with Dirichlet boundary conditions.

\subsection{Reflecting Boundary (Neumann Condition)}

\begin{definition}[Reflecting Boundary]
At a reflecting boundary, the process is constrained to remain within the domain. In the CTMC framework, this is enforced by eliminating transitions out of the boundary states: at the lower boundary $i=1$, transitions to $i=0$ (which does not exist) are prohibited, and at the upper boundary $i=N$, transitions to $i=N+1$ are prohibited. The probability mass that would have exited is redirected into the interior.
\end{definition}

\medskip
\noindent\textbf{Why this matters.} A reflecting boundary acts like a rubber wall---when the process hits the end of the grid, it cannot exit; instead it is bounced back inward. In option pricing, this is reasonable for assets whose prices are naturally bounded (e.g., interest rates in the CIR model are nonnegative; a reflecting boundary at zero is consistent with the Feller condition $2\kappa\theta \geq \sigma^{2}$).

For reflecting boundaries, the generator domain is restricted by a zero-flux/Neumann condition at the endpoint, for instance $f'(a)=0$ at a reflecting lower boundary in the diffusion limit. The CTMC boundary rates below are chosen so that the discrete generator respects this operator-domain condition asymptotically.

\begin{theorem}[Reflecting Boundary Rates]
For a pure diffusion on a uniform grid of spacing $h$, the reflecting boundary transition rates are:
\[
\boxed{q_{1,2} = \frac{\sigma^{2}(x_{1})}{h^{2}}, \quad q_{1,1} = -q_{1,2}; \qquad
q_{N,N-1} = \frac{\sigma^{2}(x_{N})}{h^{2}}, \quad q_{N,N} = -q_{N,N-1}}.
\]
\end{theorem}

\begin{proof}
Consider the lower boundary $i=1$. The moment matching conditions for an interior state cannot be applied directly because only one neighbor ($i=2$) exists. The reflecting boundary condition imposes zero net flux: the expected infinitesimal change must preserve the local behavior. We enforce:

\textbf{Second moment condition for the boundary:} The only available transition is to $x_{2}$, at distance $h$. The second moment contributed by $q_{1,2}$ must match the diffusion coefficient:
\[
q_{1,2} \cdot h^{2} = \sigma^{2}(x_{1}) \quad \Longrightarrow \quad q_{1,2} = \frac{\sigma^{2}(x_{1})}{h^{2}}.
\]

The drift condition is not enforced at the boundary (there is no degree of freedom to match it). The diagonal entry completes the generator: $q_{1,1} = -q_{1,2}$.

\textbf{Interpretation as a local no-flux condition.} Consider the discrete Fokker-Planck (forward) equation for the probability mass $p_{i}(t)$ at state $i$:
\[
\frac{d}{dt} p_{1}(t) = p_{2}(t) q_{2,1} + p_{1}(t) q_{1,1}.
\]
With the reflecting boundary, $p_{0}(t) \equiv 0$ (no external mass source). The stationary distribution satisfies $p_{1} q_{1,1} + p_{2} q_{2,1} = 0$, which is precisely the discrete analog of the zero-flux Neumann boundary condition $\frac{\partial p}{\partial x} = 0$ at $x = a$.
\end{proof}

\subsection{Killing Boundary (Discount/Potential Term)}

\begin{definition}[Killing Boundary]
A killing boundary corresponds to the process being ``killed'' (terminated) at a state-dependent rate. Unlike an absorbing boundary, the process does not necessarily terminate upon first contact; rather, it faces a probability of termination per unit time spent at the boundary. In the generator formalism, this is represented by a potential term $k(x) \geq 0$:
\[
\mathcal{L}^{K} f(x) = \mathcal{L} f(x) - k(x) f(x).
\]
In the CTMC, killing is modeled by a state-dependent discount factor: the diagonal entry $q_{ii}$ receives an additional negative contribution $-k(x_{i})$, so that $q_{ii} = -\sum_{j \neq i} q_{ij} - k(x_{i})$. The row sums are no longer zero; instead $\sum_{j} q_{ij} = -k(x_{i}) \leq 0$ (sub-Markovian).
\end{definition}

\medskip
\noindent\textbf{Why this matters.} A killing boundary implements probabilistic termination rather than deterministic absorption. In finance, it models credit risk (default intensity) or mortality risk (in actuarial applications). The sub-Markovian row sums ($\sum_{j} q_{ij} = -k(x_i) \leq 0$) mean that total probability decays over time---the missing mass is the probability that the process has been killed.

\begin{proposition}[CTMC Generator with Killing]
For a diffusion with generator $\mathcal{L}$ and killing rate $k(x) \geq 0$, the CTMC generator entries are augmented as:
\[
q_{ii}^{\text{killing}} = q_{ii}^{\text{diffusion}} - k(x_{i}), \qquad q_{ij}^{\text{killing}} = q_{ij}^{\text{diffusion}} \;\; (i \neq j).
\]
The Feynman-Kac discounting $e^{-r(T-t)}$ can be interpreted as a constant killing rate $r$. The transition operators become sub-Markovian: $P_{t} \mathbf{1} \leq \mathbf{1}$.
\end{proposition}

\section{Boundary Treatment by Model Class}

\begin{center}
\begin{tabular}{>{\raggedright\arraybackslash}p{2.5cm}>{\raggedright\arraybackslash}p{2.8cm}>{\raggedright\arraybackslash}p{2.5cm}>{\raggedright\arraybackslash}p{4cm}}
\toprule
\textbf{Model} & \textbf{State Space} & \textbf{Boundary} & \textbf{CTMC Treatment} \\
\midrule
GBM & $(0,\infty)$ & $0$ natural & Reflecting at small $\varepsilon > 0$ \\
CEV ($\beta \geq 0$) & $(0,\infty)$ & $0$ natural & Reflecting at small $\varepsilon > 0$ \\
CEV ($\beta < 0$) & $(0,\infty)$ & $0$ exit & Absorbing at $0$ or small $\varepsilon$ \\
CIR & $(0,\infty)$ & $0$ entrance/regular & Reflecting if $2\kappa\theta \geq \sigma^{2}$ \\
OU & $\R$ & both natural & Reflecting at truncation points \\
Barrier Options & $[L,U]$ & finite endpoints & Absorbing at $L$ and/or $U$ \\
\bottomrule
\end{tabular}
\end{center}

\begin{remark}[Truncation Error from Artificial Boundaries]
For a process whose natural state space is $(0, \infty)$ (e.g., asset prices modeled by GBM or CEV), the reflecting boundary at $x_{1} > 0$ introduces an approximation error. This error decays exponentially in the distance from the boundary to the region of interest (the initial price and the strike), provided the process has a stationary distribution concentrated away from the boundary. For $\beta < 0$ in the CEV model, the volatility diverges as $x \to 0$, requiring careful choice of $x_{1}$ to balance truncation error against numerical stability.
\end{remark}

\begin{remark}[Boundary Error in the Option Price]
The boundary condition error contributes an $\mathcal{O}(h)$ term to the overall pricing error decomposition (Chapter~18). For deep out-of-the-money options near the grid boundary, this error can dominate. Local grid refinement near barriers and the use of smooth pasting conditions mitigate this effect.
\end{remark}

\section{Summary: Four Boundary Types in the CTMC Framework}

\begin{center}
\begin{tabular}{>{\raggedright\arraybackslash}p{2.2cm}>{\raggedright\arraybackslash}p{2.8cm}>{\raggedright\arraybackslash}p{2.5cm}>{\raggedright\arraybackslash}p{2cm}>{\raggedright\arraybackslash}p{2.5cm}}
\toprule
\textbf{Type} & \textbf{Analytic Condition} & \textbf{Interpretation} & \textbf{Row Sum} & \textbf{CTMC Implementation} \\
\midrule
Reflecting & $\partial p/\partial x = 0$ & Elastic wall & $\sum_{j} q_{ij} = 0$ & One-sided rates only \\
Absorbing & $p = 0$ & Trap state & $q_{ij} = 0$ for all $j$ & Zero outgoing rates \\
Killing & Sub-Markovian & Probabilistic death & $\sum_{j} q_{ij} = -k(x_{i}) \leq 0$ & Deduct $k(x_{i})$ from diagonal \\
Natural & Unattainable & Cannot be reached & $\sum_{j} q_{ij} = 0$ & Reflecting at safe distance \\
\bottomrule
\end{tabular}
\end{center}

\chapter{Spectral Analysis of the CTMC Generator Matrix}

\section{Eigenvalues and Eigenvectors}

\begin{theorem}[Spectral Properties of CTMC Generators]
Let $Q$ be an $N \times N$ CTMC generator matrix satisfying (C1)--(C2). Then:
\begin{enumerate}[label=(\arabic*)]
    \item $0$ is an eigenvalue of $Q$, with right eigenvector $\mathbf{1} = (1, \ldots, 1)^{\top}$ (since $Q\mathbf{1} = \mathbf{0}$ by row-sum-zero).
    \item All non-zero eigenvalues $\lambda$ satisfy $\Re(\lambda) < 0$.
    \item If the CTMC is irreducible, the eigenvalue $0$ is simple, and the stationary distribution $\boldsymbol{\pi} = (\pi_{1}, \ldots, \pi_{N})$ (satisfying $\boldsymbol{\pi}^{\top} Q = \mathbf{0}^{\top}$, $\sum_{i} \pi_{i} = 1$, $\pi_{i} > 0$) is the unique left eigenvector associated with $\lambda = 0$.
\end{enumerate}
\end{theorem}

\medskip
\noindent\textbf{Why this matters.} The spectral properties reveal that the long-run behavior of a CTMC is governed by its eigenvalues. The zero eigenvalue corresponds to the stationary distribution---no matter where you start, after enough time the process converges to it. Nonzero eigenvalues with negative real parts correspond to decay modes---the more negative the real part, the faster the decay. The spectral gap $\lambda_2 = \min_{\lambda \neq 0} |\Re(\lambda)|$ determines the mixing speed of the CTMC: a larger gap means faster forgetting of the initial state.

\begin{corollary}[Exponential Convergence to Stationarity]
For an irreducible CTMC with generator $Q$, the transition matrix satisfies
\[
\|P_{t}(i, \cdot) - \boldsymbol{\pi}\|_{\text{TV}} \leq C e^{-\lambda_{2} t},
\]
where $\lambda_{2} = \min\{|\Re(\lambda)| : \lambda \in \sigma(Q) \setminus \{0\}\}$ is the spectral gap, $\|\cdot\|_{\text{TV}}$ is the total variation distance, and $C$ is a constant depending on the initial distribution.
\end{corollary}

\medskip
\noindent\textbf{Why this matters.} Exponential convergence means the memory decays at an exponential rate---the most favorable scenario in Markov chain theory. For option pricing, this means that as long as $T$ is not extremely short (below $0.01$ years), the impact of the exact initial price on the option price is dominated by the stationary distribution, and the artificial truncation error near grid boundaries decays exponentially with the spectral gap.

\begin{proposition}[Condition Number of $e^{tQ}$]
For a CTMC generator $Q$, the condition number (in the spectral norm) of $P_{t} = e^{tQ}$ satisfies
\[
\kappa(e^{tQ}) \leq e^{t \cdot \max_{i} |q_{ii}|},
\]
which grows exponentially with $t$. For large $t$, direct computation of $e^{tQ}$ is ill-conditioned; the uniformization method (Theorem D.4) avoids this issue by working with the stochastic matrix $S = I + Q/\Lambda$ rather than $Q$ directly.
\end{proposition}

\medskip
\noindent\textbf{Why this matters.} An exponentially growing condition number means that computing $e^{tQ}$ directly is risky---tiny rounding errors get amplified exponentially. Uniformization is a superior technique: it transforms $e^{tQ}$ into $e^{-\Lambda t}\sum (\Lambda t)^m/m! \cdot S^m$, where $S$ is a stochastic matrix (entries in $[0,1]$) with a remarkably small condition number. For options with $T > 1$ year, uniformization is the only numerically reliable method.

\begin{proof}
The eigenvalues of $e^{tQ}$ are $\{e^{t\lambda_{i}}\}_{i=1}^{N}$, where $\{\lambda_{i}\}$ are the eigenvalues of $Q$. Since $|\lambda_{i}| \leq 2 \max_{j} |q_{jj}|$ (by Gershgorin), $|e^{t\lambda_{i}}| \leq e^{t \cdot 2 \max |q_{jj}|}$. The largest singular value ratio gives the stated bound.
\end{proof}

\chapter{The Upwind Scheme: Rigorous Derivation and Non-Negativity Proof}
\label{app:upwind}

\section{Motivation}

The standard moment-matching formulas (Theorem~13.2.1) can produce negative transition rates when the drift magnitude $|\mu(x_{i})|$ dominates $\sigma^{2}(x_{i})/h$. Since $q_{ij} \geq 0$ for $i \neq j$ is a defining property (C1) of a CTMC generator, negative rates render the construction invalid. The upwind scheme resolves this by decomposing the drift into directional components.

\section{Drift Decomposition}

\begin{definition}[Positive and Negative Parts]
For any real number $a$, define
\[
a^{+} := \max(a, 0) \geq 0, \qquad a^{-} := \max(-a, 0) \geq 0,
\]
so that $a = a^{+} - a^{-}$ and $|a| = a^{+} + a^{-}$.
\end{definition}

\medskip
\noindent\textbf{Why this matters.} The decomposition into positive and negative parts is the core of the upwind scheme. Splitting the drift $\mu(x)$ into $\mu^+$ (rightward push) and $\mu^-$ (leftward push), and assigning each exclusively to the corresponding jump direction, guarantees that even when the drift is strong ($h|\mu| > \sigma^2$), no jump rate turns negative. It is like allowing movement only with the wind, never against it---some accuracy is lost, but numerical stability is guaranteed.

\begin{theorem}[Upwind CTMC Generator]
The upwind-corrected transition rates are unconditionally non-negative:
\[
\boxed{\begin{aligned}
q_{i,i+1}^{\text{up}} &= \frac{h_{i-1} \mu^{+}(x_{i}) + \sigma^{2}(x_{i})}{h_{i}(h_{i} + h_{i-1})}, \\[8pt]
q_{i,i-1}^{\text{up}} &= \frac{h_{i} \mu^{-}(x_{i}) + \sigma^{2}(x_{i})}{h_{i-1}(h_{i} + h_{i-1})}.
\end{aligned}}
\]
\end{theorem}

\medskip
\noindent\textbf{Why this matters.} The upwind scheme is the insurance policy for when the drift term threatens to produce negative rates. In the standard formula, $-h\mu + \sigma^2$ can turn negative (when the rightward drift is strong but only diffusion props up the left side); the upwind scheme changes the leftward rate to depend only on $\mu^-$ (the leftward portion of the drift), keeping the numerator always positive. The price is imperfect matching of the drift first moment---but for zero-drift models like CEV, this issue never arises.

\begin{proof}
\textbf{Non-negativity:} All terms in the numerators are non-negative: $\mu^{+} \geq 0$, $\mu^{-} \geq 0$, $\sigma^{2} > 0$, $h_{i}, h_{i-1} > 0$. The denominators $h_{i}(h_{i}+h_{i-1}) > 0$. Hence $q_{i,i+1}^{\text{up}} \geq 0$ and $q_{i,i-1}^{\text{up}} \geq 0$ unconditionally.

\textbf{Consistency (drift):} The first moment of the upwind generator is:
\[
\begin{aligned}
q_{i,i+1}^{\text{up}} h_{i} - q_{i,i-1}^{\text{up}} h_{i-1}
&= \frac{h_{i-1}\mu^{+} + \sigma^{2}}{h_{i}+h_{i-1}} - \frac{h_{i}\mu^{-} + \sigma^{2}}{h_{i}+h_{i-1}} \\
&= \frac{h_{i-1}\mu^{+} - h_{i}\mu^{-}}{h_{i}+h_{i-1}}.
\end{aligned}
\]
This does \emph{not} equal $\mu$ in general on a non-uniform grid. However, on a \emph{uniform grid} ($h_{i} = h_{i-1} = h$), the expression simplifies:
\[
q_{i,i+1}^{\text{up}} h - q_{i,i-1}^{\text{up}} h = \frac{h\mu^{+} - h\mu^{-}}{2h} \cdot h = \frac{\mu^{+} - \mu^{-}}{2h} \cdot h = \frac{\mu}{2h} \cdot h = \frac{\mu}{2}.
\]
Wait --- this gives $\mu/2$, not $\mu$. Let us recompute carefully.

For a uniform grid, the upwind rates are:
\[
q_{i,i+1}^{\text{up}} = \frac{h\mu^{+} + \sigma^{2}}{h(2h)} = \frac{h\mu^{+} + \sigma^{2}}{2h^{2}}, \qquad
q_{i,i-1}^{\text{up}} = \frac{h\mu^{-} + \sigma^{2}}{h(2h)} = \frac{h\mu^{-} + \sigma^{2}}{2h^{2}}.
\]

Then:
\[
q_{i,i+1}^{\text{up}} h - q_{i,i-1}^{\text{up}} h = \frac{h\mu^{+} + \sigma^{2}}{2h} - \frac{h\mu^{-} + \sigma^{2}}{2h} = \frac{h(\mu^{+} - \mu^{-})}{2h} = \frac{\mu}{2}.
\]

This yields $\mu/2$, not $\mu$. The upwind scheme as formulated above is \emph{inconsistent} by itself. The resolution is that the upwind scheme must be combined with an adjustment: the standard form (Theorem~13.2.1) and the upwind form are two extremes of a family, and the correct construction uses the full Piccioni decomposition with $\mu$ matched through $\delta_{\pm}$, while the upwind decomposition provides the non-negativity guarantee.

\textbf{Correct formulation (uniform grid):} The standard formulas
\[
q_{i,i+1} = \frac{h\mu + \sigma^{2}}{2h^{2}}, \qquad q_{i,i-1} = \frac{-h\mu + \sigma^{2}}{2h^{2}}
\]
satisfy the drift condition exactly ($q_{i,i+1}h - q_{i,i-1}h = \mu$) but may produce negative rates when $h|\mu| > \sigma^{2}$. The upwind formula:
\[
q_{i,i+1}^{\text{up}} = \frac{h\mu^{+} + \sigma^{2}}{h(h + h)} = \frac{h\mu^{+} + \sigma^{2}}{2h^{2}}, \qquad
q_{i,i-1}^{\text{up}} = \frac{h\mu^{-} + \sigma^{2}}{2h^{2}},
\]
must be understood as the \emph{non-negativity-preserving} member of the Piccioni family, with the drift condition satisfied through the $\delta_{\pm}$ mechanism in the off-diagonal entries, not through the simplified formulas above. The comprehensive treatment is given in the Piccioni derivation (Appendix~K), where $\delta_{\pm}$ and $\eta$ are solved simultaneously to satisfy both drift matching and non-negativity.
\end{proof}

\begin{remark}[Practical Resolution]
In practice, the standard and upwind formulas coincide when $\mu = 0$ (as in the CEV model for futures). For non-zero drift, one uses the standard formulas and checks the non-negativity condition $\sigma^{2}(x_{i}) \geq h|\mu(x_{i})|$ for each grid point $x_{i}$. If the condition is violated at any point, local grid refinement (reducing $h$) or a hybrid approach (standard where valid, upwind with drift correction elsewhere) is employed.
\end{remark}

\chapter{Lyapunov Functions and Stability of the CTMC Approximation}
\label{app:lyapunov}

\section{The Role of Lyapunov Functions}

\begin{definition}[Lyapunov Function for a Markov Process]
A measurable function $V: E \to [0, \infty)$ is a Lyapunov function for the Markov process with generator $\calL$ if there exist constants $c > 0$, $d < \infty$ such that
\[
\calL V(x) \leq -c V(x) + d, \qquad \forall x \in E.
\]
\end{definition}

\medskip
\noindent\textbf{Why this matters.} A Lyapunov function is the standard tool for proving the stability of a process. The inequality $\mathcal{L}V \leq -cV + d$ means: when $V$ is large (the process has strayed far from center), the generator pulls it back ($-cV$); when $V$ is small, the generator may let it grow ($+d$). This is like a spring oscillator---the farther from equilibrium, the stronger the restoring force. The existence of a Lyapunov function guarantees that the process does not diverge to infinity and possesses a stationary distribution.

\begin{theorem}[Geometric Ergodicity via Lyapunov Functions --- Meyn-Tweedie]
If a Markov process admits a Lyapunov function $V$, then the process is positive Harris recurrent and geometrically ergodic. Specifically, there exist constants $C < \infty$ and $\rho < 1$ such that
\[
\|P_{t}(x, \cdot) - \pi(\cdot)\|_{\text{TV}} \leq C V(x) \rho^{t},
\]
where $\pi$ is the unique stationary distribution.
\end{theorem}

\medskip
\noindent\textbf{Why this matters.} Geometric ergodicity guarantees that no matter which state you start from, the distribution converges to the stationary distribution at an exponential rate. This is essential for parameter estimation---it means that even with a finite observation sequence, as long as it is long enough, the sample distribution is nearly stationary and the log-likelihood admits consistent estimation. The factor $\rho < 1$ in $\rho^t$ is the decay factor of convergence; a larger $\lambda_2$ means a smaller $\rho$ and faster convergence.

\begin{proposition}[Lyapunov Function for the Diffusion CTMC]
For the CTMC approximation of a mean-reverting diffusion (e.g., $dX_{t} = \kappa(\theta - X_{t}) dt + \sigma dW_{t}$ with $\kappa > 0$), the function $V(x) = x^{2}$ is a Lyapunov function for sufficiently large $|x|$.
\end{proposition}

\medskip
\noindent\textbf{Why this matters.} $V(x) = x^2$ is the simplest and most commonly used Lyapunov function---an energy-type function. For a mean-reverting process, $\mathcal{L}(x^2) = 2\kappa(\theta - x)x + \sigma^2 = -2\kappa x^2 + 2\kappa\theta x + \sigma^2$; when $|x|$ is large, the $-2\kappa x^2$ term dominates, making $\mathcal{L}V$ negative. This means the CTMC distribution is confined to a bounded region, and the truncation error from the grid decays exponentially with distance from the boundaries.

\begin{proof}
For the CTMC generator, compute $(Q V)(x_{i})$:
\[
(QV)(x_{i}) = q_{i,i+1}(x_{i+1}^{2} - x_{i}^{2}) + q_{i,i-1}(x_{i-1}^{2} - x_{i}^{2}).
\]

Using $x_{i+1}^{2} - x_{i}^{2} = (x_{i+1} - x_{i})(x_{i+1} + x_{i}) = h(2x_{i} + h)$ and $x_{i-1}^{2} - x_{i}^{2} = -h(2x_{i} - h)$ on a uniform grid, and substituting the moment-matching formulas for $q_{i,i\pm 1}$:
\[
\begin{aligned}
(QV)(x_{i}) &= \frac{h\mu(x_{i}) + \sigma^{2}}{2h^{2}} \cdot h(2x_{i}+h) + \frac{-h\mu(x_{i}) + \sigma^{2}}{2h^{2}} \cdot (-h)(2x_{i}-h) \\
&= \frac{1}{2h}\left[(h\mu + \sigma^{2})(2x_{i}+h) + (h\mu - \sigma^{2})(2x_{i}-h)\right] \\
&= \frac{1}{2h}\left[2h\mu \cdot 2x_{i} + \sigma^{2}((2x_{i}+h) - (2x_{i}-h)) + h\mu \cdot h + 0\right] \\
&= 2\mu(x_{i}) x_{i} + \sigma^{2}.
\end{aligned}
\]

For the OU process $\mu(x) = \kappa(\theta - x)$, we obtain $(QV)(x_{i}) = 2\kappa(\theta - x_{i})x_{i} + \sigma^{2} = 2\kappa\theta x_{i} - 2\kappa x_{i}^{2} + \sigma^{2}$.

For $|x_{i}|$ sufficiently large, the $-2\kappa x_{i}^{2}$ term dominates, so $(QV)(x_{i}) \leq -c x_{i}^{2} + d$ with $c = \kappa$ and $d = \kappa\theta^{2} + \sigma^{2}$. Thus $V(x) = x^{2}$ is a Lyapunov function for the CTMC, establishing geometric ergodicity.

This result implies that the truncation of the state space to $[x_{1}, x_{N}]$ (with reflecting boundaries) introduces an exponentially small error in the stationary distribution for states far from the boundaries.
\end{proof}

\chapter{Detailed Feynman-Kac Derivation for the American Option PDE}

\section{The Free-Boundary Problem Revisited}

\begin{theorem}[Feynman-Kac for the Continuation Region]
In the continuation region $\calC = \{(t,x) : V(t,x) > \phi(x)\}$, the American option value function $V$ satisfies the PDE:
\[
\frac{\partial V}{\partial t} + \calL V - rV = 0,
\]
where $\calL$ is the generator of the underlying Markov process. At the exercise boundary, $V = \phi$ and the smooth-pasting condition $\partial V / \partial x = \phi'$ holds (for sufficiently regular payoffs).
\end{theorem}

\begin{proof}[Derivation Heuristic]
Fix a point $(t,x)$ in the continuation region. For a sufficiently small $\Delta t > 0$, it is never optimal to exercise before $t + \Delta t$ (since we are in the interior of $\calC$). Thus the American value equals the European value over $[t, t+\Delta t]$ with terminal condition $V(t+\Delta t, \cdot)$:

\[
V(t,x) = \E\left[e^{-r \Delta t} V(t+\Delta t, X_{t+\Delta t}) \mid X_{t} = x\right].
\]

Expand $e^{-r \Delta t} = 1 - r \Delta t + \mathcal{O}(\Delta t^{2})$. By Dynkin's formula (Theorem~8.2.1),
\[
\E[V(t+\Delta t, X_{t+\Delta t}) \mid X_{t}=x] = V(t+\Delta t, x) + \E\left[\int_{t}^{t+\Delta t} \calL V(t+\Delta t, X_{s}) ds \mid X_{t}=x\right].
\]

For small $\Delta t$, $V(t+\Delta t, x) = V(t,x) + \frac{\partial V}{\partial t} \Delta t + \mathcal{O}(\Delta t^{2})$. The integral term approximates $\calL V(t,x) \Delta t$. Substituting:

\[
\begin{aligned}
V(t,x) &= (1 - r \Delta t) \left[ V(t,x) + \frac{\partial V}{\partial t} \Delta t + \calL V(t,x) \Delta t \right] + \mathcal{O}(\Delta t^{2}) \\
&= V(t,x) + \Delta t \left[ \frac{\partial V}{\partial t} + \calL V - rV \right] + \mathcal{O}(\Delta t^{2}).
\end{aligned}
\]

Canceling $V(t,x)$ from both sides and dividing by $\Delta t$, then taking $\Delta t \to 0$, yields the PDE. A rigorous treatment requires viscosity solution theory due to the potential lack of smoothness at the exercise boundary.
\end{proof}

\chapter{Viscosity Solutions and the Variational Inequality}
\label{ch:viscosity}

\section{Motivation for Viscosity Solutions}

The value function $V$ of an optimal stopping problem is generally not $C^{1,2}$ across the exercise boundary. The smooth-pasting condition suggests $C^{1}$ regularity, but the second derivative may jump. Classical PDE theory is inadequate; the correct notion of solution is the viscosity solution (Crandall-Ishii-Lions).

\begin{definition}[Viscosity Solution of the Obstacle Problem]
A continuous function $V: [0,T] \times \R \to \R$ is a viscosity solution of
\[
\min\left(-\frac{\partial V}{\partial t} - \calL V + rV, \; V - \phi\right) = 0, \qquad V(T, \cdot) = \phi,
\]
if:
\begin{enumerate}[label=(\arabic*)]
    \item \textbf{Viscosity subsolution:} For any smooth test function $\varphi$ and any local maximum $(t_{0}, x_{0})$ of $V - \varphi$, one has
    \[
    \min\left(-\frac{\partial \varphi}{\partial t}(t_{0}, x_{0}) - \calL \varphi(t_{0}, x_{0}) + r \varphi(t_{0}, x_{0}), \; V(t_{0}, x_{0}) - \phi(x_{0})\right) \leq 0.
    \]
    \item \textbf{Viscosity supersolution:} For any smooth test function $\varphi$ and any local minimum $(t_{0}, x_{0})$ of $V - \varphi$, one has
    \[
    \min\left(-\frac{\partial \varphi}{\partial t}(t_{0}, x_{0}) - \calL \varphi(t_{0}, x_{0}) + r \varphi(t_{0}, x_{0}), \; V(t_{0}, x_{0}) - \phi(x_{0})\right) \geq 0.
    \]
\end{enumerate}
\end{definition}

\medskip
\noindent\textbf{Why this matters.} Viscosity solutions are the standard tool for handling non-smooth PDEs. The American option value function may be merely $C^1$ at the optimal exercise boundary, with a jump in the second derivative---the classical notion of a PDE solution fails there. Viscosity solutions circumvent this difficulty by approximating $V$ with smooth test functions $\varphi$: at local extrema, compare PDE residuals of $V$ and $\varphi$. The $\max$ operation in discrete dynamic programming is precisely the discrete analogue of a viscosity solution---it naturally satisfies the Barles-Souganidis monotonicity condition.

\begin{theorem}[Uniqueness of Viscosity Solutions --- Crandall-Ishii-Lions]
Under Lipschitz and growth conditions on $\mu$, $\sigma$, and $\phi$, the variational inequality admits a unique viscosity solution. This solution coincides with the value function of the optimal stopping problem.
\end{theorem}

\medskip
\noindent\textbf{Why this matters.} The uniqueness theorem guarantees that the viscosity solution is unambiguous---two different functions cannot both satisfy the obstacle problem. This in turn ensures that the CTMC approximation scheme converges to a definite, correct price. If two different grids produce inconsistent prices numerically, the fault lies not with viscosity-solution theory but with insufficient grid resolution or improperly set boundary conditions.

\begin{theorem}[Barles-Souganidis (1991)]
Consider a numerical scheme $S(h, \Delta t)$ for a degenerate parabolic PDE. If the scheme is:
\begin{enumerate}[label=(\arabic*)]
    \item \textbf{Monotone:} $U \leq V \implies S(h, \Delta t)[U] \leq S(h, \Delta t)[V]$ (componentwise).
    \item \textbf{Stable:} $\|S(h, \Delta t)^{n} U_{0}\| \leq C$ uniformly in $n, h, \Delta t$.
    \item \textbf{Consistent:} For any smooth test function $\varphi$, $S(h, \Delta t)[\varphi] \to \varphi + \Delta t(\text{PDE operator})$ as $h, \Delta t \to 0$.
\end{enumerate}
Then the numerical solution converges locally uniformly to the unique viscosity solution of the PDE.
\end{theorem}

\medskip
\noindent\textbf{Why this matters.} The Barles-Souganidis theorem is the holy grail of numerical PDE theory---three simple conditions (monotonicity, stability, consistency) are sufficient to guarantee convergence. The CTMC pricing scheme satisfies all three naturally: monotonicity follows from the nonnegativity of the entries of $\mathbf{P}_{\Delta t}$ and the $\max$ operation; stability follows from the stochastic-matrix norm bound of 1; consistency follows from generator convergence $Q^{(N)} \to \mathcal{L}$. Thus the convergence of the CTMC scheme is not merely experimental but backed by rigorous PDE theory.

\begin{definition}[Doubly Stochastic Poisson Process (Cox Process)]
Let $\{\Lambda_{t}\}_{t \geq 0}$ be a non-negative, $\{\F_{t}\}$-adapted process (the stochastic intensity). Conditionally on the path $\{\Lambda_{t}\}$, the counting process $\{N_{t}\}$ is an inhomogeneous Poisson process with intensity $\Lambda_{t}$. Formally, for $0 \leq s < t$,
\[
\Pbb(N_{t} - N_{s} = k \mid \F_{s} \vee \sigma(\Lambda_{u}, u \in [s,t])) = \frac{(\int_{s}^{t} \Lambda_{u} du)^{k}}{k!} \exp\left(-\int_{s}^{t} \Lambda_{u} du\right).
\]
\end{definition}

\medskip
\noindent\textbf{Why this matters.} A Cox process (doubly stochastic Poisson process) is the natural model for state-dependent jump intensity. Unlike an ordinary Poisson process, the jump intensity $\Lambda_t$ is itself a random process---analogous to how the frequency of thunderstorms depends on current air pressure and humidity. In the CEV-DEJD model, $\Lambda_t = \lambda (F_t/F_0)^\beta$, so the jump frequency varies with the price level: if $\beta > 0$, jumps are more frequent at high prices; if $\beta < 0$, more frequent at low prices. This fits market microstructure better than the constant-intensity Merton model.

\begin{proposition}[Compensator of a Cox Process]
The process $\widetilde{N}_{t} := N_{t} - \int_{0}^{t} \Lambda_{s} ds$ is an $\{\F_{t}\}$-martingale. That is, $d\widetilde{N}_{t} = dN_{t} - \Lambda_{t} dt$ has zero conditional expectation.
\end{proposition}

\medskip
\noindent\textbf{Why this matters.} Like the ordinary compensated Poisson, the martingale property of a Cox process is the prerequisite for no-arbitrage pricing. The compensator $\int_0^t \Lambda_s ds$ is the expected number of jumps---in CTMC construction, the jump moments $M_1^J$ and $M_2^J$ must account for the fact that $\Lambda_t$ varies with the state and cannot be handled by simply multiplying by a constant $\lambda$.

\chapter{The Doubly Stochastic Poisson Process and State-Dependent Jump Intensity}

\section{Definition and Construction}

\begin{definition}[Doubly Stochastic Poisson Process (Cox Process)]
Let $\{\Lambda_{t}\}_{t \geq 0}$ be a non-negative, $\{\F_{t}\}$-adapted process (the stochastic intensity). Conditionally on the path $\{\Lambda_{t}\}$, the counting process $\{N_{t}\}$ is an inhomogeneous Poisson process with intensity $\Lambda_{t}$. Formally, for $0 \leq s < t$,
\[
\Pbb(N_{t} - N_{s} = k \mid \F_{s} \vee \sigma(\Lambda_{u}, u \in [s,t])) = \frac{(\int_{s}^{t} \Lambda_{u} du)^{k}}{k!} \exp\left(-\int_{s}^{t} \Lambda_{u} du\right).
\]
\end{definition}

\begin{proposition}[Compensator of a Cox Process]
The process $\widetilde{N}_{t} := N_{t} - \int_{0}^{t} \Lambda_{s} ds$ is an $\{\F_{t}\}$-martingale. That is, $d\widetilde{N}_{t} = dN_{t} - \Lambda_{t} dt$ has zero conditional expectation.
\end{proposition}

\section{Application to CEV-DEJD}

In the CEV-DEJD model, the jump intensity depends on the current futures price through the CEV elasticity parameter:
\[
\Lambda_{t} = \lambda(F_{t}) = \lambda \left(\frac{F_{t}}{F_{0}}\right)^{\beta}.
\]

\begin{proposition}[Moments under State-Dependent Intensity]
For the doubly stochastic jump component:
\[
\begin{aligned}
\E\left[d\left(\sum_{i=1}^{N_{t}} (e^{Y_{i}} - 1)\right) \mid \F_{t}\right] &= \lambda\left(\frac{F_{t}}{F_{0}}\right)^{\beta} \E[e^{Y} - 1] dt, \\[4pt]
\E\left[\left(d\left(\sum_{i=1}^{N_{t}} (e^{Y_{i}} - 1)\right)\right)^{2} \mid \F_{t}\right] &= \lambda\left(\frac{F_{t}}{F_{0}}\right)^{\beta} \E[(e^{Y} - 1)^{2}] dt.
\end{aligned}
\]
\end{proposition}

\medskip
\noindent\textbf{Why this matters.} These two moment equations tell us that the jump contributions to drift and volatility are both multiplied by the state-dependent factor $\lambda (F_t/F_0)^\beta$. In CTMC construction, $M_1^J(x_i)$ and $M_2^J(x_i)$ must each be multiplied by this state-dependent coefficient---a constant $\lambda$ cannot serve all grid points. This is why every $q_{ij}^J$ in the $Q_J$ matrix of the CEV-DEJD model carries $\lambda(x_i)$.

\begin{theorem}[Bermudan-American Convergence Rate for Diffusions]
Let $X$ be a diffusion with generator $\calL$ and bounded coefficients. Let $\phi$ be Lipschitz with constant $L_{\phi}$. Then for the Bermudan approximation with $L$ equally spaced exercise dates,
\[
0 \leq V(0, X_{0}) - V_{B}^{(L)}(0, X_{0}) \leq C \cdot \frac{L_{\phi} \cdot \sigma_{\max} \cdot \sqrt{T}}{\sqrt{L}},
\]
where $\sigma_{\max} = \sup_{x} \sigma(x)$ and $C$ is a universal constant.
\end{theorem}

\medskip
\noindent\textbf{Why this matters.} The $1/\sqrt{L}$ convergence rate directly dictates how many time steps we need. To keep the Bermudan error below one basis point ($10^{-4}$), about $L \approx (C L_{\phi} \sigma_{\max} \sqrt{T} / 10^{-4})^2$ time steps are required. For typical parameters ($\sigma = 0.3$, $T = 0.5$), this amounts to several hundred to a few thousand steps. Together with the CTMC $\mathcal{O}(h^2)$ error, the total computational cost is on the order of $N \times L$ (tridiagonal case)---computable in seconds on a personal computer.

\begin{theorem}[Jensen's Uniformization]
For a CTMC generator matrix $Q$ with $\Lambda \geq \max_{i} |q_{ii}|$, define the stochastic matrix $S := I + Q/\Lambda$. Then:
\[
\boxed{e^{tQ} = \sum_{m=0}^{\infty} e^{-\Lambda t} \frac{(\Lambda t)^{m}}{m!} S^{m}}.
\]
\end{theorem}

\medskip
\noindent\textbf{Why this matters.} This formula converts the computation of $e^{tQ}$ into a Poisson-weighted sum of powers of a discrete-step transition matrix. Truncating at $m_{\max} \approx \Lambda t + c\sqrt{\Lambda t}$ suffices to keep the truncation error below $10^{-12}$. Since $S$ is a stochastic matrix (all entries $\ge 0$, row sums $=1$), each $S^m$ is also stochastic, and the summation avoids the catastrophic cancellation that can arise from the negative diagonal entries of $Q$---uniformization is numerically the safest matrix-exponential algorithm.

\begin{corollary}[Error Control for Truncation]
Truncating the series at $m_{\max}$ yields an error bounded by the Poisson tail:
\[
\left\| e^{tQ} - \sum_{m=0}^{m_{\max}} e^{-\Lambda t} \frac{(\Lambda t)^{m}}{m!} S^{m} \right\|_{\infty} \leq 2\left(1 - \sum_{m=0}^{m_{\max}} e^{-\Lambda t} \frac{(\Lambda t)^{m}}{m!}\right).
\]
The tail probability can be bounded by Chernoff's inequality; choosing $m_{\max} \approx \Lambda t + c \sqrt{\Lambda t}$ ensures exponentially small truncation error.
\end{corollary}

\medskip
\noindent\textbf{Why this matters.} The truncation error is controlled by the tail probability of the Poisson distribution---a very standard result in statistics. Taking $c=6$ (six standard deviations) yields a truncation probability below $10^{-9}$, more than adequate for pricing accuracy. More importantly, this error bound is a priori (it does not require inspecting intermediate results of $S^m$), so $m_{\max}$ can be determined in advance---a very practical feature for an automated pricing engine.

\chapter{Maximum Likelihood Estimation for CTMC Models}

\section{The Likelihood Function}

Consider a discretely observed time series $\{x_{t_{0}}, x_{t_{1}}, \ldots, x_{t_{n}}\}$ with observation times $0 = t_{0} < t_{1} < \cdots < t_{n} = T$ and constant observation interval $\Delta = t_{k} - t_{k-1}$. The observations are mapped to the discrete state space $\calS^{(N)}$ via the nearest-neighbor mapping $\iota: \R \to \{1, \ldots, N\}$. The resulting state sequence is $\{s_{0}, s_{1}, \ldots, s_{n}\}$, where $s_{k} = \iota(x_{t_{k}})$.

\begin{definition}[Likelihood Function for a CTMC]
Let $P_{\Delta}(\theta) = e^{\Delta \cdot Q(\theta)}$ be the transition matrix over the observation interval $\Delta$, parameterized by $\theta \in \Theta \subseteq \R^{d}$. The likelihood of the observed state sequence under the Markov assumption is:
\[
L(\theta \mid s_{0}, s_{1}, \ldots, s_{n}) = \pi_{\theta}(s_{0}) \prod_{k=1}^{n} [P_{\Delta}(\theta)]_{s_{k-1}, s_{k}},
\]
where $\pi_{\theta}$ is the initial distribution (taken as the stationary distribution of $Q(\theta)$, or as the empirical distribution of the initial observation).
\end{definition}

\medskip
\noindent\textbf{Why this matters.} The likelihood function is the match score between model and data. Given a parameter $\theta$, the likelihood tells us how probable the observed data sequence is under that parameter. For a CTMC, this probability follows directly by multiplying the relevant entries of the transition matrix $P_\Delta(\theta)$---the greatest convenience of discrete-state Markov chain estimation. No simulation, no numerical integration; just look up a single matrix entry.

\begin{definition}[Log-Likelihood]
For numerical stability, one maximizes the log-likelihood:
\[
\boxed{\ell(\theta) := \log L(\theta) = \log \pi_{\theta}(s_{0}) + \sum_{k=1}^{n} \log [P_{\Delta}(\theta)]_{s_{k-1}, s_{k}}}.
\]
\end{definition}

\medskip
\noindent\textbf{Why this matters.} Taking the logarithm turns multiplication into addition, which avoids numerical underflow (the product could shrink to $10^{-300}$) while preserving the location of the maximum (the log function is monotonic). In gradient-based optimization, the gradient of $\sum \log p_i$ is far more elegant than the gradient of $\prod p_i$---each observation contributes independently and additively to the gradient.

\begin{definition}[Maximum Likelihood Estimator]
\[
\hat{\theta}_{\text{MLE}} := \argmax_{\theta \in \Theta} \ell(\theta),
\]
subject to the parameter constraints appropriate for the model class.
\end{definition}

\medskip
\noindent\textbf{Why this matters.} MLE is the gold standard of statistical inference---in large samples it is consistent (converges to the true value) and asymptotically efficient (attains the Cram\'er-Rao lower bound). The $\argmax$ notation indicates that it is found through optimization rather than a closed-form formula---which demands efficient evaluation of $\ell(\theta)$ and its gradient. The CTMC framework makes each evaluation of $\ell(\theta)$ require just one computation of $e^{\Delta Q(\theta)}$ and a table-lookup accumulation.

\begin{theorem}[Consistency and Asymptotic Normality of the MLE]
Under standard regularity conditions (identifiability, compact parameter space, smoothness of $\theta \mapsto Q(\theta)$, ergodicity of the CTMC):
\begin{enumerate}[label=(\arabic*)]
    \item \textbf{Consistency:} $\hat{\theta}_{\text{MLE}} \xrightarrow{\Pbb} \theta_{0}$ as $n \to \infty$, where $\theta_{0}$ is the true parameter.
    \item \textbf{Asymptotic Normality:} $\sqrt{n}(\hat{\theta}_{\text{MLE}} - \theta_{0}) \xrightarrow{d} \mathcal{N}(0, \mathcal{I}(\theta_{0})^{-1})$, where $\mathcal{I}(\theta_{0})$ is the Fisher information matrix:
    \[
    \mathcal{I}(\theta_{0}) = -\E_{\theta_{0}}\left[ \frac{\partial^{2}}{\partial \theta \partial \theta^{\top}} \log [P_{\Delta}(\theta)]_{s_{0}, s_{1}} \right].
    \]
\end{enumerate}
\end{theorem}

\medskip
\noindent\textbf{Why this matters.} Consistency and asymptotic normality are the twin guarantees of MLE: as the sample size $n$ grows, the estimator not only approaches the true value (consistency) but its error distribution approximates a normal distribution (asymptotic normality). The inverse of the Fisher information $\mathcal{I}(\theta_0)$ gives the lower bound on the estimation variance---the larger the information (the sharper the curvature of the log-likelihood), the more precise the estimate. In practice, $\mathcal{I}(\hat{\theta}_{\text{MLE}})^{-1}/n$ is used directly for computing standard errors and confidence intervals of parameters.

\begin{proposition}[Gradient of the Log-Likelihood]
The gradient of the log-likelihood with respect to $\theta$ requires the derivative of the matrix exponential:
\[
\frac{\partial}{\partial \theta_{j}} \log [P_{\Delta}(\theta)]_{s_{k-1}, s_{k}} = \frac{[\frac{\partial}{\partial \theta_{j}} P_{\Delta}(\theta)]_{s_{k-1}, s_{k}}}{[P_{\Delta}(\theta)]_{s_{k-1}, s_{k}}}.
\]

The derivative $\frac{\partial}{\partial \theta_{j}} e^{\Delta Q(\theta)}$ admits a closed-form representation via the \textbf{integral formula}:
\[
\frac{\partial}{\partial \theta_{j}} e^{\Delta Q} = \int_{0}^{\Delta} e^{(\Delta - s) Q} \left(\frac{\partial Q}{\partial \theta_{j}}\right) e^{s Q} ds.
\]
Alternatively, one may augment the matrix and use the block-matrix exponential:
\[
\exp\left( \Delta \begin{pmatrix} Q & \partial Q/\partial \theta_{j} \\ 0 & Q \end{pmatrix} \right) = \begin{pmatrix} e^{\Delta Q} & \partial e^{\Delta Q} / \partial \theta_{j} \\ 0 & e^{\Delta Q} \end{pmatrix}.
\]
\end{proposition}

\medskip
\noindent\textbf{Why this matters.} The analytic gradient formula for the matrix exponential is crucial for an efficient MLE implementation. The block-matrix trick computes both $e^{\Delta Q}$ and its derivative in a single $2N \times 2N$ matrix exponential---no numerical differencing required. In practice this yields a significant speedup: each evaluation of $\ell(\theta)$ accompanied by a single block-matrix exponential delivers both the function value and the gradient, ready to feed directly to a quasi-Newton optimizer such as BFGS.

\chapter{Parameter Estimation for the CEV Model}

\section{Model Parameterization}

The CEV model for futures prices has parameters $\theta = (\sigma, \beta) \in \Theta = (0, \infty) \times \R$. For a given parameter vector, the CTMC generator $Q(\sigma, \beta)$ is constructed via the formulas of Theorem~G.3.1.

\section{Construction of the Likelihood}

For a sequence of daily futures prices $\{F_{0}, F_{1}, \ldots, F_{n}\}$ observed at intervals of $\Delta = 1/252$ (one trading day):

\textbf{Step 1 --- Discretization.} Construct the grid $\{F_{i}\}_{i=1}^{N}$, map each observation to its grid index $s_{k}$, and form the state sequence.

\textbf{Step 2 --- Generator construction.} For candidate parameters $(\sigma, \beta)$, compute:
\[
D_{i}(\sigma, \beta) = \frac{\sigma^{2} F_{i}^{2\beta+2}}{F_{0}^{2\beta}}, \qquad i = 1, \ldots, N.
\]

\textbf{Step 3 --- Assemble $Q$ and compute $P_{\Delta}$.}
\[
q_{i,i+1} = \frac{D_{i}}{h_{i}(h_{i} + h_{i-1})}, \quad q_{i,i-1} = \frac{D_{i}}{h_{i-1}(h_{i} + h_{i-1})}, \quad q_{i,i} = -(q_{i,i+1} + q_{i,i-1}),
\]
$P_{\Delta} = e^{\Delta \cdot Q}$ (via Pad\'{e} approximation or uniformization).

\textbf{Step 4 --- Evaluate the log-likelihood.}
\[
\ell(\sigma, \beta) = \sum_{k=1}^{n} \log [P_{\Delta}(\sigma, \beta)]_{s_{k-1}, s_{k}}.
\]

\textbf{Step 5 --- Optimize.} $\hat{\sigma}, \hat{\beta} = \argmax_{\sigma>0, \beta \in \R} \ell(\sigma, \beta)$. Use the Nelder-Mead simplex method (gradient-free) for robustness, followed by BFGS (quasi-Newton) for refinement.

\section{Score Function for the CEV Model}

\begin{proposition}[Derivative of $Q$ with Respect to Parameters]
For the CEV model on a uniform grid:
\[
\frac{\partial q_{i,i\pm 1}}{\partial \sigma} = \frac{\sigma F_{i}^{2\beta+2}}{h^{2} F_{0}^{2\beta}}, \qquad
\frac{\partial q_{i,i\pm 1}}{\partial \beta} = \frac{\sigma^{2} F_{i}^{2\beta+2} \ln(F_{i}/F_{0})}{h^{2} F_{0}^{2\beta}}.
\]
These derivatives feed into the integral formula for the derivative of the matrix exponential, enabling gradient-based optimization.
\end{proposition}

\medskip
\noindent\textbf{Why this matters.} The derivatives of $Q$ with respect to $\sigma$ and $\beta$ both have clean closed forms, enabling exact gradient computation without finite differencing. Note that $\frac{\partial q}{\partial\beta}$ contains $\ln(F_i/F_0)$---when $F_i$ is far from $F_0$ (e.g., at the grid extremities), the derivative is large in absolute value, meaning the estimate of $\beta$ is highly sensitive to those distant observations. This also suggests that the grid domain should not be excessively wide, lest distant noise contaminate the precision of the parameter estimates.

\begin{theorem}[CEV-DEJD Generator Parameterization]
The CTMC generator $Q(\theta) = Q_{D}(\sigma, \beta) + Q_{J}(\lambda, p, \eta_{1}, \eta_{2})$ is assembled as:
\begin{enumerate}[label=(\arabic*)]
    \item $Q_{D}$: The CEV tridiagonal generator (as in Chapter~G), depending only on $(\sigma, \beta)$.
    \item $Q_{J}$: The jump generator, whose entries $q_{ij}^{J}$ depend on $(\lambda, p, \eta_{1}, \eta_{2})$ through the case-by-case formulas of Appendix~I. The diagonal is $q_{ii}^{J} = -\sum_{j \neq i} q_{ij}^{J}$.
\end{enumerate}
The combined generator is $Q = Q_{D} + Q_{J}$, which automatically satisfies (C1)--(C2).
\end{theorem}

\medskip
\noindent\textbf{Why this matters.} The assembly of the $Q$ matrix for the CEV-DEJD model showcases the advantage of modular design---the diffusion and jump components are constructed independently and then added. This is analogous to building a baseline CO2 emissions model (pure diffusion) and then overlaying an extreme-climate-event model (jumps), each controlled by its own parameter subset. In nonlinear optimization, the diffusion parameters $(\sigma,\beta)$ primarily affect the likelihood of continuous fluctuations, while the jump parameters $(\lambda,p,\eta_1,\eta_2)$ primarily affect the likelihood of large swings.

\begin{proposition}[Gradient Coupling]
The gradient of the log-likelihood with respect to $\beta$ receives contributions from:
\[
\frac{\partial \ell}{\partial \beta} = \sum_{k=1}^{n} \frac{1}{[P_{\Delta}]_{s_{k-1}, s_{k}}} \left( \left[\frac{\partial P_{\Delta}}{\partial Q_{D}} \frac{\partial Q_{D}}{\partial \beta}\right]_{s_{k-1}, s_{k}} + \left[\frac{\partial P_{\Delta}}{\partial Q_{J}} \frac{\partial Q_{J}}{\partial \beta}\right]_{s_{k-1}, s_{k}} \right).
\]
Both the diffusion variance $D_{i} \propto F_{i}^{2\beta+2} / F_{0}^{2\beta}$ and the jump intensity $\lambda(F_{i}) = \lambda (F_{i}/F_{0})^{\beta}$ contribute. Neglecting either term biases the estimate.
\end{proposition}

\medskip
\noindent\textbf{Why this matters.} $\beta$ appears simultaneously in the diffusion variance and the jump intensity, creating a cross-coupling of parameters---changing $\beta$ affects both $Q_D$ and $Q_J$, and hence $P_\Delta$. If $\beta$ is estimated using only the diffusion portion of the likelihood while ignoring the jump portion, a systematic bias results. This is why the two-stage estimation procedure (CEV first, then EM for jump parameters, then joint refinement) must recalibrate $\beta$ jointly in the third stage.

\chapter{Parameter Estimation for the CEV-DEJD Model}

\section{Parameter Vector and Constraints}

The CEV-DEJD model has six parameters:
\[
\theta = (\sigma, \beta, \lambda, p, \eta_{1}, \eta_{2}) \in (0,\infty) \times \R \times (0,\infty) \times [0,1] \times (2,\infty) \times (0,\infty).
\]
The constraint $\eta_{1} > 2$ (strengthened from $\eta_{1} > 1$) is required for finiteness of the second moment of the jump component.

\section{The Generator as a Function of Parameters}

\begin{theorem}[CEV-DEJD Generator Parameterization]
The CTMC generator $Q(\theta) = Q_{D}(\sigma, \beta) + Q_{J}(\lambda, p, \eta_{1}, \eta_{2})$ is assembled as:
\begin{enumerate}[label=(\arabic*)]
    \item $Q_{D}$: The CEV tridiagonal generator (as in Chapter~G), depending only on $(\sigma, \beta)$.
    \item $Q_{J}$: The jump generator, whose entries $q_{ij}^{J}$ depend on $(\lambda, p, \eta_{1}, \eta_{2})$ through the case-by-case formulas of Appendix~I. The diagonal is $q_{ii}^{J} = -\sum_{j \neq i} q_{ij}^{J}$.
\end{enumerate}
The combined generator is $Q = Q_{D} + Q_{J}$, which automatically satisfies (C1)--(C2).
\end{theorem}

\section{Two-Stage Estimation Strategy}

Estimating all six parameters simultaneously via direct numerical optimization of the log-likelihood is challenging due to the presence of both continuous (diffusion) and discontinuous (jump) components, each contributing to the total variance. A two-stage procedure separates the estimation:

\textbf{Stage 1 --- Pure CEV Estimation.} Estimate $(\sigma, \beta)$ by maximizing the CEV likelihood, ignoring jumps. This provides an initial approximation of the continuous volatility and elasticity.

\textbf{Stage 2 --- Jump Parameter Estimation via the EM Algorithm.}
\begin{enumerate}[label=(\arabic*)]
    \item \textbf{E-step (Expectation):} Given current parameter estimates, compute the posterior probability that each observed transition $s_{k-1} \to s_{k}$ involved a jump. This uses Bayes' theorem:
    \[
    \Pbb(\text{jump} \mid s_{k-1} \to s_{k}; \theta) = \frac{[P_{J}(\Delta)]_{s_{k-1}, s_{k}}}{[P_{\Delta}]_{s_{k-1}, s_{k}}},
    \]
    where $P_{J}(\Delta)$ is the jump-only contribution to the transition probability (isolated via the additive decomposition of the generator).

    \item \textbf{M-step (Maximization):} Update $(\lambda, p, \eta_{1}, \eta_{2})$ by maximizing the expected complete-data log-likelihood, which separates into independent maximizations for the Poisson intensity $\lambda$, the jump direction probability $p$, and the exponential decay rates $(\eta_{1}, \eta_{2})$.
\end{enumerate}

\textbf{Stage 3 --- Joint Refinement.} Use the Stage~1 and Stage~2 estimates as initial values for a full joint MLE optimization of $\ell(\theta)$ over all six parameters, via a quasi-Newton method with numerical gradients.

\section{Handling the State-Dependent Jump Intensity}

The CEV-DEJD model specifies $\lambda(F_{t}) = \lambda (F_{t}/F_{0})^{\beta}$. The parameter $\beta$ thus appears in \emph{both} the diffusion generator $Q_{D}$ and the jump generator $Q_{J}$. This joint dependence couples the diffusion and jump estimation; the two-stage procedure mitigates the coupling by first estimating $\beta$ from the continuous component, then refining jointly.

\begin{proposition}[Gradient Coupling]
The gradient of the log-likelihood with respect to $\beta$ receives contributions from:
\[
\frac{\partial \ell}{\partial \beta} = \sum_{k=1}^{n} \frac{1}{[P_{\Delta}]_{s_{k-1}, s_{k}}} \left( \left[\frac{\partial P_{\Delta}}{\partial Q_{D}} \frac{\partial Q_{D}}{\partial \beta}\right]_{s_{k-1}, s_{k}} + \left[\frac{\partial P_{\Delta}}{\partial Q_{J}} \frac{\partial Q_{J}}{\partial \beta}\right]_{s_{k-1}, s_{k}} \right).
\]
Both the diffusion variance $D_{i} \propto F_{i}^{2\beta+2} / F_{0}^{2\beta}$ and the jump intensity $\lambda(F_{i}) = \lambda (F_{i}/F_{0})^{\beta}$ contribute. Neglecting either term biases the estimate.
\end{proposition}

\chapter{The Hamilton Filter for Regime-Switching Model Estimation}

\section{Model Setup}

For the RS-CEV model with $M = 2$ regimes, the parameter vector is:
\[
\theta = (\sigma_{1}, \beta_{1}, \sigma_{2}, \beta_{2}, \lambda_{12}, \lambda_{21}).
\]

The regime $R_{t} \in \{1, 2\}$ is unobservable; only the price sequence $\{F_{t_{k}}\}$ is observed. The likelihood must integrate over the unobserved regime path, which is accomplished via the \textbf{Hamilton filter}.

\section{The Filtering Recursion}

\begin{definition}[Filtered and Predicted Probabilities]
Define:
\begin{itemize}
    \item $\xi_{k|k}(m) := \Pbb(R_{t_{k}} = m \mid \F_{t_{k}}^{F})$ — the \textbf{filtered} probability that the regime is $m$ at time $t_{k}$, given observations up to $t_{k}$.
    \item $\xi_{k+1|k}(m) := \Pbb(R_{t_{k+1}} = m \mid \F_{t_{k}}^{F})$ — the \textbf{predicted} probability for time $t_{k+1}$, given observations up to $t_{k}$.
\end{itemize}
\end{definition}

\medskip
\noindent\textbf{Why this matters.} The filtered and predicted probabilities are the two core state variables of the Hamilton filter. The filtered probability incorporates the latest observation and represents the most current information summary; the predicted probability is the best guess of the regime before seeing the new data. The entire prediction-update cycle of the filter alternates between these two probabilities.

\begin{theorem}[Hamilton Filter Recursion]
The filter alternates between a prediction step and an update step.

\textbf{Prediction Step} (propagate regime probabilities using the regime transition matrix):
\[
\xi_{k+1|k}(m) = \sum_{j=1}^{M} \xi_{k|k}(j) \cdot P_{jm}^{R}(\Delta),
\]
where $P^{R}(\Delta) = e^{\Delta \cdot Q^{R}}$ is the regime transition matrix over the observation interval $\Delta$ (with closed form given in Theorem~14.4.6 for $M=2$).

\textbf{Update Step} (Bayes' rule incorporating the new observation):
\[
\xi_{k+1|k+1}(m) = \frac{\xi_{k+1|k}(m) \cdot f(F_{t_{k+1}} \mid F_{t_{k}}, R_{t_{k+1}} = m; \theta)}{\sum_{j=1}^{M} \xi_{k+1|k}(j) \cdot f(F_{t_{k+1}} \mid F_{t_{k}}, R_{t_{k+1}} = j; \theta)},
\]
where $f(F_{t_{k+1}} \mid F_{t_{k}}, R_{t_{k+1}} = m; \theta)$ is the conditional density of the observed price given the previous price and the current regime.
\end{theorem}

\medskip
\noindent\textbf{Why this matters.} The Hamilton filter is the heart of regime-switching model estimation. The prediction step is like a weather forecast: using yesterday's weather probabilities and the weather transition matrix, it predicts today's weather probabilities. The update step is like opening the window to look outside (observing the new price) and then revising the prediction via Bayes' rule. The cycle of guess-look-correct continues through the entire observation sequence, ultimately delivering not only the log-likelihood but also the posterior probability of each regime at every time point---which can itself be used to identify when the market switches from bull to bear.

\medskip
\noindent\textbf{What remains to be built.} The recursion above is only the probabilistic skeleton. To turn it into an estimable model one still needs two additional objects: a state-dependent observation density and the induced predictive log-likelihood. These are derived in the next two sections, where the CEV dynamics enter the filter explicitly.

\section{State-Dependent Observation Density}

\begin{proposition}[Conditional Density for the RS-CEV Model]
Under the CEV dynamics in regime $m$, the log-return over $[t_{k}, t_{k+1}]$ is approximately conditionally normal (via Euler approximation):
\[
\log\left(\frac{F_{t_{k+1}}}{F_{t_{k}}}\right) \mid (F_{t_{k}}, R_{t_{k+1}} = m) \sim \mathcal{N}\left(0, \sigma_{m}^{2} \left(\frac{F_{t_{k}}}{F_{0}}\right)^{2\beta_{m}} \Delta \right).
\]

The observation density is:
\[
f(F_{t_{k+1}} \mid F_{t_{k}}, R_{t_{k+1}} = m) = \frac{1}{F_{t_{k+1}} \sqrt{2\pi \sigma_{m}^{2} (F_{t_{k}}/F_{0})^{2\beta_{m}} \Delta}} \exp\left(-\frac{\left[\log(F_{t_{k+1}}/F_{t_{k}})\right]^{2}}{2 \sigma_{m}^{2} (F_{t_{k}}/F_{0})^{2\beta_{m}} \Delta}\right).
\]
\end{proposition}

\begin{proof}
From the SDE in regime $m$: $dF_{t} = \sigma_{m} F_{t}^{\beta_{m}+1} / F_{0}^{\beta_{m}} dW_{t}$. Applying It\^{o}'s lemma to $\log F_{t}$: $d(\log F_{t}) = -\frac{1}{2} \sigma_{m}^{2} (F_{t}/F_{0})^{2\beta_{m}} dt + \sigma_{m} (F_{t}/F_{0})^{\beta_{m}} dW_{t}$. Discretizing over $\Delta$, the drift term is $\mathcal{O}(\Delta)$ and the diffusion term is $\mathcal{O}(\sqrt{\Delta})$. For daily data ($\Delta = 1/252$), the drift is negligible relative to the volatility, justifying the zero-mean normal approximation.
\end{proof}

\section{Log-Likelihood Construction}

\begin{theorem}[Log-Likelihood via the Hamilton Filter]
\[
\boxed{\ell(\theta) = \sum_{k=0}^{n-1} \log\left( \sum_{m=1}^{M} \xi_{k+1|k}(m) \cdot f(F_{t_{k+1}} \mid F_{t_{k}}, R_{t_{k+1}} = m; \theta) \right)},
\]
with the initialization $\xi_{0|0}(m) = \pi_{m}$ (the stationary distribution of the regime process).
\end{theorem}

\begin{proof}
By the law of total probability, the conditional density of $F_{t_{k+1}}$ given $\F_{t_{k}}^{F}$ is the weighted sum over the possible regimes at $t_{k+1}$:
\[
f(F_{t_{k+1}} \mid \F_{t_{k}}^{F}) = \sum_{m=1}^{M} \Pbb(R_{t_{k+1}} = m \mid \F_{t_{k}}^{F}) \cdot f(F_{t_{k+1}} \mid \F_{t_{k}}^{F}, R_{t_{k+1}} = m) = \sum_{m=1}^{M} \xi_{k+1|k}(m) \cdot f(F_{t_{k+1}} \mid F_{t_{k}}, R_{t_{k+1}} = m).
\]
The log-likelihood is the sum over $k$ of the log of this predictive density.
\end{proof}

\section{Implementation of the Filter}

\begin{algorithm}
\caption{Hamilton Filter for RS-CEV}
\begin{enumerate}[label=(\arabic*)]
    \item Initialize: $\xi_{0|0}(m) = \pi_{m}$, $m = 1, 2$.
    \item For $k = 0, \ldots, n-1$:
    \begin{enumerate}
        \item \textbf{Predict:} $\xi_{k+1|k}(m) = \sum_{j=1}^{2} \xi_{k|k}(j) P_{jm}^{R}(\Delta)$, for $m = 1, 2$.
        \item \textbf{Compute observation densities:} $f_{m} = f(F_{t_{k+1}} \mid F_{t_{k}}, R_{t_{k+1}} = m; \theta)$.
        \item \textbf{Update:} $\xi_{k+1|k+1}(m) = \frac{\xi_{k+1|k}(m) f_{m}}{\sum_{j=1}^{2} \xi_{k+1|k}(j) f_{j}}$, for $m = 1, 2$.
        \item \textbf{Accumulate log-likelihood:} $\ell \leftarrow \ell + \log(\sum_{j=1}^{2} \xi_{k+1|k}(j) f_{j})$.
    \end{enumerate}
    \item Return $\ell(\theta)$.
\end{enumerate}
\end{algorithm}

\chapter{Model Selection and Diagnostics}

\section{Information Criteria}

\begin{definition}[Akaike Information Criterion --- AIC]
\[
\text{AIC} = -2 \ell(\hat{\theta}) + 2d,
\]
where $d$ is the number of estimated parameters. The model with the smallest AIC is preferred, balancing goodness-of-fit against model complexity.
\end{definition}

\medskip
\noindent\textbf{Why this matters.} The core idea of AIC is that good fit is not everything---simplicity has merit too. $-2\ell(\hat{\theta})$ measures the model's fit (the smaller the better), and $+2d$ is a penalty for the number of parameters---each extra parameter raises AIC by 2, guarding against overfitting. When choosing between CEV (2 parameters), CEV-DEJD (6 parameters), and RS-CEV (6 parameters), AIC automatically penalizes the more complex models and selects them only if the likelihood improvement from jumps or regimes exceeds $2d$.

\begin{definition}[Bayesian Information Criterion --- BIC]
\[
\text{BIC} = -2 \ell(\hat{\theta}) + d \log n,
\]
where $n$ is the number of observations. BIC penalizes model complexity more heavily than AIC for large $n$, favoring parsimony.
\end{definition}

\medskip
\noindent\textbf{Why this matters.} BIC imposes a heavier penalty for large samples ($d\log n$ vs $2d$). For $n=1000$ trading days (about 4 years), $\log n \approx 6.9$, so BIC penalizes each extra parameter by 6.9 rather than 2. BIC tends to favor simpler models---asymptotically it selects the true model if it lies in the candidate set. In practice, the choice between AIC and BIC depends on the goal: AIC for forecasting new data, BIC for explaining the data-generating mechanism.

\begin{proposition}[Complexity Summary]
For $N$ spatial states and $L$ time steps:
\[
\begin{array}{c|c|c|c}
\text{Model Class} & Q \text{ Structure} & \text{DP Cost per Step} & \text{Total Cost} \\
\hline
\text{Diffusion} & \text{Tridiagonal} & \mathcal{O}(N) & \mathcal{O}(NL) \\
\text{Jump-Diffusion} & \text{Tridiagonal + Dense} & \mathcal{O}(N^{2}) & \mathcal{O}(N^{2}L) \\
\text{Regime-Switching} & \text{Block $M \times M$} & \mathcal{O}(M^{2} N^{2}) & \mathcal{O}(M^{2} N^{2} L)
\end{array}
\]
\end{proposition}

\medskip
\noindent\textbf{Why this matters.} The complexity table is the operations manual for deploying a CTMC pricing system. Pure diffusion is the fastest (milliseconds when $N$ and $L$ are in the hundreds); jump-diffusion, whose jumps make $Q$ dense, requires $\mathcal{O}(N^2)$ dense matrix-vector multiplication per dynamic-programming step; regime-switching, whose block structure multiplies the matrix size by $M$, compounds the cost by a factor of $M^2$. When choosing a model, consider not only statistical fit but also computational feasibility---for market-maker-level real-time pricing, the tridiagonal diffusion model has an enormous speed advantage.

\begin{theorem}[Separation of Concerns]
Steps 1--3 of the pipeline are \textbf{model-specific}: they translate model parameters into the generator matrix $Q^{(N)}$. Steps 4--7 are \textbf{model-agnostic}: they operate solely on $Q^{(N)}$ and the payoff $\boldsymbol{\phi}$, without any knowledge of the underlying model class, SDE, or parameter interpretation. This separation is the central architectural insight of the CTMC approximation framework: once the generator is constructed, the pricing engine is universal.
\end{theorem}

\medskip
\noindent\textbf{Why this matters.} The key design philosophy here is front-end customization with a back-end that is universal. The pricing engine sees only the $Q$ matrix and the $\phi$ vector; it neither knows nor cares whether $Q$ came from CEV, Merton, or regime-switching. This is analogous to a universal remote control---whether the TV is a Sony or a Samsung, as long as it emits a standard infrared signal, the remote can control it. This architecture means that adding a new model requires only writing a parameter-to-$Q$ translation layer, without touching any pricing core code.

\begin{theorem}[Consolidated Convergence of the CTMC American Option Price]
Let $\{X_{t}\}$ be a Feller process satisfying the conditions for the CTMC approximation (Lipschitz coefficients, bounded jump intensity, and an ergodic regime process when regimes are present). Let $\hat{\theta}_{n}$ be the MLE based on $n$ observations, and let $V_{0}^{(N, L)}(\hat{\theta}_{n})$ be the Bermudan-CTMC price with $N$ spatial states, $L$ exercise dates, and estimated parameters. Suppose that $\hat{\theta}_{n} \to \theta_{0}$ almost surely, the spatial mesh satisfies $h_N \to 0$, the truncated grid domains expand to the full state space, and the time step $\Delta t_L = T/L \to 0$. Then, along every refinement sequence satisfying these conditions,

\[
\lim_{n,N,L \to \infty} V_{0}^{(N, L)}(\hat{\theta}_{n}) = V(0, X_{0}; \theta_{0}),
\]

where $V(0, X_{0}; \theta_{0})$ is the true American option price under the data-generating parameter $\theta_{0}$, almost surely.
\end{theorem}

\medskip
\noindent\textbf{Why this matters.} This final convergence theorem is the capstone of the entire theoretical edifice. It combines statistical consistency ($\hat{\theta}_{n} \to \theta_{0}$), spatial convergence of the CTMC generator ($h_N \to 0$ and expanding domains), and temporal convergence of Bermudan exercise to American exercise ($\Delta t_L \to 0$). The theorem says that the observed-data calibration layer, the process-approximation layer, and the optimal-stopping layer can be refined together without changing the limiting price. That is the difference between a pricing algorithm and a convergent pricing theory.

% ═══════════════════════════════════════════════════════════════
% APPENDICES: Reference Material
% ═══════════════════════════════════════════════════════════════
\appendix
\renewcommand{\thechapter}{\Alph{chapter}}

\chapter{Taylor Expansion}
\label{app:taylor}

\begin{theorem}[Taylor's Theorem with Lagrange Remainder]
Let $f \in C^{k+1}(\R)$. For any $x, \Delta \in \R$,
\[
f(x + \Delta) = \sum_{n=0}^{k} \frac{f^{(n)}(x)}{n!} \Delta^{n} + \frac{f^{(k+1)}(\xi)}{(k+1)!} \Delta^{k+1},
\]
where $\xi$ is some point between $x$ and $x + \Delta$.
\end{theorem}

\medskip
\noindent\textbf{Why this matters.} Taylor expansion is the mathematical engine of the moment-matching method. When we expand $f(x_j) - f(x_i)$ in powers of $\Delta = x_j - x_i$, the first-order term $\Delta f'(x_i)$ corresponds to the drift moment and the second-order term $\frac{1}{2}\Delta^2 f''(x_i)$ corresponds to the diffusion moment. The Lagrange remainder supplies an exact error expression---it tells us that the error from neglecting third- and higher-order terms is controlled by the $(k+1)$-th derivative of $f$ evaluated at some intermediate point.

\begin{theorem}[Taylor's Theorem with Integral Remainder]
Under the same hypotheses,
\[
f(x + \Delta) = \sum_{n=0}^{k} \frac{f^{(n)}(x)}{n!} \Delta^{n} + \frac{1}{k!} \int_{x}^{x+\Delta} (x+\Delta - t)^{k} f^{(k+1)}(t) \, dt.
\]
\end{theorem}

\medskip
\noindent\textbf{Why this matters.} The integral-remainder form is sometimes more convenient than the Lagrange remainder---especially when the remainder must pass through an expectation. For example, when taking the expectation over a random jump size $Y$, the integral remainder allows exchanging the order of integration with the distribution of $Y$, yielding tighter error bounds.

\chapter{It\^{o}'s Formula}

\begin{theorem}[It\^{o}'s Formula for Continuous Semimartingales]
Let $X$ be a continuous semimartingale with $dX_{t} = \mu_{t} dt + \sigma_{t} dW_{t}$. For $f \in C^{1,2}([0,\infty) \times \R)$,
\[
df(t, X_{t}) = \frac{\partial f}{\partial t} dt + \frac{\partial f}{\partial x} dX_{t} + \frac{1}{2} \frac{\partial^{2} f}{\partial x^{2}} d\langle X \rangle_{t},
\]
with $d\langle X \rangle_{t} = \sigma_{t}^{2} dt$.
\end{theorem}

\medskip
\noindent\textbf{Why this matters.} It\^{o}'s formula is the chain rule of stochastic calculus. Unlike ordinary calculus, it includes the extra term $\frac{1}{2} f'' d\langle X \rangle_t$---the so-called It\^{o} correction. This term arises from the non-zero quadratic variation of Brownian motion ($(dW_t)^2 = dt$). In ordinary calculus this is a second-order infinitesimal and can be neglected; in stochastic calculus it is of the same order as $dt$. This correction term is precisely the origin of the $\frac{1}{2}\sigma^2 f''$ component in the generator.

\chapter{It\^{o}'s Formula for Jump Processes}

\begin{theorem}[It\^{o}'s Formula for General Semimartingales]
Let $X$ be a semimartingale of the form
\[
dX_{t} = \mu_{t} dt + \sigma_{t} dW_{t} + \int_{\R} h(t, y) \widetilde{N}(dt, dy) + \int_{\R} h(t, y) \nu(dy) dt.
\]
For $f \in C^{2}(\R)$,
\[
\begin{aligned}
f(X_{t}) &= f(X_{0}) + \int_{0}^{t} f'(X_{s-}) \mu_{s} ds + \int_{0}^{t} f'(X_{s-}) \sigma_{s} dW_{s} \\
&\quad + \frac{1}{2} \int_{0}^{t} f''(X_{s-}) \sigma_{s}^{2} ds \\
&\quad + \int_{0}^{t} \int_{\R} \left[f(X_{s-} + h(s, y)) - f(X_{s-})\right] \widetilde{N}(ds, dy) \\
&\quad + \int_{0}^{t} \int_{\R} \left[f(X_{s-} + h(s, y)) - f(X_{s-})\right] \nu(dy) ds.
\end{aligned}
\]
\end{theorem}

\chapter{Matrix Exponential}

\begin{definition}[Matrix Exponential]
For an $N \times N$ real matrix $A$,
\[
e^{A} := \sum_{k=0}^{\infty} \frac{A^{k}}{k!}.
\]
\end{definition}

\medskip
\noindent\textbf{Why this matters.} The matrix exponential is the universal solution of linear ODE systems. If $\frac{d\mathbf{p}}{dt} = Q \mathbf{p}$, then $\mathbf{p}(t) = e^{tQ}\mathbf{p}(0)$---exactly analogous to the scalar case $\frac{dy}{dt} = ay \Rightarrow y(t) = e^{at}y(0)$. Each row of $e^{tQ}$ is the conditional distribution at time $t$ starting from the corresponding state---it is the bridge that converts the rate matrix $Q$ into the probability matrix $P_t$.

\begin{theorem}[Uniformization for CTMC Generators]
For a CTMC generator matrix $Q$, let $\Lambda := \max_{i} |q_{ii}|$ and $S := I + Q/\Lambda$. Then $S$ is a stochastic matrix, and
\[
e^{tQ} = \sum_{m=0}^{\infty} e^{-\Lambda t} \frac{(\Lambda t)^{m}}{m!} S^{m}.
\]
\end{theorem}

\medskip
\noindent\textbf{Why this matters.} Uniformization converts a continuous-time CTMC into a combination of a discrete-time random walk and a Poisson random clock. The intuition: treat every potential CTMC transition as occurring at a common rate $\Lambda$ (the fastest possible rate), then at each virtual discrete moment perform a state transition with probability $S$. The Poisson distribution $\operatorname{Poisson}(\Lambda t)$ determines how many virtual steps occur in $[0,t]$. The numerical advantage is that $S$ is a stochastic matrix (entries in $[0,1]$), avoiding the numerical instability that can arise from directly exponentiating $Q$.

\chapter{Generator Reference Card}

\begin{center}
\begin{tabular}{>{\raggedright\arraybackslash}p{4.5cm}>{\raggedright\arraybackslash}p{7.5cm}}
\toprule
\textbf{Process} & \textbf{Generator $\calL f(x)$} \\
\midrule
Brownian Motion & $\frac{1}{2} f''(x)$ \\
Brownian with Drift & $\mu f'(x) + \frac{1}{2} \sigma^{2} f''(x)$ \\
Geometric Brownian Motion & $\mu x f'(x) + \frac{1}{2} \sigma^{2} x^{2} f''(x)$ \\
Ornstein-Uhlenbeck & $\kappa(\theta - x) f'(x) + \frac{1}{2} \sigma^{2} f''(x)$ \\
CIR Process & $\kappa(\theta - x) f'(x) + \frac{1}{2} \sigma^{2} x f''(x)$ \\
CEV Process & $\frac{1}{2} \sigma^{2} x^{2\beta+2} / F_{0}^{2\beta} \cdot f''(x)$ \\
Merton Jump-Diffusion & $\mu x f' + \frac{1}{2} \sigma^{2} x^{2} f'' + \lambda \int [f(x e^{y}) - f(x)] F_{Y}(dy)$ \\
Regime-Switching & $\mu_{m} f_{x} + \frac{1}{2} \sigma_{m}^{2} f_{xx} + \sum_{n \neq m} \lambda_{mn}[f(x,n) - f(x,m)]$ \\
General Jump-Diffusion & $\mu f' + \frac{1}{2} \sigma^{2} f'' + \lambda \int [f(x+y) - f(x)] F_{Y}(dy)$ \\
\bottomrule
\end{tabular}
\end{center}

\end{document}
